% SPDX-License-Identifier: LPPL-1.3c
% Copyright (C) 2026 Christophe Jorssen
% Author and Current Maintainer: Christophe Jorssen <christophe.jorssen@gmail.com>
% LPPL maintenance status: maintained
% This file is part of luafemm. See LICENSE and LICENSES.md for its terms.
% Build from the repository root with make manual. Examples execute on LuaTeX.
% The PGF documentation environments also used by tikz-ext provide indexed
% keys, library/command signatures and adjacent source/rendered examples.
\documentclass[a4paper,11pt]{ltxdoc}
\usepackage[margin=23mm]{geometry}
\usepackage{fontspec}
\usepackage{amsmath,amssymb,booktabs,longtable,multicol}
\usepackage{xkeyval,calc,listings,fp,pifont,xxcolor}
\usepackage{tikz,pgfplots}
\usepackage{makeidx,needspace,adjustbox}
\makeindex
\usepackage[hyperindex=false]{hyperref}
\usepackage{xurl}
\hypersetup{colorlinks=true,linkcolor=blue!55!black,urlcolor=blue!55!black,
  pdftitle={The luafemm TikZ and PGFPlots libraries},pdfauthor={Christophe Jorssen}}
\input{pgfmanual-en-macros.tex}
% SPDX-License-Identifier: LPPL-1.3c
% Copyright (C) 2026 Christophe Jorssen
% Author and Current Maintainer: Christophe Jorssen <christophe.jorssen@gmail.com>
% LPPL maintenance status: maintained
% This file is part of luafemm. See LICENSE and LICENSES.md for its terms.
% Full sources are highlighted by PGF. Their bodies execute in this TeX run;
% only the document wrapper and preamble-only package loading are omitted.
% Groups isolate example styles and temporary macros. Long source listings can
% span pages; the rendered result is scaled only to fit the available paper.
\newcommand\luafemmsource[1]{%
  \par\medskip\noindent\textbf{Source: }\nolinkurl{examples/#1.tex}\par
  \begingroup
    \codeexample[multipage,from file=examples/#1.tex]
  \endgroup
}
\expandafter\def\csname luafemmexamplewidth@materials-profiles\endcsname{180mm}
\newcommand\luafemmexample[1]{%
  \label{ex:#1}%
  \luafemmsource{#1}%
  \par\medskip\noindent\begin{minipage}{\linewidth}
  \textbf{Result}\label{ex:#1:result}\par\nobreak
  \begingroup
    \edef\jobname{manual-\detokenize{#1}}%
    \begin{adjustbox}{max totalsize={.96\linewidth}{.78\textheight},center}
      \ifcsname luafemmexamplewidth@#1\endcsname
        \edef\luafemmexamplewidth{\csname luafemmexamplewidth@#1\endcsname}%
      \else\def\luafemmexamplewidth{160mm}\fi
      \begin{minipage}{\luafemmexamplewidth}
        \input{build/doc/bodies/#1.tex}%
      \end{minipage}
    \end{adjustbox}
  \endgroup
  \end{minipage}\par
}
% Named exception: a LaTeX manual cannot execute plain/ConTeXt document
% wrappers. The manual build compiles these three format demonstrations from
% their displayed source in the actual engines and verifies fresh output.
\newcommand\luafemmformatexample[1]{%
  \label{ex:#1}%
  \luafemmsource{#1}%
  \par\noindent\textit{Format demonstration: executed separately in its
    native format, since this manual uses LuaLaTeX.}\par
  \includegraphics[width=.96\linewidth,height=.7\textheight,keepaspectratio]
    {build/doc/examples/#1.pdf}\par
}

\usetikzlibrary{femm,calc}
\usepgfplotslibrary{femm}
\pgfplotsset{compat=1.18}
\tracinglostchars=3
\setlength\emergencystretch{2em}
\title{The \textsf{luafemm} TikZ and PGFPlots libraries\\
  \large Generic LuaTeX magnetic fields and quantitative profiles}
\author{Christophe Jorssen\\
  \normalsize\href{mailto:christophe.jorssen@gmail.com}{\nolinkurl{christophe.jorssen@gmail.com}}}
\date{Version \luafemmversion\quad---\quad\luafemmversiondate}
\begin{document}
\maketitle
\begin{abstract}
Declare magnetic media on ordinary TikZ paths, solve a finite-element
problem (static or linear harmonic, planar or axisymmetric) entirely in Lua, and draw oriented field lines or PGFPlots profiles.
The same generic implementation works with plain LuaTeX, LuaLaTeX and
ConTeXt MkIV. This manual documents every supported public key, the numerical
conventions, and the limitations of the supported models.
\end{abstract}
% SPDX-License-Identifier: LPPL-1.3c
% Copyright (C) 2026 Christophe Jorssen
% Author and Current Maintainer: Christophe Jorssen <christophe.jorssen@gmail.com>
% LPPL maintenance status: maintained
% This file is part of luafemm. See LICENSE and LICENSES.md for its terms.
\section*{Author, licence and scientific validation}
Copyright \textcopyright\ 2026 Christophe Jorssen, author and Current Maintainer
of \textsf{luafemm}. Contact:
\href{mailto:christophe.jorssen@gmail.com}{\nolinkurl{christophe.jorssen@gmail.com}}.
The original package code and documentation are distributed under the
\href{https://www.latex-project.org/lppl/lppl-1-3c.txt}{LaTeX Project Public License,
version 1.3c or later}; the LPPL maintenance status is \texttt{maintained}.
The file inventory and third-party exceptions are recorded in
\texttt{MANIFEST.md} and \texttt{LICENSES.md}.

Publication is planned on CTAN. Development is hosted at
\href{https://github.com/cjorssen/luafemm}{github.com/cjorssen/luafemm};
please report bugs through the
\href{https://github.com/cjorssen/luafemm/issues}{GitHub issue tracker}.
The development is based on studying the source code of FEMM/xfemm and
Triangle and reimplementing selected formulations and algorithms in
Lua/LuaTeX. No native solver or mesher is embedded, and this is not a
complete port of either program.

\medskip
\noindent This package was \textbf{vibe-coded with OpenAI Codex and the Astra-6 model},
under Christophe Jorssen's direction. Neither AI assistance nor publication
certifies scientific correctness. Independently validate fields, material laws
and mesh/domain convergence against analytical or trusted reference results.

\medskip
The imported FEMM material data and their generated derivatives retain
separate upstream terms, including the Aladdin Free Public License.
They are not relicensed as LPPL data. The distribution therefore contains mixed licensing; these separate terms
apply to the catalogue and its derivatives. Reimplementation
in another language does not itself grant permission to relicense
third-party material. See the provenance section for the retained notices.

\clearpage
\tableofcontents
\clearpage
% SPDX-License-Identifier: LPPL-1.3c
% Copyright (C) 2026 Christophe Jorssen
% Author and Current Maintainer: Christophe Jorssen <christophe.jorssen@gmail.com>
% LPPL maintenance status: maintained
% This file is part of luafemm. See LICENSE and LICENSES.md for its terms.
\section{Getting started}
\label{sec:start}
\begin{tikzlibrary}{femm}
  This library gives evaluated TikZ paths a physical meaning. It supplies
  planar and axisymmetric magnetostatic models, a constrained triangular mesh, material
  choices, oriented field lines and measurement paths. Computation runs
  inside LuaTeX; no native solver or shell escape is required.
\end{tikzlibrary}

\subsection{Requirements and installation}
Use a LuaTeX engine with Lua 5.3 or later, PGF/TikZ, and optionally PGFPlots
for axes. The release was tested with LuaTeX/LuaHBTeX 1.24.0, PGF 3.1.12,
PGFPlots 1.18.3 and ConTeXt MkIV 2026.09.10. It does not support pdfTeX,
XeTeX, or ConTeXt MkII. ConTeXt LMTX is not part of the release test matrix.

Copy the distribution's |tex| tree into a TEXMF tree; keep all Lua and
generic TeX files together under |tex/luatex/luafemm|. The LaTeX wrapper
belongs in |tex/luatex/luafemm| and the ConTeXt wrapper in
|tex/luatex/luafemm|. Refresh the filename database if your TeX
distribution requires it. Keep the supplied licence and provenance files
with redistributed copies. From a source checkout, |make test|,
|make examples| and |make manual| set up the search paths for you.

\subsection{LaTeX}
Either load the convenience package |luafemm| or load the libraries directly:
\begin{codeexample}[code only]
\documentclass{article}
\usepackage{tikz,pgfplots}
\usetikzlibrary{femm}
\usepgfplotslibrary{femm}
\pgfplotsset{compat=1.18}
\begin{document}
% Place a tikzpicture here.
\end{document}
\end{codeexample}
Compile with |lualatex --no-shell-escape document.tex|.
The package wrapper only loads TikZ and the generic library; PGFPlots
remains optional.

\subsection{Plain LuaTeX}
\begin{codeexample}[code only]
\input tikz.tex
\input pgfplots.tex
\usetikzlibrary{femm}
\usepgfplotslibrary{femm}
\pgfplotsset{compat=1.18}
\tikzpicture[femm/problem]
  % Physical paths and field pictures.
\endtikzpicture
\bye
\end{codeexample}
Use |luatex --no-shell-escape document.tex|. The macros are identical to
the LaTeX ones; |\tikzpicture|/|\endtikzpicture| and |\axis|/|\endaxis|
replace the environment syntax.

\subsection{ConTeXt MkIV}
\begin{codeexample}[code only]
\usemodule[luafemm]
\usemodule[pgfplots]
\usepgfplotslibrary{femm}
\pgfplotsset{compat=1.18}
\starttext
\starttikzpicture[femm/problem]
  % Physical paths and field pictures.
\stoptikzpicture
\stoptext
\end{codeexample}
Compile with |context --luatex document.tex|, explicitly selecting MkIV.
Use |\startaxis|/|\stopaxis| inside a |\starttikzpicture| for a graph.
Loading |\usemodule[tikz]| followed by |\usetikzlibrary[femm]| is an
alternative to the wrapper. Repeated loads are harmless.

Appendix~\ref{sec:format-examples} provides the full source and rendered result
for each format, including the shared scene. Throughout this manual, short
self-contained examples show their result beside the code. Longer examples
show the complete source followed by its compiled output; the listings are
read directly from the distributed files. Empty paper margins are cropped
when including a result, and large figures are scaled to fit the manual.
Intermediate tutorial steps and
configuration fragments are identified by their surrounding explanation.

\subsection{A complete small model}
The following example is executed while this manual is built. Its red
rectangle is a permanent magnet pointing up. Every path declared as a
region is also used by the mesher.
\begin{codeexample}[width=5cm]
\begin{tikzpicture}[
  femm/problem={xmin=-15,xmax=15,
    ymin=-15,ymax=15,mesh size=2}]
  \filldraw[femm/region={
    material=cast-alnico-5-lng37,
    magnetization angle=90},
    fill=red!20]
    (-3,-5) rectangle (3,5);
  \pic[femm/lines=13] {femm field};
\end{tikzpicture}
\end{codeexample}

\subsection{The order of operations}
\label{sec:workflow}
The document declares a model progressively. The following order is
essential, including when the calculation comes from a cache:
\begin{enumerate}
\item Define reusable TikZ and material styles. A style definition alone
  neither creates a model nor registers a physical region.
\item Put |femm/problem| and any |femm/cache| options on the magnetic
  |tikzpicture|. TikZ first reads its options and styles. luafemm then
  creates one model in its begin-picture hook, using the final problem
  options, and applies queued |femm/excitation| and |femm/boundary| declarations.
  Boundary keys may occur before or after |femm/problem| in the option list.
\item In the picture body, declare \emph{all} physical paths and component
  nodes, and any additional boundary conditions or ideal-window excitations. A physical path is
  registered when its capture postaction runs, after TikZ has resolved
  its geometry. Later regions take precedence wherever their filled
  portions overlap (Section~\ref{sec:domains}).
\item Request a mesh, a solution or a computed field. The first such
  request fixes the physical model: materials, regions and boundary
  declarations must not follow it. A mesh picture alone already reaches
  this point, even though it does not solve the field equations.
\item Draw more field views, inspect numeric results and plot registered
  measurement profiles. These reuse the prepared model. A profile may be
  registered before or after meshing; its registration alone does not
  compute anything. Its first sampling or plotting request solves its
  source model if needed.
\end{enumerate}

\begin{center}\small
\begin{tabular}{p{.49\linewidth}p{.42\linewidth}}
\toprule Operation & Numerical effect \\
\midrule
|femm/region| or a component node & Register physical geometry \\
|femm/profile| & Register a measurement curve only \\
|\pic {femm mesh};| & Mesh if needed; freeze physical declarations \\
|\femmsolve| or |\path[femm/solve];| & Mesh if needed, then solve \\
|\pic {femm field};| & Solve if needed, then draw contours \\
|\pic {femm mean path};| & Solve if needed, draw and optionally register a profile \\
|\femmflux[profile=...]| & Solve the profile's model if needed, integrate section flux \\
|\femmplot| inside an axis & Solve the profile's model if needed, then sample \\
|\femmB|, |\femmstat| & Read an already solved model \\
\bottomrule
\end{tabular}
\end{center}
For example, the safe placement of an optional diagnostic mesh is
\begin{codeexample}[width=5cm]
\begin{tikzpicture}[
  femm/boundary={side=left,type=neumann},
  femm/problem={xmin=-20,xmax=20,ymin=-20,ymax=20,mesh size=3}]
  \path[femm/region={material=pure-iron}]
    (-10,-10) rectangle (10,10);
  \path[femm/profile={name=scan}] (-15,12)--(15,12);
  % All physical declarations are now complete.
  \pic {femm mesh};
  \pic {femm field};
\end{tikzpicture}
\end{codeexample}

Placing that mesh picture before a later material path is an error.
Moving a boundary declaration into a graphical |scope| does not make it
temporary: once registered, it belongs to the model. Ordinary drawing
options remain governed by TeX grouping. Registering another profile or
drawing labels after solving is allowed; adding physical keys to those
labels or to |every path| can inadvertently attempt a late declaration.

For custom |execute at begin picture| code that needs the model, put the
hook after |femm/problem|. The order independence described above applies
to the declarative boundary keys, not to arbitrary executable hooks.
Cache hits obey the same lifecycle. Complete the declarations before any
cache lookup triggered by a mesh, solve or sampling request; the cache
cannot predict regions that will appear later in the document.

\subsection{Notation, scope and stability}
Options have the usual PGF |key=value| syntax. Most public keys start with
|/tikz/femm|; lists introduced by |femm/problem|, |femm/region| and
|femm/profile| have their own documented subkeys. Option values follow TeX
grouping and may be inherited from styles or scopes; registered physical
declarations persist in their model. A new named profile
is a Lua object with document-wide lifetime; this is deliberately different
from a local TeX option.

The defaults in this manual are initial values unless described as a reset
performed by a command. A ``number'' means a finite decimal or scientific
literal without a TeX unit. Only keys explicitly documented as accepting
PGF expressions evaluate such expressions. Graphic dimensions such as
|arrow length| follow ordinary TikZ dimension syntax.

The public interface comprises the keys and commands
documented here. Macros containing |\luafemm@|, capture postactions, and
the internal mesh layout are not extension interfaces. The Lua API is
documented separately in Section~\ref{sec:development}; its internal
tables should not be edited after meshing.

\subsection{Order for a moving machine}
A sliding-band declaration picture is mesh-only: declare all media and sources,
then capture with |femm/sliding band| at its end. Request fields and measurements
in subsequent pictures selecting that template and a |rotor angle|. A field or
measurement before capture meshes too early and is rejected. See
Section~\ref{sec:sliding} for the complete source/output example.

% SPDX-License-Identifier: LPPL-1.3c
% Copyright (C) 2026 Christophe Jorssen
% Author and Current Maintainer: Christophe Jorssen <christophe.jorssen@gmail.com>
% LPPL maintenance status: maintained
% This file is part of luafemm. See LICENSE and LICENSES.md for its terms.
% Original narrative and examples, following the progressive presentation of
% the PGF tutorials (a concrete task, small additions, then complete results).
\clearpage
\section{Tutorial: Two magnets for Maya}
\label{sec:tutorial}

Maya is preparing a lesson on magnetic circuits. She has a drawing of a
U-shaped iron core, a coil, and an armature above the two poles. The drawing
is quite satisfactory, except for the field lines: she has drawn them by hand,
and moving the armature means drawing them all again. She would also like to
answer a question that the drawing alone cannot settle: how strong is the
field in each air gap?

Her plan is to let the paths in her TikZ picture describe the physical model,
ask LuaTeX to compute the field, and use PGFPlots to graph measurements along
a contour. Later she will remove the coil and replace part of the iron with
alnico. Most of the drawing should survive this change.

We shall follow her through both examples. The complete, compilable sources
are \nolinkurl{examples/tutorial-coil.tex} and
\nolinkurl{examples/tutorial-alnico.tex}. They use the same catalogue material,
|pure-iron|, for the soft iron, so there is no B--H table to type. The earlier
|u-electromagnet.tex| example instead uses an illustrative iron curve; its
numerical values need not agree with this tutorial.

\subsection{A useful first decision: do less work}

Before choosing colours, Maya chooses a fairly coarse mesh. She will move
labels, adjust the drawing, and change her mind about the graph several
times. Rebuilding a very fine finite-element model after each such change
would be a poor use of her time.

Her starting value is |mesh size=5|, in millimetres. This is a target spacing,
not a requested number of triangles. Decreasing it generally increases the
number of triangles and the cost of the solve; in a fairly uniform planar
mesh, halving the spacing gives roughly four times as many elements. Geometry
constraints and local refinement change the actual count. The solution time
need not scale by the same factor.

\begin{quote}
\textbf{Draft the picture with a coarse global mesh. Refine the important
regions as needed, then refine and check the numerical result when the
geometry and presentation are nearly final.}
\end{quote}

There is one qualification: coarse must not mean blind to the physics.
Maya's air gaps are narrow. She keeps them explicitly in the geometry and
uses smaller elements there, even in a draft. A plausible-looking picture
on a coarse mesh is useful for composition; it is not yet an accurate
field measurement. More arrows, smoother plot styling, or more measurement
samples would not repair an insufficient mesh.

\subsection{Setting up the document}

Maya uses LuaLaTeX and loads two graphics libraries:
\begin{codeexample}[code only]
\documentclass{article}
\usepackage{tikz,pgfplots}
\usetikzlibrary{femm}
\usepgfplotslibrary{femm}
\pgfplotsset{compat=1.18}
\begin{document}
% The model picture, followed by a separate PGFPlots picture.
\end{document}
\end{codeexample}
Compile with |lualatex --no-shell-escape|. Nothing is sent to an external
solver. The complete example files use the |standalone| class to crop their
output; |article| is sufficient for following the steps here. Plain LuaTeX
and ConTeXt MkIV use the loading commands from Section~\ref{sec:start};
the physical keys and the plotting commands are the same.

The following code fragments are successive additions to the same document.
In particular, the physical paths go in one model picture, while the axes
will go in separate pictures after it. A graph is not a new magnetic problem.

\subsection{First, just draw the iron}

Maya begins with ordinary TikZ paths. The U is a single closed polygon, and
the armature is a rectangle above it. The top of each pole is at $y=30$ mm;
the bottom of the armature is at $y=33$ mm, leaving two 3 mm gaps.
\begin{codeexample}[width=5cm]
\begin{tikzpicture}[x=1mm,y=1mm,scale=.5]
  \filldraw[fill=gray!15]
    (-40,-30)--(40,-30)--(40,30)--(25,30)
    --(25,-15)--(-25,-15)--(-25,30)
    --(-40,30)--cycle;
  \filldraw[fill=gray!15]
    (-40,33) rectangle (40,48);
\end{tikzpicture}
\end{codeexample}

So far, this is only a drawing. The gray fill does not tell the solver that
there is iron inside it. To give the drawing that meaning, Maya introduces
a style that contains both physical and graphical options:
\begin{codeexample}[code only]
\tikzset{
  core/.style={femm/region={material=pure-iron},
    draw=black!70,fill=gray!15}
}
\end{codeexample}
She replaces |[fill=gray!15]| by |[core]| on both paths. The choice
|material=pure-iron| now selects a nonlinear material from the FEMM catalogue.
The colour remains an ordinary TikZ choice; changing it will not change the
material law.

Next she replaces the options of the enclosing picture:
\begin{codeexample}[code only]
\begin{tikzpicture}[femm/problem={xmin=-110,xmax=110,
  ymin=-95,ymax=115,mesh size=5}]
  % The two \filldraw[core] paths go here.
\end{tikzpicture}
\end{codeexample}
The |femm/problem| key installs millimetre coordinates and creates a rectangular
air domain larger than the device. Its outside boundary has $A_z=0$.
There is no need to draw that rectangle, nor to fill every unassigned part
with an air region. In fact, a large air region drawn \emph{after} the core
would replace the iron there: later physical regions take precedence.

Maya removes |scale=.5| for the working picture. A scale applied to the whole
model picture would only change its displayed size; it would not halve its
physical dimensions. By contrast, scaling a region inside the picture changes
the model geometry. Keeping this distinction in mind saves some surprisingly
unproductive experiments.

\subsection{Giving the coil a current}

In this two-dimensional section, the winding appears as two rectangles on
opposite sides of the left leg. Current leaves the page through one section
and enters it through the other. Maya first defines their common appearance
and material:
\begin{codeexample}[code only]
\tikzset{
  winding/.style={femm/region={material=copper},
    draw=orange!70!black,fill=orange!25}
}
\end{codeexample}
Then she adds two paths, inside the model picture and after the iron:
\begin{codeexample}[code only]
\filldraw[winding,femm/current density=6000/(6*26)*1000000]
  (-48,-10) rectangle (-42,16);
\filldraw[winding,femm/current density=-6000/(6*26)*1000000]
  (-23,-10) rectangle (-17,16);
\end{codeexample}
Each section is $6\times26$ mm$^2$. For $NI=6000$ ampere-turns, its imposed
current density is
\[
 J_z=\frac{NI}{S}=\frac{6000}{6\times26}\times10^6\ \mathrm{A/m^2},
\]
with the opposite sign in the return section. The factor $10^6$ converts
A/mm$^2$ to A/m$^2$. The key accepts this PGF expression directly. Copper
alone does not create this winding excitation; the signed current density does.

These are homogenised winding sections, not a turn-by-turn conductor model.
If Maya changes their areas but wants to keep $NI$ fixed, she must change
the current-density expression too. The solver does not infer that intention
from the picture.

\subsection{Looking at the mesh before asking for the field}

A 5 mm global spacing is quite coarse compared with a 3 mm gap. Maya therefore
adds two invisible air rectangles with a smaller local spacing:
\begin{codeexample}[code only]
\tikzset{
  gap mesh/.style={femm/region={material=air,mesh size=1}}
}
\end{codeexample}
\begin{codeexample}[code only]
\path[gap mesh] (-40,30) rectangle (-25,33);
\path[gap mesh] (25,30) rectangle (40,33);
\end{codeexample}
These paths both assign air and request local refinement. Here they occupy
only the gaps, so placing them after the solids is safe. An invisible path
still has a physical effect when it carries |femm/region|.

To see what she has asked the mesher to do, she can finish the picture with
\begin{codeexample}[code only]
\begin{scope}
  \clip (-58,-42) rectangle (58,57);
  \pic[femm/mesh style={black!25}] {femm mesh};
\end{scope}
\end{codeexample}
The clip keeps the distant air out of the illustration. It does not reduce
the computational domain. A mesh picture builds the triangulation but does
not solve for the field, which makes it a useful intermediate check.

All material paths must come \emph{before} this request. Once meshing has
started, it is too late to add the return conductor or a forgotten refinement
region to that model. Maya now knows where her last physical path belongs.

\subsection{The first computed field lines}

The interesting part is pleasantly short. Maya replaces the mesh picture
by a field picture, keeping the same clip:
\begin{codeexample}[code only]
\begin{scope}
  \clip (-58,-42) rectangle (58,57);
  \pic[femm/lines=17,femm/arrow spacing=26mm,
    femm/field style={teal!70!black}] {femm field};
\end{scope}
\end{codeexample}
This request meshes and solves the model if necessary, then draws contours of
$A_z$. Their arrows follow $\mathbf B$. Maya may put the mesh picture first
and the field picture second to see both; the second request reuses the mesh.
Labels and the $\odot$, $\otimes$ conductor symbols are added afterwards so
they remain readable.

The 17 in |femm/lines=17| counts potential levels, not triangles. Changing it
changes the drawing, not the finite-element resolution. Likewise,
|femm/arrow spacing| controls the displayed arrows. During a compilation,
further field pictures and profile plots reuse the model already solved;
a new compilation builds and solves it again. There is no persistent solver
cache between document runs.

\subsection{Measuring across both gaps}

Field lines show where the flux goes. Maya still needs numbers. She draws a
straight measurement contour halfway up the gaps:
\begin{codeexample}[code only]
\draw[femm/profile={name=coil-gap,samples=101},
  violet,dashed,thick,-Stealth] (-48,31.5)--(48,31.5);
\end{codeexample}
This path belongs inside the model picture. The |femm/profile| key records
it without creating another medium. It may come after the field picture,
since it does not change the material geometry. The chosen height, 31.5 mm,
lies inside the air gaps rather than exactly on an iron/air interface.
A measurement contour need not be closed and need not follow a field line.

After closing the model picture, Maya opens a new picture containing an axis:
\begin{codeexample}[code only]
\begin{tikzpicture}
\begin{axis}[width=13cm,height=5cm,grid=major,no markers,
  xlabel={$x$ (mm)},ylabel={Field (T)},
  legend style={at={(.5,1.03)},anchor=south,legend columns=3}]
  \femmplot[profile=coil-gap,component=B,abscissa=x,
    plot options={black,thick}]
  \addlegendentry{$|\mathbf B|$}
  \femmplot[profile=coil-gap,component=Bx,abscissa=x,
    plot options={blue}]
  \addlegendentry{$B_x$}
  \femmplot[profile=coil-gap,component=By,abscissa=x,
    plot options={red}]
  \addlegendentry{$B_y$}
\end{axis}
\end{tikzpicture}
\end{codeexample}
There is no |femm/problem| option on this graph. Each |\femmplot| names the
profile to sample; it already knows its source model. The command emits a
complete PGFPlots plot, including its terminating semicolon, so Maya adds
only the legend entry afterwards.

|component=B| means the magnitude; |Bx| and |By| retain their signs.
|abscissa=x| plots physical $x$ in millimetres. With the stated winding signs,
the dominant vertical component points upwards in the left gap and downwards
in the right one. The two gaps need not have exactly equal magnitudes:
the winding is on the left leg and the finite-domain solution includes leakage.

\clearpage
\subsection{The electromagnet, assembled}
\label{sec:tutorial-coil-result}
The complete file \nolinkurl{examples/tutorial-coil.tex} appears below, followed
by its picture and graph, computed with the draft settings just introduced.
The small jumps in the curves come from the
piecewise-constant element field; PGFPlots joins the samples with line segments.
\luafemmexample{tutorial-coil}

The triangle count and gap-centre values beneath the graph are printed by
|\femmstat{triangles}| and |\femmB|. They are useful for checking that a mesh
change actually took effect. The latter prints only three decimal places;
use profiles or the Lua API when more numerical precision is needed.

\clearpage
\subsection{Replacing the coil by a permanent magnet}

Maya would now like a field without supplying a current. She keeps the U,
removes \emph{both} copper rectangles, and inserts an alnico crosspiece at its
base. This time the armature will leave 1 mm gaps. She starts a new model
picture with the same outer air domain and the same 5 mm global mesh.

The |core| style remains unchanged. A new style supplies the hard material
and its magnetisation direction:
\begin{codeexample}[code only]
\tikzset{
  magnet/.style={femm/region={material=cast-alnico-5-lng37,
    magnetization angle=180},draw=red!70!black,fill=red!25},
  gap mesh/.style={femm/region={material=air,mesh size=.5}}
}
\end{codeexample}
After drawing the \emph{same U polygon}, she adds these paths in this order:
\begin{codeexample}[code only]
\filldraw[magnet] (-25,-30) rectangle (25,-15);
\filldraw[core] (-40,31) rectangle (40,46);
\path[gap mesh] (-40,30) rectangle (-25,31);
\path[gap mesh] (25,30) rectangle (40,31);
\end{codeexample}
The alnico rectangle overlaps the bottom of the original U. That is deliberate:
because it comes later, it replaces the iron in that part of the crosspiece.
The armature now extends from $y=31$ to $46$ mm. The gap rectangles must change
with it; leaving the old 3 mm refinement rectangles would paint air over part
of the new armature.

There are no current-density declarations in this model. The selected materials
have zero source current. The alnico record supplies the coercivity and shifted
B--H table, so there is no need to guess an equivalent coil current or enter a
second remanence value. The angle $180^\circ$ points the prescribed magnetisation
to the left. With this orientation, flux rises through the left leg, crosses the
armature, and returns through the right leg.

The material law describes a prescribed static branch. Changing the drawing
does not simulate the magnetic history of a real alnico magnet or irreversible
demagnetisation. Those limits still apply even when the computed field lines
look convincing.

\subsection{Keeping one contour and bending another}

Maya uses the same field-picture command as before. Her straight profile needs
a new name and a new height, halfway through the narrower gaps:
\begin{codeexample}[code only]
\draw[femm/profile={name=magnet-gap,samples=101},
  violet,dashed,thick,-Stealth] (-48,30.5)--(48,30.5);
\end{codeexample}
She copies the previous axis and replaces |profile=coil-gap| by
|profile=magnet-gap| in its three plots. Distinct names let both models and
both graphs coexist in one document. The two devices have different gaps and
different excitation; their field magnitudes are not expected to match.

She also wonders how the field varies along a curved contour in the central
air space. The profile key works on a Bezier path too:
\begin{codeexample}[code only]
\draw[femm/profile={name=magnet-arc,samples=101,
  curve tolerance=.03},orange!85!black,dashed,thick,-Stealth]
  (-18,-8) .. controls (-22,22) and (22,22) .. (18,-8);
\end{codeexample}
Again this is inside the model picture, after all physical regions. The
contour is a measurement path chosen by Maya, not a computed field line.
For its graph she chooses arc length and components relative to the path:
\begin{codeexample}[code only]
\begin{tikzpicture}
\begin{axis}[width=13cm,height=5cm,grid=major,no markers,
  xlabel={Arc length $s$ (mm)},ylabel={Field (T)},
  legend style={at={(.5,1.03)},anchor=south,legend columns=3}]
  \femmplot[profile=magnet-arc,component=B,
    plot options={black,thick}]
  \addlegendentry{$|\mathbf B|$}
  \femmplot[profile=magnet-arc,component=Bt,
    plot options={orange!85!black}]
  \addlegendentry{$B_t$}
  \femmplot[profile=magnet-arc,component=Bn,
    plot options={teal!70!black}]
  \addlegendentry{$B_n$}
\end{axis}
\end{tikzpicture}
\end{codeexample}
The default abscissa is |s|, the distance travelled along the flattened path,
in millimetres. It is not the Bezier parameter. |Bt| is the tangential component
and |Bn| is the component along the left normal. Reversing the traversal
reverses those component signs at a common point. Cartesian components remain
available by using |Bx| or |By| instead.

If Maya needs a closed measurement contour, she can use a path ending in
|-- cycle| with the same profile key. The first position then appears again
at the last sample. At a sharp corner the tangent changes, so local components
can jump even in a smooth physical field. On a material interface the sampling
side is ambiguous; placing the contour inside the desired medium or using a
small signed |offset| makes the choice explicit.

\clearpage
\subsection{The permanent magnet, assembled}
The file \nolinkurl{examples/tutorial-alnico.tex} combines the device and both
profiles; its full source and output follow. The orange curve lies in the
central air space, whereas the violet
line crosses the pole gaps; they sample different parts of the same field.
\luafemmexample{tutorial-alnico}

\clearpage
\subsection{Keeping a calculation while editing the picture}
Maya is now adjusting labels and choosing colours. These edits do not change
her magnetic problem, so she adds one option to its picture:
\begin{codeexample}[code only]
\begin{tikzpicture}[
  femm/problem={xmin=-110,xmax=110,ymin=-95,ymax=115,mesh size=5},
  femm/cache={file=\jobname-field,mode=auto}]
  % The same material paths, field picture and measurement paths.
\end{tikzpicture}
\end{codeexample}
The first compilation writes the mesh and solution; subsequent compilations
reuse them. Maya can change field-line counts, profile sampling or PGFPlots
labels without solving again. Changing the current reuses the mesh but
requires a new solution; refining the mesh requires both steps again.
The two complete tutorial files enable this option. Distinct magnetic
problems in the same document need distinct cache filenames.

When the physical inputs are final, |mode=frozen| makes an absent or outdated
result an error instead of starting another calculation. It still allows
Maya to edit the presentation. It does not certify numerical accuracy:
the next step is precisely to investigate that accuracy. Full options,
file locations and diagnostics are in Section~\ref{sec:cache}.

\subsection{Finishing the picture: now refine the computation}
\label{sec:tutorial-refinement}

Maya is happy with the geometry, labels and axes. Now it is worth spending more
time on the finite-element solution. She first refines the global spacing,
then the gaps, comparing results after each change instead of assuming that
one particular setting is sufficient.

\begin{center}
\begin{tabular}{llll}
\toprule Setting & Working draft & First refinement & Further gap refinement\\
\midrule
Global |mesh size| & 5 & 3 & 2\\
Coil gap |mesh size| & 1 & .5 & .25\\
Alnico gap |mesh size| & .5 & .25 & .125\\
\bottomrule
\end{tabular}
\end{center}
All entries are in millimetres. These are suggested trials, not accuracy
prescriptions or guaranteed mesh-quality bounds. For example, the first
refinement of the coil model consists of changing these two settings:
\begin{codeexample}[code only]
% In the model picture options:
femm/problem={xmin=-110,xmax=110,
  ymin=-95,ymax=115,mesh size=3}
% In the style definition:
gap mesh/.style={femm/region={material=air,mesh size=.5}}
\end{codeexample}
The physical geometry, material laws and current density stay fixed during
this comparison. Each trial must construct a new model, normally by recompiling
the document. Changing a key after |\pic {femm field};| does not remesh the
model that has already been solved.

She compares the field at fixed positions inside both gaps and the curves as
a whole. She looks especially at the places where the graph changes rapidly.
Then, keeping the mesh spacings fixed, she enlarges the outer air domain
by moving all four boundaries outwards. This separate check tells her whether
the artificial $A_z=0$ boundary still influences the displayed region.
A small Newton residual alone answers neither question.

For a direct overlay in one document, Maya can draw two complete model pictures,
using names such as |coil-coarse| and |coil-fine| for the same physical contour.
After both pictures, she plots them on one axis:
\begin{codeexample}[code only]
\femmplot[profile=coil-coarse,component=By,abscissa=x,
  plot options={gray,dashed}]
\addlegendentry{Coarse mesh}
\femmplot[profile=coil-fine,component=By,abscissa=x,
  plot options={red}]
\addlegendentry{Refined mesh}
\end{codeexample}
This fragment assumes both profiles have been registered under those names.
Each keeps its own solved model. Merely changing |samples=101| to |samples=301|
on the coarse model would sample the same approximate solution more densely;
it would not produce the fine-mesh comparison.

Once the values of interest are stable enough for her purpose, Maya can increase
the contour levels or profile samples for the final presentation. She keeps
the mesh/domain sensitivity in mind and retains the source settings alongside
the figure. The workflow has been economical: ordinary TikZ while arranging
the device, a coarse solve while developing the explanation, and more triangles
when there is a numerical result worth checking carefully.

The remaining sections provide the complete key reference and numerical
conventions. This tutorial follows the progressive, task-based presentation
of the PGF tutorials; its narrative and magnetic examples are original.
\clearpage

% SPDX-License-Identifier: LPPL-1.3c
% Copyright (C) 2026 Christophe Jorssen
% Author and Current Maintainer: Christophe Jorssen <christophe.jorssen@gmail.com>
% LPPL maintenance status: maintained
% This file is part of luafemm. See LICENSE and LICENSES.md for its terms.
% Original tutorial following PGF's progressive, task-led presentation.
\clearpage
\section{Tutorial: A synchronous machine for Kamil}
\label{sec:tutorial-machines}

Kamil is preparing a lesson on synchronous machines. His first drawing has a
circular rotor, an iron stator and two small rectangles for the sides of a
coil. He can explain which way the current flows. What he cannot yet show is
why the field around the air gap should resemble the sine wave on the next
page of his notes. Two rectangles do not look particularly sinusoidal.

He decides to start with those two rectangles, calculate the field, and let
the graph tell him what needs changing. We shall follow him from a concentrated
winding to regularly spaced slots, then to weighted turns and deliberately
unequal slot spacing. Every complete source below is distributed with the
package and executed when this manual is built.

\subsection{Deciding what the drawing represents}

The drawing is a transverse section of a cylindrical machine. The nominal
air-gap surfaces are concentric circles, with no salient poles. Current flows
parallel to the machine axis, out of or into the page. The model has a constant
axial depth. This is a planar field calculation; it is not an axial-flux
machine with a gap between two disks.

Kamil first excites only the stator. The passive iron rotor remains essential:
it changes the magnetic circuit even though its current is zero. Each figure
represents a magnetostatic state. A conducting material does not acquire an
induced current simply because its conductivity is nonzero. End windings,
eddy currents and time-dependent motion are outside this example.

For these comparisons he chooses a linear teaching material with relative
permeability 2000. This avoids mixing winding-distribution effects with
saturation. Later he can replace its identifier with |pure-iron| or another
catalogue material and solve the nonlinear problem.

\subsection{A drawing that is inexpensive to change}

Kamil remembers Maya's first decision: use a modest global mesh while composing
the drawing. Here the global spacing is 4 mm, with 0.7 mm in an air-only band
inside the 2 mm gap. These are working settings, not certified accuracy bounds.
The 0.03 mm curve tolerance controls the polygonal representation of the
circular material interfaces independently of the mesh spacing.

He defines reusable TikZ styles. The component keys use model millimetres,
so |radius=28| is a physical radius; |scale=.65| on the picture only adjusts
the printed diagram. The physical geometry does not grow when he changes its
labels or selects a larger font.

\begin{codeexample}[code only]
\tikzset{
  femm/materials/teaching-iron/.style={mu r=2000},
  kamil rotor/.style={femm/rotor={
    radius=28,core material=teaching-iron}},
  kamil stator/.style={femm/stator={
    bore radius=30,outer radius=45,
    core material=teaching-iron,
    slots={count=2,start angle=90,
      depth=8,neck depth=.8,
      opening width=1,width=3.5,bottom width=4.5},
    coil={slots={1,2},ampere turns=200}}}
}
\end{codeexample}

A slot is a cavity, not a current source. The |coil| record associates two
conductor sections: slot 1 carries $+200$ A\,turns and slot 2 carries
$-200$ A\,turns. Positive current points out of the page. The conductors need
not have equal areas: their current densities are calculated from their
individual captured areas.

\subsection{Two nodes, a gap, and one circular measurement}

The order is worth writing down. Kamil declares the two nodes, adds the gap
refinement, registers the measurement, and finally asks for field lines.
The field request starts the solve. Afterwards he may plot as many components
of the registered profile as he likes, but he cannot add another physical
part to that solved model.

\begin{codeexample}[code only]
\node[kamil stator] (S) {};
\node[kamil rotor] (R) {};
\pic[femm/air gap={rotor=R,stator=S,mesh size=.7}]
  {femm air gap};
\pic[femm/gap profile={rotor=R,stator=S,
  name=gap,samples=721}] {femm gap profile};
\pic[femm/lines=19] {femm field};
\end{codeexample}

The gap profile lies halfway between the nominal surfaces. It samples an
exact circle, uniformly in mechanical angle. Its red drawing is only an
annotation; it does not impose a boundary condition.

The component Kamil wants is $B_r$, positive away from the common centre.
Plotting $|\mathbf B|$ would conceal the change of sign between north and
south poles. |Bn| is also a different convention: the left normal of a
counterclockwise circular path points inward, whereas $B_r$ points outward.

\begin{codeexample}[code only]
\begin{axis}[xlabel={Mechanical angle (degrees)},ylabel={$B_r$ (T)}]
  \femmplot[profile=gap,component=Br,abscissa=angle]
\end{axis}
\end{codeexample}

Here is his complete first document, including the calculated field, the
angular scan and its harmonic spectrum. Its profile is called |kamil-single|
to keep it distinct from the following examples.
\luafemmexample{tutorial-machine-single}

\subsection{More slots do not specify a winding}

Kamil now puts 24 slots around the stator. Their centres start at
$7.5^\circ$ and advance in steps of $15^\circ$. He keeps the opening and body
widths unchanged. With |distribution=uniform|, conductor sections in one
half of the machine carry equal positive ampere-turns, and those in the
opposite half carry equal negative ampere-turns. The selection is based on
the sign of $\sin(\theta)$ for |axis=0,pole pairs=1|.

\begin{codeexample}[code only]
slots={count=24,start angle=7.5,
  depth=8,neck depth=.8,opening width=1,
  width=3.5,bottom width=4.5},
winding={distribution=uniform,ampere turns=200}
\end{codeexample}

The total positive excitation remains 200 A\,turns, as in the concentrated
coil. It is not 200 A\,turns in every slot. This normalization makes the
comparison meaningful, although the fundamental amplitudes need not be equal.
The distributed source produces a different magnetomotive-force waveform.
It does not prescribe the resulting flux density.
\luafemmexample{tutorial-machine-uniform}

\subsection{Changing the turns before moving the slots}

Kamil tries |distribution=sine| with exactly the same slots. For each slot
centre $\theta_i$, the source weight is
\[
 w_i=\sin\bigl(p(\theta_i-\theta_0)\bigr),\qquad
 (NI)_i=(NI)_+\frac{w_i}{\sum_j\max(w_j,0)}.
\]
Here $p$ is |pole pairs| and $\theta_0$ is |axis|. The distribution must be
balanced: $\sum_i w_i=0$ within numerical tolerance. The package rejects an
unbalanced set instead of inventing a return conductor. Multiple layers
share a slot's assigned ampere-turns equally.

In the usual thin-gap, highly permeable-iron approximation, the spatial
variation of magnetomotive force is related to the current distribution.
A sinusoidal current sheet therefore motivates this construction. The FEM
calculation still includes the finite permeability and the actual slots.
This is why the graph need not be a perfect sinusoid. For the continuous
current-sheet model, see the synchronous-machine example in
\href{https://web.mit.edu/6.013_book/www/chapter11/11.7.html}{Haus and Melcher,
\emph{Electromagnetic Fields and Energy}, Section 11.7}.

\begin{codeexample}[code only]
winding={distribution=sine,ampere turns=200,
  pole pairs=1,axis=0}
\end{codeexample}

These weights describe homogenised ampere-turns and can imply noninteger
turn counts at a chosen current. For a manufacturable winding, Kamil would
choose integer turns with explicit |coil| records, then recalculate the
harmonics. No turn rounding is hidden inside |distribution=sine|.
\luafemmexample{tutorial-machine-sine}

\subsection{Keeping equal turns and choosing the positions}

Kamil also wants a construction with identical coil sides. He retains twelve
positive and twelve negative sections, but changes their positions. In the
positive half, he places slot centres at the midpoint quantiles of a density
proportional to $\sin\theta$:
\[
 \theta_i=\arccos\left(1-2\frac{i-1/2}{12}\right),
 \qquad i=1,\ldots,12.
\]
The remaining positions are $\theta_i+180^\circ$. Pairing each slot with its
opposite gives equal positive and negative ampere-turns. The slots are denser
near $90^\circ$ and $270^\circ$, where the continuous source is strongest.

The source below contains the resulting explicit angles. The geometric
checks still apply: adding slots indefinitely would eventually leave too
little room for the teeth. The construction approximates a distribution;
it is not an optimizer and does not guarantee cancellation of particular
harmonics. Its spectrum is the test of the idea.
\luafemmexample{tutorial-machine-positioned}

\subsection{Asking how close is close enough}

The spectrum reports $B_n/B_1$ through order 15. For a component expanded as
$B(\theta)=a_0+\sum_n(a_n\cos n\theta+b_n\sin n\theta)$, the reported amplitude
is $B_n=\sqrt{a_n^2+b_n^2}$. The spatial distortion shown below each example is
\[
 \mathrm{THD}_{15}=\frac{\sqrt{\sum_{n=2}^{15}B_n^2}}{B_1}.
\]
It excludes DC and harmonics above order 15. In particular, a small value
does not imply that slot-scale ripple at higher orders has disappeared.
The output is a ratio; multiply by 100 for percent.

Before using the final figures, Kamil repeats the calculations with global
spacing 2 mm, gap spacing 0.35 mm, and a smaller curve tolerance. He compares
fundamental amplitude, phase, and selected harmonics. He then checks a larger
outer air domain to assess the imposed outer boundary. Increasing |samples|
only improves the scan and Fourier quadrature; increasing |femm/lines| only
adds drawn contours. Neither refines the finite-element solution.

He can finally replace |teaching-iron| with a nonlinear material, change the
stator to a square outer shape, or excite a slotted rotor. Each change is a
new physical comparison. Styles let him keep the same labels, measurement
commands and PGFPlots presentation throughout.

% SPDX-License-Identifier: LPPL-1.3c
% Copyright (C) 2026 Christophe Jorssen
% Author and Current Maintainer: Christophe Jorssen <christophe.jorssen@gmail.com>
% LPPL maintenance status: maintained
% This file is part of luafemm. See LICENSE and LICENSES.md for its terms.
% Original task-led tutorial; every complete source is compiled with the manual.
\clearpage
\section{Tutorial: Kamil looks for a defect without cutting the tube}
\label{sec:tutorial-ndt}

Kamil has a copper tube, a small coil and a new question for his students.
Can the coil tell him that some metal is missing from the tube wall? He would
like to show the experiment in three stages: a drawing of the apparatus,
a field picture explaining the interaction, and the signal that an instrument
could actually measure. A coloured defect in a drawing is easy to find.
The more interesting part is finding it in the coil's electrical response.

We shall build an encircling-coil inspection of a tube with a circumferential
groove. We then move the probe, compare each position with a sound reference,
and plot the two components of the impedance change. The three complete
examples below contain their own preambles and styles. They use ordinary
TikZ paths, the existing |femm| keys and PGFPlots; no document-side Lua,
native solver or external data file is required.

\subsection{First decide what the defect means}
A meridional $(r,z)$ section of a circular tube is a rectangle. Revolving a
small rectangular recess in its outer edge makes an \emph{annular groove}.
This is a useful model of circumferential material loss. It is also compatible
with the axisymmetric harmonic solver: the coil, specimen and defect all
retain rotational symmetry throughout the scan.

A notch on just one side of the tube, an off-centre probe or a longitudinal
slit would break that symmetry. Drawing one of those in this meridional
picture would still revolve it into a ring. Kamil therefore labels his defect
as a groove and keeps the coil coaxial. The remaining azimuthal currents still
form closed circles; the groove changes the available conducting volume and
its coupling to the coil. It does not force these circles to detour around a
local three-dimensional crack.

The apparatus uses the following dimensions. Coordinates and local mesh
spacings are in millimetres, while conductivity is in SI units.
\begin{center}
\begin{tabular}{ll}
\toprule Setting & Value\\
\midrule
Tube inner / outer radius & $5/8\,\mathrm{mm}$\\
Tube axial extent & $-25\leq z\leq25\,\mathrm{mm}$\\
Groove width / centre & $3\,\mathrm{mm}$ / $z=0$\\
Groove depths to compare & $0.5$ and $1.5\,\mathrm{mm}$\\
Coil radial extent & $9\leq r\leq11\,\mathrm{mm}$\\
Coil axial extent & $s-1.5\leq z\leq s+1.5\,\mathrm{mm}$\\
Coil turns / copper fill & $60$ / $0.6$\\
Excitation & $f=2\,\mathrm{kHz}$, $\widehat I=1\,\mathrm A$ peak\\
Copper & $\mu_r=1$, $\sigma=58\,\mathrm{MS/m}$\\
Air-domain truncation & $0\leq r\leq45$, $-45\leq z\leq45\,\mathrm{mm}$\\
\bottomrule
\end{tabular}
\end{center}
In the apparatus drawing, two orange rectangles are opposite sections of the
\emph{same} encircling winding. The solver's meridional drawing needs only the
positive-radius rectangle. The transverse view shows that the missing metal
extends around the whole circumference. This first drawing is purely
geometrical; the following two examples calculate the fields and signals.
\luafemmexample{tutorial-ndt-setup}

\subsection{Prepare a model that is inexpensive to edit}
Kamil starts with a coarse field mesh while placing labels and deciding how
much of the apparatus to show. The far air needs fewer triangles than the
conducting wall or the narrow groove. At $2\,\mathrm{kHz}$, copper's standard
skin depth is approximately
\[
 \delta=\sqrt{\frac{2}{\omega\mu_0\sigma}}\simeq1.48\,\mathrm{mm}.
\]
He can try |\Mesh=.6| while composing the example, then use |.4| for the
published pictures and compare with |.25| before interpreting small signal
changes. The groove receives half this spacing. These are resolution choices,
not a built-in adaptive skin-depth estimator.

He puts the common physical settings in a TikZ style:
\begin{codeexample}[code only]
\def\Mesh{.4}\def\Frequency{2000}
\pgfmathsetmacro{\NearMesh}{1.5*\Mesh}
\pgfmathsetmacro{\GrooveMesh}{\Mesh/2}
\tikzset{ndt/model/.style={
  femm/problem={type=axisymmetric,frequency=\Frequency,
    xmax=45,ymin=-45,ymax=45,mesh size=2},
  femm/circuit={name=probe,current=1},
  femm/materials/specimen/.style={mu r=1,conductivity=58000000}}}
\end{codeexample}
The two |\pgfmathsetmacro| calls evaluate expressions before passing their
numeric values to the mesh keys. The |ndt/| styles are local example styles,
not another installed package interface. They can be edited and composed with
normal PGF |.style| and |.append style| handlers.

|type=axisymmetric| is essential here. It solves the finite probe and its
surrounding air. The |type=axial| long-solenoid reduction of
Section~\ref{sec:currents} prescribes an axial field and has no probe terminal
impedance to measure. The outer boundary in this tutorial is the usual zero
vector-potential truncation. Enlarging the air domain and refining the wall
mesh are separate checks.

\subsection{Paint copper, remove metal with air, add the winding}
Kamil first paints a locally refined air band, then the complete copper wall.
The order matters: painting the air band last would erase the specimen in
that band. The groove is a later \emph{air region}, replacing part of the wall.
It is not an excluded mesh hole. Magnetic field must still be solved in the
space left by the removed metal.
\begin{codeexample}[code only]
\path[femm/region={material=air,mesh size=\NearMesh}]
  (0,-15) rectangle (12,15);
\filldraw[fill=blue!10,
  femm/region={material=specimen,mesh size=\Mesh}]
  (5,-25) rectangle (8,25);
% Omit this path for the sound reference:
\filldraw[fill=white,draw=red,
  femm/region={material=air,mesh size=\GrooveMesh}]
  (6.5,-1.5) rectangle (8,1.5);
\filldraw[fill=orange!25,
  femm/region={material=copper,circuit=probe,turns=60,
    winding fill=.6,mesh size=\Mesh}]
  (9,-1.5) rectangle (11,1.5);
\end{codeexample}
The specimen is a solid conductor and develops resolved eddy currents.
The named winding is stranded: it represents insulated fine wires carrying
the prescribed current, with homogenised copper resistance. It does not become
a massive copper ring just because its region uses |material=copper|.
The specified 1 A is a peak phasor, so the terminal power uses the factor $1/2$.
All sources follow $e^{+\mathrm i\omega t}$.

\subsection{Look at the field, then at the current}
The example draws magnetic flux lines in blue and normal-current symbols in
red. Both pictures use $\omega t=0$, so their arrows and symbols describe the
same instant. In the meridional section, positive $J_\varphi$ points into the
page and is drawn as a cross; a dot points out of the page. Its actual path is
a circle about the axis, not a loop lying in the $(r,z)$ picture.
\begin{codeexample}[code only]
\pic[femm/phase=0,femm/lines=17,
  femm/field style={blue!65!black},femm/arrow spacing=11mm]
  {femm field};
\pic[femm/currents={material=specimen,phase=0,spacing=.9,
  reference current=2000000,symbol radius=.9pt},
  femm/eddy style={red!70!black}] {femm eddy currents};
\end{codeexample}
The fixed current reference is in $\mathrm{A/m^2}$; it gives both panels the
same opacity scale. Magnetic contour levels are chosen separately in each
picture, so line counts are not a quantitative measure of the difference.
The view is clipped to a close-up while the complete tube and air domain
remain in the calculation. A TikZ |\clip| changes the visible drawing; region capture still passes
the complete, unclipped physical path to the solver.

A second comparison follows the current along a line inside the remaining
wall. The line lies at $r=6.25\,\mathrm{mm}$, below the floor of even the
$1.5\,\mathrm{mm}$ groove. Thus it never jumps into air at the defect:
\begin{codeexample}[code only]
\path[femm/profile={name=ndt-groove,samples=161}]
  (6.25,-10)--(6.25,10);
% Later, inside a PGFPlots axis:
\femmplot[profile=ndt-groove,component=Jeddy,abscissa=z,
  representation=amplitude,plot options={red!70!black,thick}]
\end{codeexample}
The pic is instantaneous, whereas the curve is a phasor amplitude. In the
complete example the axis converts $\mathrm{A/m^2}$ to $\mathrm{MA/m^2}$
with a normal PGFPlots |y filter|. The deeper groove concentrates current in
part of the remaining wall: the local curve may rise even while the integrated
specimen loss falls. There is no rule that every local current must decrease
when metal is removed. The plotted differences come from the solved fields.
\luafemmexample{tutorial-ndt-fields}

\subsection{What does the instrument measure?}
Kamil cannot put a current-density probe inside every triangle of the tube.
He can, however, measure the voltage needed to drive the coil. Its complex
impedance is
\[
 \widehat V=R_{\rm coil}\widehat I+\mathrm i\omega\widehat\lambda,
 \qquad Z=\frac{\widehat V}{\widehat I}=R+\mathrm iX.
\]
The flux linkage includes the reaction of the induced specimen currents.
The existing circuit measurement supplies both components:
\begin{codeexample}[code only]
\femmmeasure[name=probe-reading,quantity=circuit,circuit=probe]
% Expandable SI numbers, suitable for subsequent PGF calculations:
\femmresult{probe-reading}{impedance_real}
\femmresult{probe-reading}{impedance_imaginary}
\end{codeexample}
|impedance_real| is the total AC resistance seen at the terminals;
|resistance| alone is the winding's DC copper resistance. Kamil must retain
|impedance_real| to see specimen losses. For this one-coil, passive-conductor
model,
\[
 \tfrac12|\widehat I|^2\bigl(\mathop{\rm Re}Z-R_{\rm coil}\bigr)
 =P_{\rm specimen}.
\]
The automated tutorial test checks this identity independently of the drawing.
The coil itself is unchanged between the sound and grooved specimens, so
its copper resistance cancels in the difference.

This use of the coil's impedance is the usual eddy-current inspection
principle: a change in the specimen changes induced currents and their
magnetic reaction on the probe. The
\href{https://infotechnology.fhwa.dot.gov/eddy-current-testing/}{FHWA introduction
to eddy-current testing} describes both impedance-plane and distance-based
signal displays, as well as the influence of probe spacing and material
properties. Our example calculates a particular idealised groove response;
it does not import an instrument trace or impose a chosen signal shape.

\subsection{Move the probe and subtract a matched reference}
Let $s$ be the axial coordinate of the coil centre. Only the winding rectangle
moves:
\begin{codeexample}[code only]
\path[femm/region={material=copper,circuit=probe,turns=60,
  winding fill=.6,mesh size=\Mesh}]
  (9,{\Position-1.5}) rectangle (11,{\Position+1.5});
\femmmeasure[name=\Case-\Index,quantity=circuit,circuit=probe,
  parameter=\Position]
\end{codeexample}
Every position starts a new |tikzpicture| and declares its complete model.
Do not shift the winding after solving or add a groove after a field pic has
already triggered meshing. The scan is a sequence of steady harmonic states;
no velocity or motion-induced current term is used. Physically this describes
measurements after the electromagnetic transient has settled at each position.

For every $s$, solve the sound tube and both groove depths, then form
\[
 \Delta Z(s)=Z_{\rm groove}(s)-Z_{\rm sound}(s),\qquad
 \Delta R=\mathop{\rm Re}\Delta Z,\quad
 \Delta X=\mathop{\rm Im}\Delta Z.
\]
The sound reference is evaluated at the \emph{same position}. A finite tube
can change its coupling as the coil approaches an end; subtracting a single
reference value for the whole scan would mix that change into the defect
signal. This is a numerical reference subtraction, not a model of a physical
two-coil differential probe or a bridge circuit.

The full source uses a named TeX box to execute each model without printing
57 intermediate pictures. Named measurements remain available after the box
is discarded. Each case/position has its own native cache. On a later
compilation, changing graph styles or labels restores the solved phasors.
Changing physical geometry, frequency or mesh settings invalidates the
appropriate caches automatically.

The difference is small compared with the total impedance. Use PGF's
floating-point unit when calculating it, and convert to milliohms afterward:
\begin{codeexample}[code only]
\begingroup
  \pgfkeys{/pgf/fpu=true,/pgf/fpu/output format=fixed}
  \pgfmathparse{1000*(\femmresult{deep-10}{impedance_real}
    -\femmresult{sound-10}{impedance_real})}
  \xdef\CentreDeltaR{\pgfmathresult}
\endgroup
\end{codeexample}
In the distributed position list, index 10 is $s=0$. The floating-point keys
are deliberately grouped: they must not change the arithmetic used by later
TikZ geometry. The complete example collects coordinate lists with ordinary
TeX macros and draws them with |\addplot|. No hand-written Lua is necessary.
\luafemmexample{tutorial-ndt-scan}

\subsection{Read the three graphs together}
The left graph is the resistive channel. For these parameters it has a
negative central response and smaller positive shoulders. The middle of the
groove gives lower total specimen loss, while the off-centre response also
contains redistribution and changed coupling. A generic ``defect signal''
should not be substituted for this computed shape.

The reactive channel has a positive central response. The deeper groove has
a larger complex signal, although doubling groove depth would not in general
double either component. The impedance-plane graph uses the same computed
points, with $\Delta R$ horizontally and $\Delta X$ vertically. It does not
plot the temporal sine wave. The geometry is symmetric under $s\mapsto-s$,
so the two halves of the scan nearly retrace the same trajectory. Small
separations between them reflect the discretisation, not hysteresis.

The markers identify actual solves. Lines join adjacent samples without
claiming an interpolated extremum or adding a fitted defect model. Use more
positions near a narrow response to improve spatial sampling; use more
triangles to improve field accuracy. These are different operations.

\subsection{Before interpreting a small difference}
Kamil has three independent resolution controls:
\begin{enumerate}
\item Reduce local |mesh size| in the metal, groove and nearby air. Compare
$\Delta Z$ itself: a small change in the much larger total $Z$ does not by
itself establish convergence of the difference.
\item Enlarge the outer air domain while retaining the same local mesh
settings. The zero-potential boundary is a truncation, not a real shield.
For this example, the initially tried $25\,\mathrm{mm}$ radial extent changed
the central differential signal by about $2.7\%$ when enlarged to $45$ mm,
which is why the complete sources use the larger domain.
\item Refine the list of probe positions independently of the field mesh.
The source places extra points near the groove, where the signal changes most.
\end{enumerate}
The regression |tests/test_ndt.lua| checks the two depth responses, the power
identity, mesh refinement, outer-domain enlargement, scan symmetry and the
small far-from-groove signal. Its recorded results and tolerances are in
\nolinkurl{docs/validation.md}. They establish numerical sensitivity for this
teaching geometry, not an experimental detection threshold.

Finally, a change in impedance is not uniquely a change in groove depth.
Conductivity, frequency and coil clearance also change the coupling. In a
real inspection these require calibration; the present calculation includes
no instrument noise or uncertainty distribution. In an encircling-coil model,
increasing the coil radius changes its circumference and therefore its copper
resistance as well as its clearance. Keep that distinction when extending the
exercise. Lateral wobble and isolated cracks require a three-dimensional model.
The useful result here is a reproducible chain from a declared material defect
to fields, currents and a measurable complex coil response.

\clearpage

% SPDX-License-Identifier: LPPL-1.3c
% Copyright (C) 2026 Christophe Jorssen
% Author and Current Maintainer: Christophe Jorssen <christophe.jorssen@gmail.com>
% LPPL maintenance status: maintained
% This file is part of luafemm. See LICENSE and LICENSES.md for its terms.
\section{Problems, coordinates and physical paths}
\label{sec:model}
\begin{key}{/tikz/femm/problem=\marg{options} (default \normalfont empty)}
Create one model at the beginning of a |tikzpicture|. Put this key in the
picture options, not on a path or a nested scope. It installs |x=1mm,y=1mm|
and records the initial affine frame. There is one current model; nested
models inside the same picture are unsupported. Separate pictures create
independent models.
\end{key}

The confined |field model=ideal| and its domain membership keys are documented
in Section~\ref{sec:ideal}. The defaults below describe |field model=open|.

The following subkeys belong to |/femm/problem| and are normally supplied
inside |femm/problem={...}|. All lengths in this list are numbers in
\emph{model millimetres}, not graphic dimensions.

\begin{key}{/femm/problem/xmin=\meta{number} (initially -90)}
Left boundary of the rectangular air domain or ideal-mode staging mesh.
\end{key}
\begin{key}{/femm/problem/xmax=\meta{number} (initially 90)}
Right boundary; must exceed |xmin|. All material vertices must lie within
the domain, including its boundary.
\end{key}
\begin{key}{/femm/problem/ymin=\meta{number} (initially -75)}
Bottom boundary of the air domain.
\end{key}
\begin{key}{/femm/problem/ymax=\meta{number} (initially 95)}
Top boundary; must exceed |ymin|. By default, the potential is zero on the
entire outer rectangle. Use |\femmboundary| to change individual sides
(Section~\ref{sec:boundaries}). A finite exterior boundary is an approximation,
so enlarge the domain when testing the sensitivity of a result.
\end{key}
\begin{key}{/femm/problem/mesh size=\meta{positive number} (initially 2)}
Global target spacing of interior points and maximum length used to
subdivide input constraints. It is not a bound on every triangle diameter
or area. Smaller values increase both memory and solution time.
\end{key}
\begin{key}{/femm/problem/mesher=\meta{delaunay or grid} (initially delaunay)}
|delaunay| uses the constrained Lua triangulator and accepts general
polygonal interfaces. |grid| uses the structured algorithm,
which requires axis-aligned region edges. Local mesh sizes apply only to
|delaunay|. The standalone Lua constructor defaults to |grid|; the TikZ interface
explicitly selects |delaunay|.
\end{key}
\begin{key}{/femm/problem/mesh tolerance=\meta{nonnegative number} (initially 0)}
Vertex-merging distance. Zero selects $10^{-12}$ times the greater domain
dimension. A positive value must be below one percent of the global mesh
size. Features at this scale may merge. This tolerance controls geometric
construction, not the exact signs of the robust predicates.
\end{key}
\begin{key}{/femm/problem/curve tolerance=\meta{positive number} (initially .03)}
Default chord-distance tolerance for flattening material curves. It is
copied to the region options when the picture starts. It does not change
the original TikZ drawing and is distinct from mesh spacing. Measurement
profiles have a separate |curve tolerance| option.
\end{key}
\begin{key}{/femm/problem/minimum angle=\meta{number} (initially 0)}
Optional quality target in degrees, between 0 and 30 inclusive; zero disables
the angle criterion. Requires |mesher=delaunay| when nonzero. After constructing
the initial constrained mesh, split encroached segments and insert triangle
circumcentres until every angle meets the target, within floating-point audit
tolerance. Start with 20 degrees. A request that cannot be completed raises an
error; acute input corners may make the target impossible. See
Section~\ref{sec:mesh-quality} for the algorithm, limits and a complete example.
\end{key}
\begin{key}{/femm/problem/maximum area=\meta{nonnegative number} (initially 0)}
Optional upper bound on each triangle's area, in model millimetres squared.
Zero disables the area criterion. A nonzero value requires |mesher=delaunay|
and enables quality refinement even when |minimum angle=0|. Both criteria
apply together when enabled. This bound applies to the entire rectangular
domain, including background air; it is independent of local |mesh size|.
The Lua counterpart uses model-coordinate units squared, not SI units.
\end{key}
\begin{key}{/femm/problem/max Steiner points=\meta{nonnegative integer} (initially 5000)}
Maximum number of vertices added by quality refinement, counting both segment
midpoints and circumcentres. Initial constraint subdivision and lattice seeds
do not count towards this budget; the overall limit of 40000 mesh vertices
still applies. Zero permits no additional vertices. Exhausting the budget
raises an error, rather than accepting an under-refined mesh. The key has no
effect when both quality criteria are zero. Quality targets and this budget
are part of the persistent mesh descriptor; changes invalidate a saved mesh.
\end{key}
\begin{key}{/femm/problem/tolerance=\meta{positive number} (initially 1e-7)}
Target relative Newton residual, normalised by the initial residual. This
is an algebraic stopping criterion, not an estimate of the physical error.
\end{key}
\begin{key}{/femm/problem/max iterations=\meta{positive integer} (initially 60)}
Maximum number of Newton corrections. Nonconvergence raises an error;
the package does not silently use the last iterate.
\end{key}

\subsection{Declaring regions}
\begin{key}{/tikz/femm/region=\marg{options} (default \normalfont empty)}
Attach physical properties to one evaluated path with one or more closed
subcontours. Use |\path|,
|\draw| or |\filldraw| with |rectangle|, |circle|, |ellipse|, or |-- cycle|.
The capture happens after TikZ evaluates coordinates, transformations,
arcs, Bézier curves and rounded corners. Nodes on the path are annotations
and do not create additional regions. Stroke width does not add material.
\end{key}

\begin{key}{/femm/region/material=\meta{choice} (initially air)}
Select one of the 246 catalogue identifiers or a previously declared
custom material style. This is a PGF |.is choice| handler. Unknown values
produce an error, except that a declared |femm/materials/|\meta{name}
style is accepted as a custom choice. Catalogue identifiers are reserved.
\end{key}
\begin{key}{/femm/region/current density=\meta{expression or material} (initially material)}
Uniform signed $J_z$ in A/m$^2$. PGF expressions are accepted, including
|6000/(6*26)*1000000| for 6000 ampere-turns over a $6\times26$ mm$^2$
section. Positive current points out of the page. |material| inherits
the catalogue source; for most materials this is zero. Use an explicit
zero for a passive conductor. Currents are not renormalised if a path is
resized. The return section of a winding normally has the opposite sign.
\end{key}
\begin{key}{/femm/region/magnetization angle=\meta{expression} (initially 0)}
Prescribed direction in degrees from $+x$ towards $+y$, evaluated by PGF.
This is a direction in the model frame: rotating a local region path does
not automatically rotate its magnetisation. Rotating the entire picture
rotates both the rendered geometry and the displayed field.
\end{key}
\begin{key}{/femm/region/mesh size=\meta{nonnegative number} (initially 0)}
Local spacing inside the filled region and along every contour edge. Zero inherits the
global spacing; a positive value refines but does not coarsen that value.
Where refinement regions overlap, the smallest spacing wins. The region
still paints a material: declare a broad air refinement region before
the solids it contains, otherwise it replaces their material.
\end{key}
\begin{key}{/femm/region/curve tolerance=\meta{positive number}}
Override the picture's chord tolerance for this material path. The
initial value is copied from the problem at picture start.
\end{key}

\begin{keylist}[/tikz/femm]{material=\meta{choice},current density=\meta{expression},
  magnetization angle=\meta{expression},mesh size=\meta{number},curve tolerance=\meta{number}}
Aliases of the corresponding |/femm/region| subkeys. They can be inherited
from scopes or combined with graphics in a reusable TikZ style.
The alias |femm/mesh size| on a path means \emph{local} spacing; set the
global spacing inside |femm/problem|.
\end{keylist}

\begin{codeexample}[width=5cm]
\tikzset{passive/.style={
  femm/region={material=copper,
    current density=0},
  draw,fill=orange!25}}
\begin{tikzpicture}[femm/problem={
  xmin=-12,xmax=12,ymin=-12,ymax=12,
  mesh size=2}]
  \path[femm/region={material=air,
    mesh size=.6}]
    (-6,-5) rectangle (6,5);
  \filldraw[passive,rotate=15]
    (-4,-3) rectangle (4,3);
  \pic {femm mesh};
\end{tikzpicture}
\end{codeexample}

\subsection{Geometry and overlap rules}
A region may contain several individually simple closed contours.
Self-intersections, self-contact and backtracking within one contour are
errors. Native |even odd rule| and |nonzero rule| determine which portions
of a compound path receive its material; see Section~\ref{sec:domains}.
Distinct regions may overlap or intersect: \emph{the last declared
region whose fill contains a point determines material and current}.
The mesher splits all contour interfaces first, then classifies triangles
by their centroids. Air cavities remain part of the computed domain.

Simple node shapes with one closed background contour can also carry
|femm/region|. Their ordinary TikZ text sizing affects that contour.
For components whose dimensions must be independent of text, use the
node interfaces in Section~\ref{sec:components}. A pic may group several
physical paths, each with its own region declaration.

Use |femm/region| only on physical paths, not as |every path|. A clip
changes the visible part of a field or mesh, never the computed domain.
For profile paths, use |femm/profile| instead of |femm/region|.

The frame of the original picture is removed during capture and
reapplied to generated field and mesh graphics. Consequently |scale=2|
on the whole picture is a display change. A local scale inside the
picture changes the physical geometry. Fixed graphic dimensions, such
as the radius of |rounded corners|, retain TikZ's ordinary meaning.
Captured coordinates are rounded to $0.0001$ model unit, and points within
$0.001$ unit of a domain side are snapped to that side. Finer geometric
features cannot be recovered by reducing the mesh tolerance.

\subsection{Complete example: a nonlinear electromagnet}
This explicit-path example assembles a U core, its armature and the two
sections of a winding. Material and drawing options are collected into
TikZ styles. The displayed field and numeric diagnostics come from the
same physical paths; their declarations all precede the solve request.
\luafemmexample{u-electromagnet}

\subsection{Complete example: curved and transformed interfaces}
The next source rounds the core corners and tilts the armature with native
TikZ options. The mesh overlay exposes the flattened interfaces used by
the solver. Graphic clipping limits the displayed area without removing
any part of the computational domain.
\luafemmexample{curved-core}

% SPDX-License-Identifier: LPPL-1.3c
% Copyright (C) 2026 Christophe Jorssen
% Author and Current Maintainer: Christophe Jorssen <christophe.jorssen@gmail.com>
% LPPL maintenance status: maintained
% This file is part of luafemm. See LICENSE and LICENSES.md for its terms.
\section{Controlling mesh quality}
\label{sec:mesh-quality}
A smaller |mesh size| puts more points into the model, but it does not
necessarily improve the shape of every triangle. A narrow air gap or the
transition from a fine material region to coarse air can leave long, thin
elements. The optional |minimum angle| and |maximum area| keys address
triangle shape and area separately from nominal point spacing.

\subsection{Draft first, request quality second}
Keep both quality targets at zero while positioning parts and checking their
material declarations. Inspect the mesh, then add a reusable style with a
modest angle target, usually 20 degrees. Declare that style in the picture
options, together with the domain and spacing: refinement is part of meshing,
so it must be configured before the first mesh or field request.

The complete example below uses the same U electromagnet and 8 mm draft
spacing twice. Only the angle target changes. The two |femm/problem| keys
provided by the styles configure one model per picture; their options compose
before the model is created. |\pic {femm mesh};| performs meshing and quality
refinement, but does not solve the magnetic field.

For these sources, the minimum angle increases from about $7.29^\circ$ to
$20.00^\circ$ with 148 added points. The diagnostics are computed by the
example, rather than written into the figure. Counts are observations for
this geometry and revision, not promises for other devices.
\luafemmexample{mesh-quality}

\subsection{How the refinement proceeds}
The initial mesh still follows global and local point spacing. When either
quality target is nonzero, refinement first examines constrained segments.
A neighbouring triangle's opposite vertex encroaches on a segment if it lies
strictly inside the circle whose diameter is that segment. Such a segment
is bisected before processing deficient triangles. Points exactly on the
circle do not trigger subdivision; the disk test uses a filtered dot product
with exact-expansion fallback.

Deficient triangles are processed in order of shortest edge, with deterministic
tie breaking. Their proposed circumcentres are checked against the constrained
segments. If a candidate encroaches a segment, that segment is split instead.
The point-location walk also stops at a constrained interface and requests
its subdivision before continuing. Accepted insertions are followed by local
Lawson flips that preserve every constrained edge. Newly created triangles
and neighbouring segments are checked again. All split segments retain their
material-interface role or their exterior tag for boundary conditions.

The final audit independently checks the requested angle and area bounds,
constraint preservation, coverage, topology and constrained-Delaunay legality.
Successful refinement therefore meets the requested criteria, within the
documented floating-point audit margins ($10^{-8}$ degrees and $10^{-10}$
relative area). It does not estimate the magnetic field error. Check field
convergence with spacing and exterior-domain changes as before.

\subsection{Limits and failure diagnostics}
This is an original, bounded Lua implementation of segment-first Delaunay
refinement inspired by the algorithms studied in Triangle. It is not a full
port of Triangle's refinement machinery. In particular, it has no specialised
protection for acute input corners, no coarsening and no guaranteed smooth
gradation between local sizes. Candidate disk checks are conservative across
all segments, which can request more subdivisions than a visibility-aware test.

A material wedge with a smaller angle than the requested target cannot be
triangulated while satisfying that target everywhere. Closely spaced features
can require many extra points. Splitting stops with an error when the Steiner
budget, the total vertex cap or the geometric tolerance prevents further
progress; a candidate indistinguishable from an existing vertex also raises
an error. No partial result is silently accepted or saved as a successful mesh.
Reduce the target or review the geometry first; increase the budget only when
the extra resolution is intentional. These controls are shared by all three
TeX formats and use no external meshing program.

% SPDX-License-Identifier: LPPL-1.3c
% Copyright (C) 2026 Christophe Jorssen
% Author and Current Maintainer: Christophe Jorssen <christophe.jorssen@gmail.com>
% LPPL maintenance status: maintained
% This file is part of luafemm. See LICENSE and LICENSES.md for its terms.
\section{Nested domains and compound paths}
\label{sec:domains}
Two rectangles inside one another can describe a core and its surrounding
material. There are two complementary ways to express this in TikZ:
paint the outer region first and the inner region second, or give the
surrounding material an explicit hole. Both use the same conforming mesh.

\subsection{Two ordinary paths: later declarations win}
With separate region paths, the physical declaration order matters.
The last region containing a point supplies its material, current density
and magnetisation direction. A broad outer region therefore comes first:
\begin{codeexample}[width=5cm]
\begin{tikzpicture}[femm/problem={
  xmin=-25,xmax=25,ymin=-20,ymax=20,
  mesh size=3}]
  \filldraw[femm/region={
    material=pure-iron},fill=gray!25]
    (-20,-15) rectangle (20,15);
  \filldraw[femm/region={
    material=cast-alnico-5-lng37,
    magnetization angle=90},fill=red!25]
    (-10,-6) rectangle (10,6);
  \pic {femm mesh};
\end{tikzpicture}
\end{codeexample}
The iron occupies the space between the rectangles; alnico occupies the
inner rectangle. Reversing these two declarations would replace the
alnico with iron. Graphical opacity, clipping and drawing layers do not
perform material subtraction. An explicit later |material=air| region
paints air in exactly the same way.

\subsection{A surrounding material with its own hole}
To express the annulus directly, use one physical path with two closed
subpaths and TikZ's native |even odd rule|. No additional luafemm fill-rule
key is needed. The same rule controls the visible fill and the material:
\begin{codeexample}[width=5cm]
\tikzset{iron shell/.style={
  even odd rule,
  femm/region={material=pure-iron},
  draw,fill=gray!25}}
\begin{tikzpicture}[femm/problem={
  xmin=-25,xmax=25,ymin=-20,ymax=20,
  mesh size=3}]
  \filldraw[femm/region={
    material=cast-alnico-5-lng37,
    magnetization angle=90},fill=red!25]
    (-10,-6) rectangle (10,6);
  \path[iron shell]
    (-20,-15) rectangle (20,15)
    (-10,-6) rectangle (10,6);
  \pic {femm mesh};
\end{tikzpicture}
\end{codeexample}
Here the magnet was declared first. It survives because the iron path
does not cover the cavity. The cavity preserves the previously declared
material; if none was declared there, it contains the background air.
A later material may still paint over any part of this compound region.

\subsection{The native TikZ fill rules}
\begin{key}{/tikz/even odd rule}
For a physical region, count how many contour boundaries a ray from the
point crosses. Odd counts fill the region, even counts leave it unchanged.
Contour direction and the order of subpaths within this one path do not
matter. With three nested contours, the outer band and innermost island
receive the material, while the intermediate cavity preserves its previous
material. Disjoint contours describe separate pieces with common properties.
\end{key}
\begin{key}{/tikz/nonzero rule}
PGF's default rule uses a signed winding number. A point is filled when
that number is nonzero. Nested contours with the same direction therefore
fill the centre too; reverse the inner contour to make a hole. For nested
rectangles or circles, |even odd rule| is usually easier to read. A native
fill-rule key on a path overrides an inherited rule in the ordinary TikZ way.
\end{key}
These are existing PGF/TikZ keys, documented here for their additional
physical effect. luafemm reads the evaluated PGF fill rule during capture.
It does not depend on whether the path also requests graphical filling:
|\path[even odd rule,femm/region=...]| has the same material membership
as |\filldraw| with those options.

Every subpath must close, using |rectangle|, |circle|, |ellipse| or
|-- cycle|. Each individual contour must remain simple: self-intersections,
self-contact and backtracking are rejected after flattening. Distinct
contours may intersect or share edges; the constrained mesher subdivides
their intersections before classification. Curves use the existing
|curve tolerance| and coordinate-rounding conventions. Features lost in
that approximation cannot be recovered by refining the mesh.

All filled pieces of one compound region share its material, current
density, magnetisation direction and local mesh size. Total source current
is current density times the final painted area, not the sum of the
unsigned contour areas. Source current is not automatically renormalised
when another region paints over a winding. To use different properties
in several pieces, declare separate physical paths.

\subsection{Meshing, field sampling and cache reuse}
Both meshers include every subcontour as an interface constraint. The
grid mesher still requires axis-aligned straight edges. Delaunay local
refinement applies inside the filled portion and along all its contour
edges, including cavity walls. It does not request the same fine interior
lattice throughout a hole; another region can refine that hole if desired.
Boundary point spacing can nevertheless influence neighbouring triangles.

The full rectangular problem domain remains meshed. A hole in a material
is \emph{not an unmeshed exclusion}: air and any inner media still carry
field, and a measurement profile can cross them. No boundary condition
is added on the cavity wall. The ordinary material-interface transmission
conditions follow from the finite-element formulation. Sampling exactly
on an interface retains the documented first-containing-triangle
convention; use a profile offset when a particular side matters.

Unmeshed holes would change the PDE domain and require their own boundary
conditions and solvability checks. They are not introduced by |even odd rule|
or by a TikZ clip, and remain unsupported in this version. Measurement
profiles also retain a separate contract: each profile is one connected
curve, not a compound region boundary. Give separate contours separate
profile names.

The cache descriptor contains every contour and the fill rule. Changing
either invalidates the mesh, because local refinement also depends on
membership. Material-only changes can still reuse geometry. The cache
revision is now 4, including quality-refinement controls and diagnostics;
|auto| rebuilds earlier records and |frozen| rejects them.
The complete \texttt{examples/nested-domains.tex} example combines an alnico
magnet, an air cavity and a nonlinear iron shell with field lines and a
PGFPlots scan across all three media.

\subsection{Complete example: a magnet inside a shell}
The following source declares the alnico first and the compound iron
shell second. Its hole preserves both the inner magnet and the surrounding
air. The scan crosses the shell, air and magnet; the plotted jumps reflect
their different material laws and the elementwise field approximation.
The numbered comments also make the declaration/solve/plot order explicit.
\luafemmexample{nested-domains}

% SPDX-License-Identifier: LPPL-1.3c
% Copyright (C) 2026 Christophe Jorssen
% Author and Current Maintainer: Christophe Jorssen <christophe.jorssen@gmail.com>
% LPPL maintenance status: maintained
% This file is part of luafemm. See LICENSE and LICENSES.md for its terms.
\section{Boundary conditions}
\label{sec:boundaries}
By default, the outer rectangle is part of the physical model.
Section~\ref{sec:static} introduces custom computational domains and tagged paths. Leaving its four sides
at the default $A_z=0$ confines flux to a finite box. This is often useful
for a first calculation, but the boundary should be moved farther away
to assess its influence. Other conditions can impose a background field,
describe a symmetry plane, or specify a mixed relation between potential
and tangential field.

\subsection{A uniform field to start with}
Suppose that the empty domain should contain $B_x=0.2\,\mathrm T$ and
$B_y=0$. Since $B_x=\partial_y A_z$, choose $A_z=0.2y$, with $y$ in
metres. The rectangle below uses millimetres for geometry; the coefficient
is nevertheless entered as |.2|, because boundary data always use SI.

\begin{codeexample}[width=5cm]
\begin{tikzpicture}[femm/problem={
  xmin=-20,xmax=20,ymin=-15,ymax=15,
  mesh size=4},
  femm/boundary={type=dirichlet,
    y coefficient=.2}]
  \draw (-20,-15) rectangle (20,15);
  \pic[femm/lines=9] {femm field};
\end{tikzpicture}
\end{codeexample}

More generally, use |x coefficient=|$-B_y$ and
|y coefficient=|$B_x$ to prescribe an affine potential on the exterior.
Adding a material region changes the interior field. Imposing the
background potential on a finite rectangle still perturbs the exterior
solution: enlarge the domain and compare results.

\subsection{Command and key reference}
\begin{key}{/tikz/femm/boundary=\marg{options} (default \normalfont empty)}
Declare a boundary condition using ordinary TikZ keys. Its subkeys are
exactly those of |\femmboundary| below. In picture options, declarations
are collected and applied after the new model is created, before the
picture body starts. This key may precede or follow |femm/problem|.
Repeated declarations are applied in order; the last one for each side
wins. Setting |femm/problem| in several styles still creates only one
model, using the final problem options.

Normal TikZ |.style|, |.append style|, parameterised styles and
|every picture| are supported. Outside a picture, a direct
|\tikzset{femm/boundary=...}| records defaults for subsequent magnetic
pictures in that TeX scope. Each picture inherits its own copy; body
declarations do not change the defaults for later pictures. Pending
options, including macros and boundary-family styles, are evaluated
when the model starts. A picture without |femm/problem| does not apply
them to a previous model.

In the body of a started magnetic picture, the key applies immediately.
For example, |\tikzset{femm/boundary={side=left,type=neumann}}| is equivalent
to the corresponding command. It can also be used in |scope| options.
Physical declarations persist after the graphical scope ends: they are
not temporarily active boundary conditions. A declaration after meshing
raises an error, just as with the command.

This key selects complete sides of the model rectangle. On a drawing
path, it still modifies the selected side; it does not assign a boundary
condition to that path's contour. Prefer picture options or a dedicated
setup scope. Do not put it in |every path| or in a style used to draw
computed fields, since those paths may be drawn after the model is meshed.
Using the key without a value resets all sides to zero Dirichlet data.
\end{key}

\begin{codeexample}[code only]
\tikzset{
  half model/.style={
    femm/boundary={side=left,type=neumann}},
  background field/.style={
    femm/boundary={y coefficient=#1}},
  background field/.default=.2}
\begin{tikzpicture}[half model,
  femm/problem={xmin=0,xmax=85,ymin=-75,ymax=95,mesh size=4}]
  % Declare the half-device, then request a field or profile.
\end{tikzpicture}
\end{codeexample}

\begin{command}{\femmboundary\oarg{options}}
Declare a condition on one complete side of the rectangular domain, or
on all four sides. Call this command in the picture body, after the model
has been created and before the first mesh, solve or field request.
It produces no drawing. The side names refer to the original model frame;
a local |rotate| or |scale| does not change their meaning.

The keys belong to |/femm/boundary|. Each call first resets the six keys
to their defaults below, then processes its options. A later declaration
replaces the earlier condition on the selected side. Other sides keep
their conditions. Declarations persist in the current model even if the
command is inside a TeX group. They cannot modify a meshed model; start a
new picture for a different physical calculation.

All four numerical keys accept PGF mathematical expressions. Values are
expanded, evaluated with PGF's floating-point library, and passed to Lua
as finite numbers. No user text is evaluated as Lua code. Normal PGF
styles, including |.style| and |.append style|, are available in the family.
The command and |femm/boundary| share this implementation and can be
mixed in one model. Use the TikZ key in picture options; the command
itself requires an already created model.
\end{command}

\begin{key}{/femm/boundary/side=\meta{choice} (initially all)}
The |.is choice| alternatives are |all|, |left|, |right|, |bottom| and |top|.
They select $x=x_{\min}$, $x=x_{\max}$, $y=y_{\min}$ and $y=y_{\max}$,
respectively. A condition applies to the whole side, including material
regions which reach it. To configure several sides differently, repeat
the command. There is no automatic symmetry detection.
\end{key}
\begin{key}{/femm/boundary/type=\meta{choice} (initially dirichlet)}
A |.is choice| selector with alternatives |dirichlet|, |neumann| and
|robin|. Their physical meanings are given below. Every side initially
has homogeneous Dirichlet data, whether or not this command is used.
\end{key}
\begin{key}{/femm/boundary/value=\meta{expression} (initially 0)}
Constant term $g_0$ in $g(x,y)=g_0+g_xx+g_yy$, with coordinates in metres
relative to the model origin. Its unit is T\,m for Dirichlet data and
A/m for Neumann or Robin data.
\end{key}
\begin{keylist}[/femm/boundary]{x coefficient=\meta{expression} (initially 0),
  y coefficient=\meta{expression} (initially 0)}
The coefficients $g_x,g_y$ of the affine data. Their units are T for
Dirichlet and A/m$^2$ for Neumann or Robin. These are spatial coefficients,
not components of H; their effect also depends on the selected side.
The potential and field use the model's physical frame, independently of
graphic transformations.
\end{keylist}
\begin{key}{/femm/boundary/coefficient=\meta{nonnegative expression} (initially 0)}
The Robin coefficient $c$, in A/(T\,m$^2$). A nonzero value requires
|type=robin|. Negative coefficients are rejected because the present
solver requires a positive definite tangent after constraints are applied.
Zero makes a Robin condition identical to Neumann data.
\end{key}

\begin{codeexample}[code only]
\pgfkeys{/femm/boundary/symmetry/.style={type=neumann,value=0}}
% Inside a picture, before meshing:
\femmboundary[side=left,symmetry]
% All other sides retain their initial Az=0 condition.
\end{codeexample}

\subsection{Signs, corners and physical interpretation}
Let $\mathbf n=(n_x,n_y)$ be the outward unit normal and
$\mathbf t=(-n_y,n_x)$ the counterclockwise unit tangent. Write
$H_t=\mathbf H\cdot\mathbf t$, including the coercive contribution
$\mathbf H=\nu\mathbf B-H_c\mathbf m$ in a permanent magnet.
The three conditions are
\[
 \begin{array}{ll}
 \texttt{dirichlet}:&A_z=g,\\
 \texttt{neumann}:&H_t=g,\\
 \texttt{robin}:&H_t=cA_z+g.
 \end{array}
\]
On the left, right, bottom and top sides, $H_t$ is respectively
$-H_y$, $H_y$, $H_x$ and $-H_x$. The sign does not depend on how a
decorative rectangle is drawn. This orientation also gives Ampere's law
$\oint H_t\,\mathrm ds=\int J_z\,\mathrm dS$.

A constant Dirichlet value makes the normal component of B vanish on
that side: field lines run along it. Homogeneous Neumann data make the
tangential component of H vanish. In air and an isotropic unmagnetised
material, B is then normal to the side, which is useful for the appropriate
symmetry plane. In a permanent magnet, $H_t=0$ does not generally imply
$B_t=0$. The condition concerns the total H, not just a derivative of A.

At a corner shared by Dirichlet and weak data, the prescribed potential
wins at that node; the weak edge integral still contributes to its free
neighbours. Two Dirichlet sides must prescribe the same corner value.
Conflicts raise an error rather than silently choosing a side. Affine
values are compared with a relative tolerance of $10^{-10}$ and a
$10^{-12}$ T\,m comparison scale floor.

In an unmagnetised material, the Robin relation can also be written
\[
 \nu\,\partial_n A_z+cA_z+g=0.
\]
This corresponds to FEMM's static mixed condition with $c_0=c$ and
$c_1=g$. A positive Robin coefficient is not, by itself, an exact
open-boundary treatment. In particular, a circular asymptotic formula
cannot simply be transferred to this rectangular domain without an
approximation study.

\subsection{Solving only a symmetric half}
The complete example \texttt{examples/boundary-conditions.tex} models the
right half of a U electromagnet with mirrored windings on both legs.
The cut is $x=0$, where H has no tangential component. The only change to
the boundary setup is
\begin{codeexample}[code only]
\begin{tikzpicture}[femm/problem={
  xmin=0,xmax=85,ymin=-75,ymax=95,mesh size=4},
  femm/boundary={side=left,type=neumann}]
  % Declare the right half of the core, armature and winding.
  % Refine the right gap, then draw the field and register a profile.
\end{tikzpicture}
\end{codeexample}
The other sides remain at $A_z=0$. The example plots the gap field with
PGFPlots. The symmetry assumption concerns both geometry and excitation;
the original tutorial's single winding on one leg does not have this
exact mirror symmetry. A normal material interface needs no boundary
declaration: continuity is already part of the finite-element formulation.

The complete source and computed result follow. The style declares the
symmetry condition before |femm/problem| in the option list; it is applied
after model creation. The field lines cross the symmetry cut normally,
and the graph samples the remaining air gap.
\luafemmexample{boundary-conditions}

\subsection{When no potential is prescribed}
With only Neumann conditions (including Robin with $c=0$), A is determined
only up to a constant. The package first checks the discrete compatibility
condition $\int J_z\,\mathrm dS=\oint g\,\mathrm ds$, with a relative
tolerance of $10^{-10}$ against the sum of absolute load contributions.
It then fixes $A_z=0$ at the first mesh vertex, the lower-left corner for
the supported meshers. This is a gauge reference, not a magnetic boundary.
Incompatible data raise an error; a reference node must not conceal a
current imbalance. Any Dirichlet side or positive Robin coefficient
removes this constant nullspace and suppresses the extra reference.

Boundary types and numerical values are included in the solution cache
descriptor. Changing them reuses a matching mesh and solves again;
|mode=frozen| rejects an outdated solution. The cache revision was
increased for this feature; earlier caches are rebuilt in |auto| mode.

This chapter describes complete sides of the rectangular exterior, affine data
and constant nonnegative Robin coefficients. Native internal frontiers,
periodic/antiperiodic traces and circular multipole approximations are described
in Section~\ref{sec:static}; unmeshed holes belong to the domain chapter.
The independent circular moving-band workflow is described in
Section~\ref{sec:sliding}. Native analytical air-gap elements and exact
infinite elements remain outside the supported subset.

% SPDX-License-Identifier: LPPL-1.3c
% Copyright (C) 2026 Christophe Jorssen
% Author and Current Maintainer: Christophe Jorssen <christophe.jorssen@gmail.com>
% LPPL maintenance status: maintained
% This file is part of luafemm. See LICENSE and LICENSES.md for its terms.
\section{Static models and tagged frontiers}
\label{sec:static}
This chapter extends the planar, direct-current model: the computational domain
can follow a TikZ path; named boundaries can be paired; materials can have a
linear permeability tensor; permanent magnets can have a spatially varying
direction; and named windings can supply their current and report flux linkage.
All features use the same pure-Lua solver and generic TeX interface.
They do not introduce an axisymmetric or alternating-current formulation.

\subsection{Decide what each path means}
A magnetic drawing can contain four different kinds of path. A |femm/region|
path assigns material. A |femm/domain| path determines where the equation is
solved. A |femm/frontier| path assigns a physical boundary condition. A
|femm/profile| or |femm/selection| path measures an existing solution. Ordinary
TikZ paths only draw. Keeping these roles explicit makes the diagram easier to
edit: moving a profile cannot accidentally move a physical boundary.

Start with a coarse |mesh size| and modest |curve tolerance|. Finish the geometry,
sources and styles, then reduce both tolerances where the final result requires
it. Decreasing the curve tolerance improves the polygonal representation;
decreasing the mesh size improves the finite-element solution on that
representation. They are separate convergence checks.

The order of operations is:
\begin{enumerate}
\item Define material styles and reusable boundary/circuit styles.
\item Start |tikzpicture| with |femm/problem|. Circuit and pairing declarations in
  picture options are queued until the model exists, just like |femm/boundary|.
\item Capture the computational domain, material regions and tagged paths.
  Their declaration order is free, except that material paint order still matters.
  Declare each circuit before the mesh is generated; its paths may come first.
\item Mesh or solve. A field pic, profile sample or measurement may trigger this
  step automatically. Finish \emph{all} physical declarations first.
\item Draw the field, take measurements and plot their named snapshots.
\end{enumerate}
Physical declarations persist beyond the graphical scope that captured them.
They cannot be added after meshing. Reusable |.style|, |.append style| and
parameterised PGF styles work throughout this chapter. Numerical keys accept
PGF expressions; values passed to Lua are validated finite numbers.

\subsection{A computational domain drawn with TikZ}
\begin{codeexample}[width=6cm]
\begin{tikzpicture}[femm/problem={
  xmin=-22,xmax=22,ymin=-17,ymax=17,
  mesh size=4}]
  \path[femm/domain]
    (-18,-12)--(18,-8)--(14,13)--(-14,10)--cycle;
  \draw[thick,femm/frontier={name=outer,
    y coefficient=.2}]
    (-18,-12)--(18,-8)--(14,13)--(-14,10)--cycle;
  \pic[femm/lines=7] {femm field};
\end{tikzpicture}
\end{codeexample}

\begin{key}{/tikz/femm/domain=\marg{options} (default \normalfont empty)}
Capture one closed, possibly compound path as the computational domain. At most
one domain declaration is allowed per problem. The TikZ fill rule applies:
|even odd rule| conveniently describes holes. Holes in this domain exclude
triangles; holes in a material region merely expose the previous material.
The problem rectangle is still required as a bounding/staging box, and must
contain every physical path. Its sides cease to be physical boundaries when a
custom domain is present. Unlabelled custom-domain boundaries default to
$A_z=0$, independently of rectangle-side declarations.

With no |femm/domain|, the original rectangular domain and its side conditions
remain in effect. Path frontiers, custom domains and circuits require
|mesher=delaunay| and |field model=open|. They do not modify an imported FEM
model or an ideal-circuit model.
\end{key}
\begin{key}{/femm/domain/curve tolerance=\meta{distance} (initially .03)}
Positive flattening tolerance in model units. One closed domain can contain
several disjoint components. Ordinary TikZ curves, arcs and rounded corners are
flattened after their transformations have been evaluated.
\end{key}

\subsection{Tagging a physical boundary}
\begin{key}{/tikz/femm/frontier=\marg{options} (default \normalfont empty)}
Attach a named condition to the captured path. The path may be open or closed,
but must belong to the graph of physical boundaries/interfaces. An internal
Dirichlet contour or an internal path joining existing boundaries is supported
when it bounds valid faces. An isolated open internal constraint is rejected.
Neumann, Robin and periodic conditions require exterior edges.

The condition follows every split segment during arrangement construction and
meshing. It overrides the default condition on the same exterior edge.
Conflicting explicit tags on coincident edges are rejected. The existing key
|femm/boundary| retains its original meaning: it selects rectangle sides and
\emph{does not} tag the path on which it occurs.
\end{key}
\begin{key}{/femm/frontier/name=\meta{name} (initially empty)}
Required problem-local name. Names are unique within a problem; repeated names
are errors. Names need not be TeX control sequences.
\end{key}
\begin{key}{/femm/frontier/type=\meta{choice} (initially dirichlet)}
An |.is choice| key accepting |dirichlet|, |neumann|, |robin|, |periodic| and
|open|. With outward normal $\mathbf n$, define the positive tangent
$\mathbf t=(-n_y,n_x)$. Dirichlet prescribes $A_z=g$; Neumann prescribes
$\mathbf H\cdot\mathbf t=g$; Robin prescribes
$\mathbf H\cdot\mathbf t=cA_z+g$, with $c\geq0$.
|periodic| marks a path for a separate pairing declaration, whose sign can be
periodic or antiperiodic. |open| is the circular multipole approximation below.
\end{key}
\begin{key}{/femm/frontier/value=\meta{number} (initially 0)}
Constant term of $g=\texttt{value}+\texttt{x coefficient}\,x+
\texttt{y coefficient}\,y$, where $x,y$ are in metres. Potential data use T\,m;
Neumann/Robin data use A/m. Geometry remains in model units.
\end{key}
\begin{key}{/femm/frontier/x coefficient=\meta{number} (initially 0)}
Coefficient of the SI abscissa in $g$.
\end{key}
\begin{key}{/femm/frontier/y coefficient=\meta{number} (initially 0)}
Coefficient of the SI ordinate in $g$.
\end{key}
\begin{key}{/femm/frontier/coefficient=\meta{number} (initially 0)}
Nonnegative Robin coefficient $c$, in A/(T\,m$^2$); other types require zero.
\end{key}
\begin{key}{/femm/frontier/curve tolerance=\meta{distance} (initially .03)}
Positive path-flattening tolerance in model units. Use the same tolerance when
separate paths describe the same curved interface.
\end{key}
\begin{stylekey}{/femm/frontier/field tangent}
Homogeneous Dirichlet symmetry: $A_z=0$ and $B_n=0$.
\end{stylekey}
\begin{stylekey}{/femm/frontier/field normal}
Homogeneous natural symmetry: $H_t=0$. In isotropic, source-free air this means
$B_t=0$. In an anisotropic medium, $\mathbf H$ and $\mathbf B$ need not be parallel;
use the defining condition rather than inferring a direction for $\mathbf B$.
\end{stylekey}

\subsection{Working on a periodic sector}
An associated pair identifies the nodal potentials on two exterior traces.
Both traces are first synchronised by inserting the union of their transformed
vertices into their boundary triangles. The solver eliminates the constrained
unknowns exactly; it does not add a penalty parameter. A reversed sign is carried
through the residual, tangent, nodal updates and cache restoration.

The following sector contains an angularly magnetised permanent magnet. Its
$60^\circ$ pairing is antiperiodic. The angle law is chosen consistently with
that symmetry: the magnetisation rotates by $240^\circ$, namely the spatial
rotation plus a $180^\circ$ reversal.
\luafemmexample{static-frontiers}

\begin{key}{/tikz/femm/pair=\marg{options}}
Declare a pair through ordinary picture, scope or body keys. Picture declarations
are queued; body declarations apply immediately. The referenced paths must exist
by the time meshing starts, and each path belongs to exactly one pair.
\end{key}
\begin{command}{\femmpair\oarg{options}}
Equivalent immediate declaration in a started magnetic picture. Subkeys belong
to |/femm/pair|. Each declaration resets its options before applying the supplied
keys/styles. Multiple pairs may share corner vertices; contradictory nonzero
Dirichlet values at an antiperiodic corner are rejected.
\end{command}
\begin{key}{/femm/pair/first=\meta{name} (initially empty)}
Name of the source trace, declared with |type=periodic|.
\end{key}
\begin{key}{/femm/pair/second=\meta{name} (initially empty)}
Name of the target trace. It must be a rigid image of the first trace.
\end{key}
\begin{key}{/femm/pair/type=\meta{choice} (initially periodic)}
|periodic| imposes $A_2=A_1$; |antiperiodic| imposes $A_2=-A_1$.
The transform maps the first trace onto the second.
\end{key}
\begin{key}{/femm/pair/rotation=\meta{angle} (initially 0)}
Counterclockwise rotation in degrees about the specified centre, followed by
translation. Reflections, scaling and sliding interfaces are not included.
\end{key}
\begin{key}{/femm/pair/center x=\meta{coordinate} (initially 0)}
Abscissa of the rotation centre, in model units.
\end{key}
\begin{key}{/femm/pair/center y=\meta{coordinate} (initially 0)}
Ordinate of the rotation centre, in model units.
\end{key}
\begin{key}{/femm/pair/shift x=\meta{distance} (initially 0)}
Translation along the model $x$ axis after rotation, in model units.
\end{key}
\begin{key}{/femm/pair/shift y=\meta{distance} (initially 0)}
Translation along the model $y$ axis after rotation, in model units.
\end{key}
\begin{key}{/femm/pair/tolerance=\meta{distance} (initially .001)}
Positive matching tolerance in model units. It accommodates PGF coordinate
rounding. Transformed points within this tolerance are projected onto the actual
target trace. Choose a tolerance much smaller than the smallest relevant feature;
it is not a licence to pair physically different boundaries.
\end{key}

Symmetry of the material distribution and sources remains a modelling
responsibility. Synchronisation may reduce the minimum angle near a paired
boundary; inspect the final mesh statistics. No automatic multiplication of
energy, linkage or force measurements is performed. A sector energy can be
summed over identical sectors; force vectors must also be rotated before summing.

\begin{key}{/tikz/femm/repeat=\marg{options}}
Drawing options for the pic |femm repeated field|. The pic reuses solved contours
and the transform of a declared pair. Antiperiodic copies reverse their arrows.
It neither duplicates material regions nor creates new profiles or measurements.
\end{key}
\begin{key}{/femm/repeat/copies=\meta{integer} (initially 2)}
Number of copies including the solved sector, between 1 and 1000.
\end{key}
\begin{key}{/femm/repeat/pair=\meta{integer} (initially 1)}
One-based index of the pairing declaration whose transform is repeated.
\end{key}
For a $60^\circ$ sector, set |femm/repeat={copies=6}| on a
|femm repeated field| pic, as in the complete example above. Draw repeated
outlines with ordinary TikZ paths, without registering additional physical regions.

\subsection{Approximating an open exterior}
On a circle of radius $R$, the decaying exterior angular mode of order $n$ obeys
\[
 A_z\propto r^{-n},\qquad
 \mathbf H\cdot\mathbf t=\frac{n}{\mu_0 R}A_z.
\]
The |open| type installs this nonnegative Robin coefficient. It is exact for
that mode on a true circle in source-free air; on a flattened circle with
several angular modes it is an approximation. Enlarge the circle, refine its
polygonalisation and compare the observables of interest. The native model
checks that the adjacent elements are source-free air and that net longitudinal
current vanishes. Include the return conductor of a winding.

\begin{key}{/femm/frontier/radius=\meta{distance} (initially 1)}
Positive circular radius in model units for |type=open|. The path must be a
complete closed circle about the specified centre; its flattened vertices are
checked allowing curve/coordinate approximation. This key does not draw or
create the domain: declare the circular |femm/domain| separately.
\end{key}
\begin{key}{/femm/frontier/order=\meta{number} (initially 1)}
Positive exterior decay exponent $n$. The default is the first multipole.
Use another value only when the corresponding exterior angular symmetry is
known; this local boundary does not implement a full spectral exterior operator.
\end{key}
\begin{key}{/femm/frontier/center x=\meta{coordinate} (initially 0)}
Circular centre abscissa for |type=open|, in model units.
\end{key}
\begin{key}{/femm/frontier/center y=\meta{coordinate} (initially 0)}
Circular centre ordinate for |type=open|, in model units.
\end{key}

\subsection{Material directions and sheet mixtures}
For a linear material, the solver uses
\[
 \mathbf H=\frac{1}{\mu_0}\mathbf R(\alpha)
 \begin{pmatrix}1/\mu_x&0\\0&1/\mu_y\end{pmatrix}
 \mathbf R(\alpha)^T\mathbf B-H_c\mathbf m.
\]
The full symmetric tensor is used in assembly, field sampling and energy.
Its principal axes and the coercive direction are independent.

\begin{key}{/tikz/femm/materials/mu x=\meta{number} (initially empty)}
Positive relative permeability along the first material principal axis.
An omitted value inherits |mu r| or the selected library material.
\end{key}
\begin{key}{/tikz/femm/materials/mu y=\meta{number} (initially empty)}
Positive relative permeability along the second principal axis. These keys are
for linear media. Unequal principal permeabilities combined with a B--H curve
are rejected; nonlinear anisotropy needs a separate constitutive model.
\end{key}
\begin{key}{/tikz/femm/materials/material angle=\meta{angle} (initially 0)}
Counterclockwise orientation of the first principal axis in the physical model
frame, in degrees. A TikZ transformation of a region changes its geometry; it
does not silently rotate this numerical material direction. Set the angle
explicitly when a component's physical orientation changes. Rotating the entire
picture for presentation preserves the model-relative direction.
\end{key}
\begin{key}{/tikz/femm/materials/lamination=\meta{choice} (initially inherit)}
|inherit| uses a supported catalogue choice or defaults to |none|. The
|parallel x| and |parallel y| choices homogenise linear iron/air
sheets parallel to the selected principal axis, before applying |material angle|.
For fill $f$, permeability parallel to the sheets is $f\mu_r+1-f$; across them
it is $(f/\mu_r+1-f)^{-1}$. Nonlinear sheet mixtures and laminated permanent
magnets are rejected. The |stacked| choice mixes along both field-plane directions; its complex
harmonic extension and thickness parameter are detailed in Section~\ref{sec:losses}.
Catalogue nonlinear curves retain their existing static mixture interpretation.
\end{key}
\begin{key}{/tikz/femm/materials/lamination fill=\meta{fraction} (initially empty)}
Iron fraction $0<f\leq1$. Empty inherits the catalogue value, otherwise defaults
to one. A fraction below one requires a lamination choice.
\end{key}
\begin{key}{/tikz/femm/materials/conductivity=\meta{number} (initially empty)}
Nonnegative conductivity in S/m. An omitted value inherits the catalogue value,
converted from MS/m, or defaults to zero for a custom material. Conductivity
enters the DC winding resistance. It does not introduce eddy currents.
\end{key}

\subsection{Spatial magnetisation without executable expressions}
The region keys below describe the direction of $H_c\mathbf m$. Magnitude and
recoil/B--H law still come from the material. The existing
|magnetization angle| supplies the offset $\alpha_0$ in degrees. Centres and
angles are defined in the physical model frame, independently of drawing styles.

\begin{key}{/femm/region/magnetization=\meta{choice} (initially uniform)}
|uniform| uses the constant offset. |radial| uses
$\alpha=\alpha_0+\theta$; |tangential| uses
$\alpha=\alpha_0+90^\circ+\theta$; |angular| uses
$\alpha=\alpha_0+k\theta$, where
$\theta=\operatorname{atan2}(y-y_c,x-x_c)$ expressed in degrees.
For example |angular,magnetization order=2| describes a cylindrical Halbach
angle law. The family is passive data: arbitrary Lua expressions are not run.
\end{key}
\begin{key}{/femm/region/magnetization center x=\meta{coordinate} (initially 0)}
Abscissa $x_c$, in model units.
\end{key}
\begin{key}{/femm/region/magnetization center y=\meta{coordinate} (initially 0)}
Ordinate $y_c$, in model units.
\end{key}
\begin{key}{/femm/region/magnetization order=\meta{number} (initially 1)}
Multiplier $k$ for |angular|; ignored by the other presets. Positive, zero and
negative finite values are accepted. Noninteger orders introduce an angular
branch discontinuity; use them only on a domain where that is intentional.
\end{key}

Three symmetric quadrature points average the coercive direction in each
triangle. The weak source, reported element $\mathbf H$ and energy use the same
average. Its magnitude can be below one on coarse elements. Refine the magnet
until the angular variation is resolved. A polar law is undefined at its centre;
annular or sector-shaped magnets naturally avoid that point.

The following complete example combines an angularly magnetised annulus, a
circular open boundary and an ordinary field profile through its bore.
\luafemmexample{static-open}

\subsection{Named circuits and signed-turn windings}
A circuit supplies a terminal current or voltage. It contains one or more named
windings, each made from region legs. A leg's signed turn count determines the
sign of its longitudinal current density. Group outgoing and returning legs
under the same winding name. The default winding name is |main|.

\begin{key}{/tikz/femm/circuit=\marg{options}}
Declare a named circuit in picture options, scope options or the picture body.
Picture options are queued until the problem is created. Duplicate circuit
names are errors; every declared circuit must have at least one active leg.
\end{key}
\begin{command}{\femmcircuit\oarg{options}}
Equivalent immediate command in a magnetic picture. Subkeys are in
|/femm/circuit|, and defaults reset for each declaration.
\end{command}
\begin{key}{/femm/circuit/name=\meta{name} (initially empty)}
Required problem-local circuit identifier.
\end{key}
\begin{key}{/femm/circuit/type=\meta{choice} (initially series)}
|series| passes the terminal current through every winding. |parallel| splits
it among named windings according to their DC conductances. Legs of one winding
are always in series. Distinct parallel branches therefore need distinct
|winding| names, even if each branch contains two return-path legs.
\end{key}
\begin{key}{/femm/circuit/drive=\meta{choice} (initially current)}
|current| or |voltage|. This is a static resistive circuit model. There is no
induced voltage, skin effect or time-dependent circuit equation.
\end{key}
\begin{key}{/femm/circuit/current=\meta{number} (initially 0)}
Signed terminal current in amperes for current drive. Leave voltage at zero.
\end{key}
\begin{key}{/femm/circuit/voltage=\meta{number} (initially 0)}
Signed terminal voltage in volts for voltage drive. Leave current at zero.
\end{key}
\begin{key}{/femm/region/circuit=\meta{name} (initially empty)}
Associate a material region with a circuit. Its prescribed |current density|
must be zero; the circuit supplies its excitation. This also applies to
catalogue materials that carry a nonzero default source: explicitly set
|current density=0| for those materials.
\end{key}
\begin{key}{/femm/region/winding=\meta{name} (initially main)}
Group related outgoing/returning legs within the circuit.
\end{key}
\begin{key}{/femm/region/turns=\meta{number} (initially 1)}
Finite, nonzero signed turn count $N$. Negative turns reverse $J_z$.
A fractional effective count is allowed for homogenised winding models.
\end{key}
\begin{key}{/femm/region/winding fill=\meta{fraction} (initially 1)}
Conductor area fraction $0<f\leq1$. It changes resistance, not the total
ampere-turns or the material permeability.
\end{key}
\begin{key}{/femm/region/conductor length=\meta{distance} (initially 0)}
Electrical length per turn of this leg, in model units. Zero inherits problem
|depth|. An explicit larger value can account for allocated end-connection
resistance; those extra lengths do not create an out-of-plane magnetic model.
\end{key}

Let $S$ be a leg's \emph{active} area after all material paint operations, $\ell$
its electrical length, $\sigma$ its conductivity, and $I_w$ its winding current.
The implementation uses
\[
 J_z=\frac{NI_w}{S},\qquad
 R_{\rm leg}=\frac{N^2\ell}{\sigma fS}.
\]
Every active piece of the same leg shares this $J_z$. A completely overwritten
leg is an error. Resistances add within each winding. Winding resistances add
for a series circuit; conductances add for a parallel circuit. Positive
conductivity and length are required for every parallel or voltage-driven
winding. A current-driven series circuit can be solved without them, but its
resistance, voltage and power are then unavailable.

The complete example below sweeps current in a closed winding on an anisotropic
core and plots flux linkage using the existing PGFPlots snapshot interface.
\luafemmexample{static-circuits}

\subsection{Measuring a circuit and interpreting inductance}
\begin{key}{/femm/measure/quantity=circuit}
Resolve a named circuit after solving and store an independent result snapshot.
Use |\femmmeasure[name=coil-result,quantity=circuit,circuit=excitation]|.
The ordinary |parameter| key supplies a PGFPlots sweep abscissa.
\end{key}
\begin{key}{/femm/measure/circuit=\meta{name} (initially empty)}
Circuit identifier required by |quantity=circuit|. Region selections and profiles
do not select circuit legs: their circuit/winding associations do.
\end{key}

The scalar result components are |current| (A), |voltage| (V), |resistance|
($\Omega$), |power| (W), |linkage| (Wb) and |inductance| (H), subject to the
availability conditions below. Use |\femmresult|, |\femmvalue| and
|\femmresultplot| exactly as for other measurements. For magnetic depth $d$,
\[
 \lambda_w=d\sum_{\rm legs}\frac{N}{S}\int_S A_z\,dS.
\]
The triangle integral is exact for the P1 potential. Linkage is reported only
when the signed turns sum to zero in \emph{each} winding, so adding a constant
to $A_z$ cannot change it. An isolated conductor leg has no defined return path
and is not reported as a gauge-independent coil inductance.

For series windings the linkages add. For parallel windings the terminal linkage
uses the fixed DC conductance weights $R_{\rm eq}/R_w$, consistently with the
current split and magnetic energy. Problem depth is required for these measures.
At nonzero current, |inductance| is the secant ratio $\lambda/I$, and is omitted
when the model contains a permanent magnet. With multiple excited circuits it
includes their contribution to linkage; it is not an independently identified
self-inductance. Section~\ref{sec:operating} adds differential responses about a
saved bias and explains how to obtain self and mutual inductances by exciting
one port at a time. There is no automatic matrix-valued result. The resistance model includes only the declared conductor
lengths and conductivities; temperature dependence and AC losses remain separate.

\subsection{Interchange, caching and Lua entry points}
The FEM exporter writes diagonal permeability in file axes, the resolved DC
windings, and angular magnetisation laws. A rotated tensor with a nonzero cross
term cannot be represented by FEMM's two permeability entries and is rejected.
Linear sheet mixtures export their effective tensor. Parallel/voltage circuits
export series winding records at their resolved operating currents; they do not
retain a terminal voltage controller. Fill and electrical length are represented
by an effective DC conductivity. See Section~\ref{sec:interchange} for the full
contract and the distinction between source geometry and answer snapshots.

Periodic chains are split into pairs of straight primitives, each with one FEMM
periodic property. The importer accepts such straight pairs, diagonal linear
anisotropy, linear directional laminations and the restricted passive angular
expression grammar. Periodic circular-arc import, arbitrary executable angle
expressions and FEMM parallel-circuit import remain explicitly unsupported.

For spatial magnetisation, ANS output preserves the solved quadrature-averaged
source with an elementwise block representation. It is a field snapshot: use
the separate FEM export to edit the original angle law. ANS files also include
periodic nodal records. The native cache retains the original model semantics;
its revision is advanced so an older numerical cache cannot masquerade as a
solution of the extended model.

The main Lua entry points, independent of TeX, are:
\begin{itemize}
\item |luafemm-frontiers.domain(m, contours, fill_rule)|;
\item |luafemm-frontiers.path(m, name, points, closed, options)|, with
  |type|, |value|, |dx|, |dy|, |coefficient| and open-boundary parameters;
\item |luafemm-frontiers.pair(m, first, second, options)|, with
  |angle|, |x|, |y|, |cx|, |cy|, |sign| and |tolerance|;
\item |luafemm.material(m, name, properties)| additionally accepts
  |mur_x|, |mur_y|, |material_angle|, |lamination|, |lamination_fill| and
  |conductivity|;
\item |luafemm-static.region(region, options)| configures the returned region
  with |magnetization|, |center_x|, |center_y|, |order|, |circuit|, |winding|,
  |turns|, |fill| and |length|;
\item |luafemm-static.circuit(m, name, options)| and
  |luafemm-static.result(m, name)| declare and read a circuit.
\end{itemize}
Lua numerical arguments are numbers, not PGF expressions. Constructors copy
geometry and validate physical inputs. These modules are original Lua
implementations; neither a native mesher nor a native FEMM solver is needed
when compiling a document.

% SPDX-License-Identifier: LPPL-1.3c
% Copyright (C) 2026 Christophe Jorssen
% Author and Current Maintainer: Christophe Jorssen <christophe.jorssen@gmail.com>
% LPPL maintenance status: maintained
% This file is part of luafemm. See LICENSE and LICENSES.md for its terms.
\section{Axisymmetric magnets and circular windings}
\label{sec:axisymmetric}

\subsection{A rectangle that becomes a solenoid}
Suppose Kamil now wants to show his students the field of a finite solenoid.
A planar drawing of two opposite current rectangles would describe a device
invariant along the direction out of the page. His new coil is a body of
revolution instead. He draws only its meridional half-section: the horizontal
coordinate is the radius $r\geq0$, and the vertical coordinate is $z$.
Rotating one copper rectangle about the vertical axis produces the whole
winding volume. There is no second return rectangle at negative radius.

Kamil starts with a reusable style. It selects the geometry before setting its
bounds, so the meaning of every subsequent coordinate is clear.
\begin{codeexample}[code only]
\tikzset{solenoid/.style={
  femm/problem={type=axisymmetric,
    xmax=100,ymin=-85,ymax=85,mesh size=4},
  femm/circuit={name=coil,current=2}}}
\end{codeexample}
The coordinate unit is still one millimetre in the TikZ interface. The axis is
$r=0$, and |type=axisymmetric| sets |xmin=0|. The regularity condition there is
automatic. Other exterior sides default to zero azimuthal vector potential;
they are a finite truncation of the surrounding air, not an exact infinite
boundary. Kamil will enlarge the air domain separately from refining the mesh.

A single region now declares 300 circular turns carrying 2 A:
\begin{codeexample}[code only]
\filldraw[fill=orange!25,draw=orange!70!black,
  femm/region={material=copper,circuit=coil,turns=300,
    winding fill=.6,mesh size=1}]
  (20,-30) rectangle (30,30);
\end{codeexample}
Positive current flows along $+\varphi$ and gives a positive axial field in the
bore. The current density is $NI/S$, with $S$ the active meridional winding
area. The fill factor changes the electrical resistance; it does not reduce
these declared ampere-turns. Signed turns reverse the excitation.

Before solving, Kamil registers a profile on the axis. The existing drawing
and plotting commands keep their usual TikZ/PGFPlots syntax:
\begin{codeexample}[code only]
\path[femm/profile={name=solenoid-axis,samples=121}]
  (0,-60)--(0,60);
\pic[femm/lines=17] {femm field};
% In an ordinary PGFPlots axis:
\femmplot[profile=solenoid-axis,component=Bz,abscissa=z]
\end{codeexample}
The lines represent contours of the full flux $\Phi=2\pi rA_\varphi$.
Contours of $A_\varphi$ alone would not be magnetic field lines. Arrows follow
$(B_r,B_z)$, and equally spaced flux levels enclose equal full-revolution flux.
The profile samples the unsmoothed finite-element field, including its jumps
between elements. Increasing |samples| draws those jumps more densely; it does
not refine the field solution.

The complete example additionally refines the bore and compares its axial
profile with the Biot--Savart integral for the same uniformly filled coil in
infinite air. It also reports the winding resistance and secant inductance.
\luafemmexample{axisymmetric-solenoid}

Start with coarse global spacing and local refinement where the field is
observed. Once the drawing and excitation are settled, reduce both spacings,
then move the outer boundary. The analytic reference has no artificial outer
boundary, whereas the finite-element example does. A small discrepancy can
therefore contain both mesh and exterior-truncation error. Drawing more flux
lines or subdividing their display mesh changes neither error.

\subsection{Choosing the problem and rendering the flux}
\begin{key}{/femm/problem/type=\meta{choice} (initially planar)}
A PGF choice handler selecting |planar|, |axisymmetric| or |axial|.
The first two are vector-potential models; the transverse |axial| model is
described in Section~\ref{sec:currents}. All retain ordinary paths, compound
regions, material styles and the generic TeX loaders.
The axisymmetric choice sets |xmin=0|; place |type| before an explicit positive
|xmin| when modelling a domain that does not touch the axis. Returning to
|planar| does not reset bounds that have already been chosen.

For axisymmetry, all coordinates must have nonnegative radius. Every connected
component must touch the axis or have a Dirichlet or positive Robin anchor.
The solver rejects unanchored annular components rather than imposing
an arbitrary planar potential gauge. |depth| has no physical effect: the
problem represents a complete revolution, not a slice of prescribed thickness.
Imported FEM files select their own problem type; an explicit conflicting
|type| is rejected. Declaring only |type=axisymmetric| while importing does not
override bounds supplied by the file.
\end{key}

\begin{key}{/femm/problem/contour subdivisions=\meta{integer} (initially 2)}
For axisymmetric field rendering, subdivide every drawing triangle into four
at each level, from 0 to 4. The default therefore uses 16 drawing triangles
per solution element. The quadratic flux $2\pi rA_\varphi$ is evaluated at the
subdivision vertices and traced by the existing stitched marching-triangle
renderer. This is a controlled drawing approximation; it creates no solver
nodes, changes no measured field and does not invalidate a cached solution.
Large values increase drawing time and memory. Planar rendering ignores it.
As with other problem settings, choose it before the picture begins.
\end{key}

The field-line count, arrow spacing, arrow size, mesh and field styles remain
those of Section~\ref{sec:field}. Numerical contour levels supplied through the
Lua API are now full flux in Wb. Whole-picture rotations or scaling change the
presentation only. Geometry transformed inside the picture must still lie in
the physical half-plane; a negative-radius copy is not a valid return winding.
The existing U/E/toroidal and rotor/stator component nodes describe planar
models and are rejected by the axisymmetric formulation. Ordinary TikZ
nodes used as coordinates or annotations remain available.

\subsection{Fields, profiles and integrated quantities}
The following additions are choices of |/femm/plot/component|. They require an
axisymmetric model, except |Br| and |Hr|, which keep their existing polar meaning
for planar models.
\begin{center}
\begin{tabular}{lll}
\toprule Choice & Axisymmetric meaning & Unit\\
\midrule
|Br|, |Bz| & Radial and axial flux density & T\\
|Hr|, |Hz| & Radial and axial magnetic field & A/m\\
|Aphi| & Azimuthal vector potential & T\,m\\
|flux| & $2\pi rA_\varphi$ at the sample & Wb\\
|Bphi|, |Hphi| & Azimuthal components, identically zero & T, A/m\\
\bottomrule
\end{tabular}
\end{center}
The choices |/femm/plot/abscissa=r| and |z| return the model coordinates, in mm
for TikZ. There is no change to the default magnitude |component=B| or arc-length
|abscissa=s|. |Bx| and |By| remain aliases for the horizontal and vertical
components; |Az| is a compatibility column containing the solved potential
$A_\varphi$. Prefer the explicit cylindrical names in new axisymmetric code.
The meridional polar names |Btheta| and |Htheta| are rejected for these models:
they would be easy to confuse with the physically absent azimuthal field.
Local tangent/left-normal components and |circulation| retain their geometric
meaning in the drawn meridional plane. Circulation is integrated by Gaussian
quadrature after splitting each segment at triangle crossings.

Regional |quantity=integrals| keeps the meridional area and current
$\int J_\varphi\,\mathrm dr\,\mathrm dz$, but computes volume, energy and coenergy
over the full revolution. Field means and reported centroids are weighted by
$2\pi r$; they are cylindrical-component means, not Cartesian vector means
around a complete circle. Result aliases |Br|, |Bz|, |Hr| and |Hz| are available
through |\femmresult|, |\femmvalue| and |\femmresultplot|.

For |quantity=force|, |Fz| is the full axial resultant and aliases |Fy| in the
picture frame. The global transverse resultant |Fx| and torque vanish by
axisymmetry. They do not describe the local radial stress or hoop loading.
Weighted Maxwell stress remains the default, with its existing weighting styles.
A body may meet the symmetry axis but must otherwise be separated from other
materials and outer boundaries by source-free air. A closed Maxwell contour
may include a segment on the axis; its surface-of-revolution area is zero.
Lorentz force applies to nonmagnetic conductors and integrates
$-J_\varphi B_r$ over the physical volume. Force arrows use the existing keys
and point along the drawn axial direction.

The Lua section-flux call |femm.flux(model,points)| returns flux through the
surface obtained by revolving the path, with its right meridional normal:
\[
 \Phi_{\rm section}=2\pi\bigl(r_{\rm start}A_{\varphi,\rm start}
                    -r_{\rm end}A_{\varphi,\rm end}\bigr).
\]
It verifies that the path stays in the computational domain. Endpoint snapping
is not supported for axisymmetric sections. The corresponding TeX call is:
\begin{codeexample}[code only]
\femmflux[profile=name,snap tolerance=0]
\end{codeexample}
The default |snap tolerance=auto| already selects zero for axisymmetric
profiles, including imported models; the explicit value above documents it. For a uniform positive $B_z$, a
radial path from its outer radius towards the axis has positive flux.

\subsection{Circular windings and material directions}
Every declared winding cross-section represents complete circular turns, so
linkage requires no cancellation of signed outgoing and returning legs.
For uniform turn density on an active region $S$,
\[
 \lambda={N\over S}\int_S 2\pi r A_\varphi\,\mathrm dr\,\mathrm dz,
 \qquad R={N^2\ell\over\sigma fS},\qquad
 \ell=2\pi\bar r.
\]
The default conductor length is the mean circumference, using the final active
region after painting and subtraction of holes. An explicit |conductor length|
replaces that per-turn length for resistance only, with the usual model units.
The existing series/parallel and current/voltage choices apply unchanged.
Secant inductance remains $\lambda/I$ under the documented restrictions; it
is neither incremental inductance nor a general multiport self-inductance.

Material principal directions and |material angle| are interpreted in the
$(r,z)$ frame. |magnetization angle=90| points along $+z$; zero points along
$+r$. The existing spatial |radial|/|tangential| presets vary around a centre
\emph{in this meridional plane}. They do not mean azimuthal magnetisation about
the revolution axis. Use angle zero for an ordinary cylindrically radial source.
The solver has no azimuthal magnetic-field unknown.

The second example revolves a magnet rectangle touching the axis into a solid
cylinder. Its linear recoil law and axial magnetisation produce the following
field and axis profile. Single-valued nonlinear B--H laws are also supported.
\luafemmexample{axisymmetric-magnet}

\subsection{Formulation, axis regularity and boundary conventions}
The unknown is continuous piecewise-affine $A_\varphi$ on the physical
$(r,z)$ triangles, with
\[
 B_r=-\partial_z A_\varphi,\qquad
 B_z=\partial_r A_\varphi+{A_\varphi\over r},\qquad
 \mathrm dV=2\pi r\,\mathrm dr\,\mathrm dz.
\]
Axis nodes are fixed to zero automatically. Nonzero prescribed axis values
are errors. At an axis sample the regular one-sided limit is used:
$B_r=0$ and $B_z=2\partial_r A_\varphi$. Field values inside an element vary
with radius; the element's stored field is its centroid diagnostic, not a
replacement for point evaluation or volume quadrature.

Writing $\mathcal C v=(-\partial_zv,\partial_rv+v/r)$, the weak residual per
radian is
\[
 \int_\Omega r\bigl[\mathbf H\cdot\mathcal Cv-J_\varphi v\bigr]
       \,\mathrm dr\,\mathrm dz
 +\int_{\partial\Omega_w}r(cA_\varphi+g)v\,\mathrm ds=0.
\]
$\mathbf H$ includes the coercive source. A positive seven-point triangle rule
integrates the radius-weighted material residual and its consistent symmetric
Newton tangent. Regional energy uses the same quadrature. Linear tensor and
isotropic nonlinear material laws are shared with the planar solver. Spatial
coercivity directions retain the documented elementwise averaging rule.

Dirichlet data specify $A_\varphi$ in T\,m. For the counterclockwise meridional
tangent $\mathbf t=(-n_z,n_r)$, the weak Neumann/Robin convention is
\emph{$-\mathbf H\cdot\mathbf t=cA_\varphi+g$}. The sign differs from the planar
curl convention. All affine coefficients still use SI coordinates. The exact
radially weighted edge mass matrix is used. A uniform positive $B_z$ in air,
for example, requires |value=-Bz/mu0| on a natural outer vertical boundary.
An all-natural domain that touches the axis is regularity-anchored; its total
azimuthal current need not vanish.

\subsection{Exchange and validation limits}
The supported FEM input/output subset now accepts axisymmetric headers, geometry,
materials, circular winding currents and supported fixed/mixed boundaries.
ANS output writes nodal \emph{flux} $2\pi rA_\varphi$, following FEMM's convention.
FEMM uses different flux interpolation inside triangles, so its reconstructed
$B$ and $H$ from these nodal values are not an exact snapshot of our P1
$A_\varphi$ derivatives. Preserve the native cache for exact reuse, and compare
FEMM field samples under refinement. The uniform-field patch agrees to rounding;
nonuniform snapshots are subject to this interpolation difference.
The primary comparison is described in
\href{https://femm.info/doku/doku.php?id=axisymmetricformulation}{FEMM's
axisymmetric formulation note} and in the audited xfemm source.

Analytic tests cover a uniform axial field, full-revolution volume and energy,
flux orientation, natural and mixed boundary signs, a nonlinear magnet with
$H=0$, circular winding resistance and the identity $W=I\lambda/2$.
For the manufactured $A_\varphi=r^3$, $B_z=4r^2$ and
$J_\varphi=-8r/\mu_0$, radial-volume-weighted field errors at spacings
.2, .1 and .05 m are 9.287\%, 4.602\% and 2.304\%.
Analytical conductor forces check all three force methods and their signs.
The generic format matrix includes cylindrical profiles and cold/warm/shared
frozen winding caches under all three formats.

The optional |scripts/axisymmetric-interop.py| compares uniform fixed/mixed fields, solenoids with air or nonlinear iron cores,
and cylindrical magnets through xfemm's native reader and independent solver. It reports
its samplewise values and tolerances in |build/axisymmetric/xfemm-results.json|.
All 36 local comparisons passed; the largest peak-normalised field discrepancies
were .532\% for snapshots and 1.974\% for independent solves.
These checks establish a bounded static subset, not complete axisymmetric FEMM
parity. Periodic pairs, the planar ideal model, the circular open-boundary
approximation and FEMM exterior transformations are explicitly unavailable.
The axial H model is described separately in Section~\ref{sec:currents}.
Azimuthal magnetic fields, sliding gaps and previous
solution modes remain separate work. Use domain enlargement and mesh refinement
for exterior studies in this formulation.

% SPDX-License-Identifier: LPPL-1.3c
% Copyright (C) 2026 Christophe Jorssen
% Author and Current Maintainer: Christophe Jorssen <christophe.jorssen@gmail.com>
% LPPL maintenance status: maintained
% This file is part of luafemm. See LICENSE and LICENSES.md for its terms.
\section{Alternating fields, eddy currents and skin depth}
\label{sec:harmonic}

\subsection{Kamil turns up the frequency}
Kamil's static solenoid is ready. He now wants to explain why an alternating
field penetrates a copper body less easily as its frequency increases.
Drawing fewer field lines inside the copper would illustrate an assumption;
instead, he wants the induced currents to change the computed field.
He starts with a slab, for which he can compare the numerical result with an
analytic diffusion solution, before returning to a circular coil.

The new ingredient is |frequency|, in hertz. Zero retains the static solver;
a positive value selects the linear single-frequency diffusion problem.
The geometry and material paths remain ordinary TikZ paths. For his slab,
Kamil puts real, opposite potentials on the two vertical surfaces and natural
conditions on the horizontal surfaces. These boundary conditions describe a slab of finite horizontal thickness,
invariant along the vertical and out-of-page directions.
\begin{codeexample}[code only]
\tikzset{conducting slab/.style={
  femm/problem={frequency=400,xmin=-5,xmax=5,
    ymin=0,ymax=2,mesh size=.25,depth=10},
  femm/boundary={side=all,type=neumann},
  femm/boundary={side=left,type=dirichlet,value=-.0001},
  femm/boundary={side=right,type=dirichlet,value=.0001}},
  femm/materials/slab copper/.style={mu r=1,conductivity=58000000}}
\end{codeexample}
Here the conductivity is $58\,\mathrm{MS/m}$; the key itself expects
$\mathrm{S/m}$. The prescribed vector potentials are in $\mathrm{T\,m}$,
regardless of the millimetre drawing unit. They are phase-zero \emph{peak}
phasors. Kamil registers a profile across the slab, then plots the amplitude
of the induced current density:
\begin{codeexample}[code only]
\filldraw[fill=orange!20,
  femm/region={material=slab copper}] (-5,0) rectangle (5,2);
\path[femm/profile={name=skin,samples=101}] (-5,1)--(5,1);
\pic[femm/phase=0,femm/lines=13] {femm field};
% In a PGFPlots axis:
\femmplot[profile=skin,component=Jeddy,abscissa=x,
  representation=amplitude]
\end{codeexample}

For an isotropic conductor, the skin depth is
\[
 \delta=\sqrt{\frac{2}{\omega\mu\sigma}},\qquad \omega=2\pi f.
\]
It is about $3.3043\,\mathrm{mm}$ in this example. Kamil starts with a few
triangles per skin depth while adjusting the drawing, then refines the
conductor before trusting the loss estimate. More profile samples do not
resolve the diffusion layer: only a finer \emph{field mesh} does that.
At higher frequency he must revisit the local spacing, even if the geometry
has not changed. The automatic mesh has no skin-depth error estimator.

For half-width $a$, surface potential $A_s$ and $k=(1+\mathrm{i})/\delta$,
this slab has $\widehat A(x)=A_s\sinh(kx)/\sinh(ka)$ and
$\widehat J_{\rm eddy}=-\mathrm{i}\omega\sigma\widehat A$.
The complete example plots the finite-element amplitude against this formula
and reports the cycle-average Joule loss for the declared depth.
\luafemmexample{harmonic-skin}

\subsection{A conducting tube inside the coil}
Kamil now returns to his axisymmetric solenoid. A positive frequency works with
|type=axisymmetric|: horizontal coordinates are radii, and all physical volume
integrals include the complete revolution. He adds a rectangular copper region
between two positive radii; rotating it produces a conducting tube.

The coil and tube must have different electrical interpretations. The coil
contains insulated fine wires carrying a prescribed terminal current. A named
current-driven series circuit treats its regions as \emph{stranded}: no
resolved eddy current flows through the homogenised winding cross-section.
The tube is a \emph{solid} conductor. Its induced azimuthal currents close around
the axis and react on the coil's field.
\begin{codeexample}[code only]
\tikzset{induction/.style={
  femm/problem={type=axisymmetric,frequency=1000,
    xmax=40,ymin=-40,ymax=40,mesh size=2},
  femm/circuit={name=drive,current=2,phase=0}}}
% In this picture:
\fill[femm/region={material=copper,circuit=drive,
  turns=100,winding fill=.6,mesh size=1}]
  (14,-15) rectangle (18,15);
\fill[femm/region={material=copper,conductor=solid,mesh size=.5}]
  (7,-12) rectangle (10,12);
\end{codeexample}
The default outer boundary still truncates the surrounding air. Refining the
skin layer and enlarging the air domain are separate checks. The complete
example selects the tube for its loss integral and also measures the coil's
complex terminal response.
\luafemmexample{harmonic-shield}

\subsection{Peak phasors and the meaning of a picture}
Every harmonic source uses the convention
\[
 u(t)=\mathop{\rm Re}\{\widehat u\,e^{\mathrm{i}\omega t}\},\qquad
 \widehat u=u_{\rm re}+\mathrm{i}u_{\rm im}.
\]
A source of amplitude $I$ and phase $\alpha$ therefore has
$\widehat I=I(\cos\alpha+\mathrm{i}\sin\alpha)$. Its instantaneous value at
$\omega t=\theta$ is $I\cos(\theta+\alpha)$.
A negative signed amplitude is allowed; it reverses the source.
Inputs and amplitude plots use peak values, as does FEMM; divide a scalar
phasor amplitude by $\sqrt2$ to obtain its RMS value.
See the \href{https://www.femm.info/doku/doku.php?id=FAQ}{official FEMM FAQ}
for the upstream peak/RMS convention.

A field-line picture is an \emph{instantaneous} snapshot. At phase $\theta$,
it traces contours of $A_{\rm re}\cos\theta-A_{\rm im}\sin\theta$ in planar
mode, or of $2\pi r(A_{\rm re}\cos\theta-A_{\rm im}\sin\theta)$ in axisymmetry.
Arrows follow the corresponding instantaneous vector field. Contours of
$|\widehat A|$ are not substituted for physical field lines. Several pics can
show different phases from the same solution without solving again.
\begin{codeexample}[code only]
\pic[femm/phase=0,femm/lines=17] {femm field};
% An optional second layer, using the same solved phasors:
\pic[femm/phase=90,femm/field style={red,dashed}] {femm field};
\end{codeexample}

\subsection{Keys and order of operations}
Declare the problem, material laws, circuits, boundary conditions and physical
regions first. Register profiles and selections before or after the solve as
usual. A field pic, profile plot or measurement solves on demand. Once meshed,
physical declarations are immutable; create another problem for a frequency
sweep. Drawing phases and plot representations are presentation choices and
do not invalidate the physical cache.

\begin{key}{/femm/problem/frequency=\meta{hertz} (initially 0)}
A finite nonnegative number. Zero selects static magnetics. A positive number
selects the linear complex solver in either planar or axisymmetric geometry.
This is a temporal frequency, independent of spatial winding Fourier orders.
Frequency enters the computation fingerprint; changing it requires a new
solution. A compatible mesh may still be reused. The static |max iterations|
key controls Newton iterations only. Each harmonic field-block GMRES solve has a fixed maximum of
4000 iterations and uses the existing |tolerance| as its relative true-residual
criterion; a failed solve raises an error instead of saving a partial field.
\end{key}

The conductor choices, isolated-current constraints and coupled circuit
sources are documented in Section~\ref{sec:ac-circuits}. By default,
|conductor=auto| selects stranded interpretation for named circuit regions
and solid interpretation otherwise. An unconstrained solid retains the
zero-imposed-voltage model; insulation is not inferred from its geometry.

\begin{key}{/femm/region/current phase=\meta{degrees} (initially 0)}
Phase of the region's signed |current density| amplitude. Only harmonic
problems accept a nonzero phase. For a named winding, use the circuit phase
instead; a nonzero region phase on a circuit leg is rejected. The default
material current density remains a real amplitude unless this key rotates it.
\end{key}

\begin{key}{/femm/circuit/phase=\meta{degrees} (initially 0)}
Phase of the prescribed peak current or voltage, in degrees. Signed turns
still determine each leg's orientation. Series and parallel networks use the
coupled formulation in Section~\ref{sec:ac-circuits}.
\end{key}

\begin{key}{/tikz/femm/phase=\meta{degrees} (initially 0)}
Temporal phase $\omega t$ of a |femm field| pic. It affects both contour geometry
and arrow orientation. It does not modify the stored phasors. Static pictures
ignore temporal phase because their solution is time independent.
\end{key}

\begin{key}{/femm/plot/representation=\meta{choice} (initially snapshot)}
A choice handler accepting |snapshot|, |real|, |imaginary|, |amplitude|, |rms|
and |phase|. It acts on the chosen complex component after projection onto the
profile tangent, normal or polar basis. |phase| returns the principal argument
in degrees in $[-180,180]$; it is not unwrapped, and a zero phasor yields a
plotting gap. |snapshot| returns the instantaneous value at the plot's |phase|.

For the vector norms |B| and |H|, |amplitude| is
$\sqrt{|\widehat u_x|^2+|\widehat u_y|^2}$; |rms| divides this by $\sqrt2$.
The former is the complex vector norm, not in general the maximum instantaneous
length of an elliptically rotating field. |snapshot|, |real| and |imaginary|
return the Euclidean norm of the corresponding real vector. A vector norm has
no unique scalar phase: requesting |representation=phase| with |B| or |H| is
rejected. Choose a signed component instead. Static profiles accept |snapshot|
and |real| only. Harmonic circulation plots are deferred.
\end{key}

\begin{key}{/femm/plot/phase=\meta{degrees} (initially 0)}
Temporal phase for |representation=snapshot|. Other representations ignore
this key. It is independent of the phase selected for field-line pics.
\end{key}

\begin{key}{/femm/plot/component=\meta{component} (initially B)}
Harmonic plots retain the documented magnetic components and add |J| and
|Jeddy|, in $\mathrm{A/m^2}$. These are signed out-of-plane scalar phasors in
planar mode and signed azimuthal scalar phasors in axisymmetry. |J| includes
both impressed and induced currents; |Jeddy| contains only the latter.
For stranded regions |Jeddy=0|. The axisymmetric |flux| component is the full
complex flux $2\pi r\widehat A_\varphi$, in webers.
\end{key}

\subsection{Measurements and power balance}
The existing |\femmselect| and |\femmmeasure| commands retain their syntax.
For |quantity=integrals| on a harmonic model, the available components are
|area|, |volume|, |current_real|, |current_imaginary|, |energy|, |coenergy|,
|joule_loss| and |eddy_loss|. Energy and coenergy are cycle averages and coincide
for the supported real linear permeability. Currents integrate the total
complex current density over the section, without a $2\pi r$ volume factor.
\begin{codeexample}[code only]
\femmselect[name=tube,material=solid copper]
\femmmeasure[name=tube-loss,selection=tube]
\femmvalue[component=joule_loss]{tube-loss}
\end{codeexample}
The reported quantities are
\[
 \overline W=\frac14\int\mathop{\rm Re}
  (\widehat{\mathbf B}\mathbin{\cdot}\widehat{\mathbf H}^{*})\,dV,
 \qquad
 P_{\rm Joule}=\int\frac{|\widehat J|^2}{2\sigma}\,dV.
\]
Nonconducting regions contribute no Joule loss. Stranded regions use the
fill-adjusted conductivity for their prescribed homogenised current.
|eddy_loss| uses only $\widehat J_{\rm eddy}$; it is \emph{not an additive
partition} of |joule_loss| when an impressed current is also present, because
the latter includes an interference term. Use |joule_loss| for heat generation.

For |quantity=circuit|, |current_real|, |current_imaginary| and the available DC
|resistance| are reported. Closed windings additionally provide |linkage_real|
and |linkage_imaginary|. When resistance is known, they provide |voltage_real|,
|voltage_imaginary|, |power|, |reactive_power|, and, at nonzero drive current,
|impedance_real| and |impedance_imaginary|. Units are A, Wb, V, W, var and ohms.
The terminal convention is
\[
 \widehat V=R\widehat I+\mathrm{i}\omega\widehat\lambda,\qquad
 S=\tfrac12\widehat V\widehat I^{*}.
\]
Power is $\mathop{\rm Re}S$ and reactive power is $\mathop{\rm Im}S$.
Mutual coupling is included in linkage: one driven circuit's real power can
include exchange with another source, so it is not automatically a local loss.
No scalar inductance is inferred from this multi-source complex response.
Planar winding closure still requires balanced signed turns per winding;
axisymmetric turns are closed circles. An explicit conductor-length override
changes terminal resistance, while regional loss integrals use the actual
model volume; their sums then need not agree without including the additional
wire volume in the model.

Cycle-average and twice-frequency forces/torques are not implemented in this
formulation; harmonic force measurements raise an error. Static regional field
averages and scalar |current| are also unavailable for harmonic measurements,
avoiding an implicit choice between a real part and an amplitude.
The low-level real field sampling and section-flux routines retain the
phase-zero snapshot convention. Use phasor profiles for other representations.
For |\femmstat|, harmonic |peak| is the largest centroid complex B-vector norm;
|net_current| and |net_current_imaginary| are the real and imaginary total
section currents. |gmres| is the iteration count; static |newton| and |cg|
statistics are zero.

\subsection{Formulation, validation and current limits}
The unknown is a continuous complex P1 $A_z$ or $A_\varphi$. With $C(v)$ the
appropriate planar or cylindrical curl and real reluctivity tensor $\nu$,
assembly adds the consistent conductive mass to the static stiffness:
\[
 \int \nu C(\widehat A)\mathbin{\cdot}C(v)\,d\Omega
 +\mathrm{i}\omega\int\sigma\widehat A v\,d\Omega
 =\int\widehat J_{\rm src}v\,d\Omega+\hbox{boundary terms}.
\]
Here $d\Omega=dx\,dy$ in planar mode and $r\,dr\,dz$ in axisymmetry; the common
$2\pi$ factor is restored in physical volume observables. Axis regularity
still imposes $A_\varphi=0$. The same quadrature is used for assembly and losses.
The real diffusion stiffness and imaginary mass are solved by restarted GMRES
with complex ILU(0), reorthogonalisation and an explicit true-residual check.
The static conjugate-gradient solver is not used on this non-Hermitian system.

The solver supports real positive linear permeability tensors, finite
nonnegative scalar conductivity, real affine boundary phasors, fine-wire
prescribed series currents and conductive bodies. Automatic static potential
gauges are rejected: supply an anchored boundary. Nonlinear harmonic material
laws, detailed wire models, periodic pairs, circular open boundaries,
sliding gaps and motion-induced currents
remain separate extensions. Permanent-magnet coercivity is rejected here;
compute its DC field separately instead of oscillating a permanent source.
Section~\ref{sec:biased-harmonic} describes the supported small-signal extension
around a captured nonlinear DC bias, including permanent-magnet bias.
Section~\ref{sec:losses} adds linear complex permeability and supported insulated
sheet losses. Unsupported catalogue loss, lamination and wire models are rejected.

Public |.fem|/|.ans| interchange remains static-only. Harmonic exports fail
explicitly rather than writing a misleading frequency-zero file. The passive
native cache stores both potential arrays and fingerprints frequency, phase,
conductivity and conductor interpretation; format 13 requires recomputing older
snapshots. It is reusable across all three supported TeX formats.

\nolinkurl{tests/test_harmonic.lua} checks analytic skin-depth convergence,
phase signs, the low-frequency limit, axis regularity, copper losses, complex
terminal power balance, rejected combinations and cache round trips.
\nolinkurl{scripts/harmonic-interop.py} independently meshes and solves passive
AC test fixtures with xfemm, comparing complex A/B/H/J, energy and losses.
That test-only file adapter does not widen the public interchange contract.
The generic format matrix exercises profiles, snapshots and shared frozen
phasors in plain LuaTeX, LuaLaTeX and ConTeXt MkIV.

% SPDX-License-Identifier: LPPL-1.3c
% Copyright (C) 2026 Christophe Jorssen
% Author and Current Maintainer: Christophe Jorssen <christophe.jorssen@gmail.com>
% LPPL maintenance status: maintained
% This file is part of luafemm. See LICENSE and LICENSES.md for its terms.
\section{Harmonic circuits and isolated conductors}
\label{sec:ac-circuits}

A current source fixes a current, not its spatial distribution in a massive
conductor. A voltage source fixes a voltage, not the currents in neighbouring
windings. At nonzero frequency, both questions must be solved with the magnetic
field: skin effect, proximity effect and mutual induction change the answer.
This chapter adds those constraints to the harmonic conventions of
Section~\ref{sec:harmonic}. All amplitudes are \emph{peak} phasors with
$\exp(\mathrm i\omega t)$; displayed powers are cycle averages.

\subsection{Two parallel coils, one isolated bar}
Start with a coarse mesh to choose the layout and check signs. Declare a
circuit at picture level, then give each winding a name. In a planar model,
a complete winding has outgoing and returning legs with opposite signed
turns. Legs with the same |winding| belong to one branch; |type=parallel|
connects the branches to the same voltage. A material choice alone does not
specify an electrical connection.

Here the lower branch has 20 turns and the upper one has 40. The source imposes
$0.1\,\mathrm V$ at $2\,\mathrm{kHz}$ with a $30^\circ$ phase. The central
copper bar is electrically isolated: its total longitudinal current must vanish,
although it supports local eddy currents. The field and branch currents are
computed together. Their division is not a DC resistance ratio.

\luafemmexample{harmonic-circuits}

The example uses a reusable |/femm/region/coil| style inside each region's key
list. Declare one |femm/region| capture per path. Finish all geometry and source
declarations before the first field pic, measurement or explicit solve. Then
register profiles and plot their stored samples. Branch measurements use the
same |quantity=circuit| interface as terminal measurements.

Before trusting loss values, refine the conductor mesh relative to the skin
depth $\delta=\sqrt{2/(\omega\mu\sigma)}$ and move the artificial outer boundary.
A small algebraic residual is not evidence of spatial convergence. Constrained
solid conductors add field-response solves; begin with a few branches and a
moderate mesh, then refine the final figure.

\subsection{Sources, branches and region choices}
The existing |/femm/circuit| family is shared by DC and AC. Its declaration is
local to the problem; repeated names are rejected. Keys reset to their defaults
for each circuit declaration. Circuit legs inherit ordinary PGF region styles
within TeX groups. Region keys remain scoped to their path or surrounding scope.

\Needspace{7\baselineskip}
\begin{key}{/femm/circuit/type=\meta{series or parallel} (initially series)}
In harmonic mode, |series| imposes equal complex branch currents and adds their
terminal voltages. |parallel| imposes equal complex voltages and adds currents.
A branch is the set of regions with the same |winding| within this circuit.
Mutual induction is included through the common field solution.
\end{key}
\Needspace{7\baselineskip}
\begin{key}{/femm/circuit/drive=\meta{current or voltage} (initially current)}
Choose the prescribed whole-port quantity. |current| is in amperes and
|voltage| in volts; both default to zero. Set only the driven quantity, leaving
the other at zero. |phase|, initially zero, is its phase in degrees and accepts
PGF mathematical expressions. Thus |drive=voltage,voltage=.1,phase=15+15|
means $\widehat V=0.1\exp(\mathrm i30^\circ)\,\mathrm V$. A zero-voltage
solid ring models a shorted turn. Zero current does not mean zero local current
density in a solid conductor.
\end{key}
\Needspace{7\baselineskip}
\begin{key}{/femm/region/current constraint=\meta{none or zero} (initially none)}
A choice handler for an individual connected solid region. |zero| introduces
an unknown voltage such that $\int_S\widehat J\,dS=0$. It requires positive
frequency, positive material conductivity, and a conforming Delaunay mesh.
It cannot accompany a circuit, an impressed current density or phase, or
|conductor=stranded|. Use separate declarations for disconnected conductors;
a region split into disconnected active pieces by later paint operations is
rejected. Holes inside one connected conductor are allowed.

With |none| and no circuit, the existing solid model retains zero imposed
voltage. It does not infer insulation or a floating electrical connection from
a closed TikZ path. The key is inherited like other region keys and resets
at the end of its TeX group.
\end{key}
\Needspace{7\baselineskip}
\begin{key}{/femm/region/conductor=\meta{auto, solid or stranded} (initially auto)}
|auto| selects stranded interpretation for a named circuit region and solid
interpretation otherwise. This preserves the previous default physical model
while making an explicit |solid| choice effective on circuit legs.

|stranded| homogenises fine wires, suppresses resolved eddy currents and retains
conductivity and |winding fill| for resistance and copper loss. A solid circuit
leg instead resolves its nonuniform current density. Such a leg requires
|turns=1| or |turns=-1|, |winding fill=1|, no |conductor length| override, and
positive frequency and conductivity. Each solid leg must be connected. Planar
coupled problems require positive |depth|. In axisymmetry, constrained solid
rings must stay strictly away from $r=0$.
\end{key}

Planar voltage-driven or parallel branches require closed signed turns
($\sum N=0$) so that their linkage and voltage do not depend on the potential
gauge. An open leg remains usable with prescribed series current, but its
terminal linkage, voltage and impedance are not reported. Conductivity and
conductor lengths must be known for voltage or parallel excitation. Unspecified
conductivity is not silently interpreted as an ideal zero-resistance winding.

\subsection{Measurements and plots}
\Needspace{7\baselineskip}
\begin{key}{/femm/measure/winding=\meta{name} (initially empty)}
For |quantity=circuit| in a harmonic model, restrict the measurement to one
named branch of |circuit|. Empty means the whole port. This key resets for each
measurement; an unknown name or a static branch request raises an error. The
result is saved under |name| and may be used by |\femmvalue| in an annotation
or by |\femmresult| when building a PGFPlots table.
\end{key}
\begin{codeexample}[code only]
\femmmeasure[name=branch,quantity=circuit,circuit=supply,winding=lower]
\femmvalue[component=current_rms]{branch}
\edef\CurrentPhase{\femmresult{branch}{current_phase}}
\end{codeexample}
Available terminal phasors are |current| (A), |voltage| (V), |linkage| (Wb)
and |impedance| ($\Omega$), each with suffix |_real|, |_imaginary|,
|_amplitude| and |_phase|. Phase is in degrees and is undefined for an exactly
zero phasor. Current, voltage and linkage also have |_rms|; impedance does not.
Existing |resistance| is the DC equivalent resistance, |power| is active input
power (W), and |reactive_power| is reactive input power (var). Impedance is
undefined at zero port current. Unavailable components raise an error rather
than fabricate a zero. These conventions also apply to legacy harmonic
series-current winding measurements.

For parallel branches the reported linkage is an equivalent port linkage,
$\lambda_{\rm eq}=R_{\rm eq}\sum_b\lambda_b/R_b$, where
$R_{\rm eq}^{-1}=\sum_bR_b^{-1}$. It satisfies
$\widehat V=R_{\rm eq}\widehat I+\mathrm i\omega\widehat\lambda_{\rm eq}$;
it is not the sum of branch linkages. Select a branch when that is the quantity
needed for a winding diagram.

Field profiles retain |J| and |Jeddy| and their real, imaginary, amplitude,
RMS, phase and snapshot representations. In a constrained solid,
|Jeddy| is the zero-net-current redistribution about the DC transport profile.
It includes the compensating induced voltage, making it invariant under
$A\mapsto A+\hbox{constant}$ in planar geometry. For an isolated conductor it
is the entire current. In an unconstrained solid it retains the previous
zero-imposed-voltage definition. |femm eddy currents| uses that same field.
Regional |joule_loss| integrates the total current; |eddy_loss| integrates the
redistribution. For a uniform solid, their difference is its DC transport
loss. Stranded copper loss uses the fill-adjusted conductivity.

\subsection{Coupled equations and verification}
For each solid leg let $\ell=d$ in a planar problem and $\ell=2\pi r$ in
axisymmetry. Its solved voltage $\widehat V_l$ gives
\[
 \widehat J=\frac{\sigma}{\ell}\widehat V_l
                 -\mathrm i\omega\sigma\widehat A,\qquad
 G_l\widehat V_l-\mathrm i\omega\int_{S_l}\sigma\widehat A\,dS
       =s_l\widehat I_b,\qquad
 G_l=\int_{S_l}\frac{\sigma}{\ell}\,dS.
\]
For an isolated conductor the right side is zero. A branch voltage sums its
signed solid-leg voltages and the stranded contributions
$R\widehat I+\mathrm i\omega\widehat\lambda$. The field mass, load, linkage,
current and loss integrals share quadrature. Response columns of the sparse
complex field operator reduce the terminal equations to a small dense complex
system with scaling, pivoting and a backward-residual check. The recombined
field is checked against the requested true-residual tolerance. Response columns
use a tighter tolerance; the reported GMRES iteration count adds their work.
Very fine meshes may exhaust the per-solve iteration limit; this raises an error
instead of storing an unconverged result.

The native cache stores the real and imaginary potentials, then reconstructs
terminal currents and conductor voltages on restoration. Revision 13 includes
current-constraint semantics and invalidates older numerical snapshots. A
changed terminal drive can reuse a compatible mesh; frozen reads neither solve
nor modify the cache. Source, phase and conductor choices enter the solution
fingerprint.

The regression suite checks the analytic transport-current slab under mesh
refinement, isolated-conductor gauge invariance, voltage/current equivalence,
parallel currents against independently solved unit-current ports, axisymmetric
shorted rings, active and reactive power identities, rejected combinations and
cache reconstruction. The optional xfemm oracle independently meshes and solves
transport-current and isolated-conductor fixtures, comparing complex A/B/H/J,
energy and loss with unchanged 2\% scalar and 4\% derivative gates. The generic format matrix checks identical restored
terminal results under plain LuaTeX, LuaLaTeX and ConTeXt MkIV.

The circuit interface does not support arbitrary external RLC netlists, nonlinear
large-signal harmonic laws, time-domain hysteresis, transient motion, or public harmonic
|.fem|/|.ans| exchange. Captured incremental operating-point workflows retain
their prescribed series-current interface. Their bias is not redefined by a
new AC voltage constraint.

% SPDX-License-Identifier: LPPL-1.3c
% Copyright (C) 2026 Christophe Jorssen
% Author and Current Maintainer: Christophe Jorssen <christophe.jorssen@gmail.com>
% LPPL maintenance status: maintained
% This file is part of luafemm. See LICENSE and LICENSES.md for its terms.
\section{Complex permeability and laminated-core losses}
\label{sec:losses}

A conducting core can dissipate energy in two different ways: through resolved
currents in the body, and through its magnetic constitutive law. A stack of
insulated sheets adds a third scale: currents circulate within each sheet,
although they cannot travel freely across the insulation. Meshing every sheet
is often unnecessary. This chapter supplies a linear effective law, including
its loss decomposition, and retains the peak $\exp(\mathrm i\omega t)$ convention.

\subsection{A thinner stack, a different frequency response}
Imagine replacing a core made from 0.5 mm sheets by an otherwise identical
stack of 0.2 mm sheets. Give both stacks the same intrinsic material and the
same stacking factor. A material style is the natural place for these choices:
\begin{codeexample}[code only]
\tikzset{femm/materials/core/.style={
  mu r=300,conductivity=2000000,loss angle=10,
  lamination=stacked,lamination fill=.95,lamination thickness=.5}}
\end{codeexample}
The permeability magnitude is 300, the conductivity is $2\,\mathrm{MS/m}$,
and the intrinsic loss angle is $10^\circ$. These are illustrative parameters,
not a fitted electrical-steel grade. |lamination thickness| is in millimetres,
even when the drawing uses another model unit. The stacking factor is the
fraction of magnetic sheet material; the remaining volume is nonmagnetic
insulation with relative permeability one.

First choose a coarse core mesh and a short list of frequencies. Once the
geometry, signs and circuit are satisfactory, refine the core and the air
around it. A sheet thickness is a \emph{material parameter}: it does not set
triangle size. Mesh refinement resolves variations of the averaged field,
while the effective law resolves the one-dimensional variation through a sheet.
The following complete example uses ordinary PGF loops, named measurements
and |\femmresultplot| to compare both stacks and separate the losses.

\luafemmexample{harmonic-materials}

The excitation is a prescribed peak current at every frequency. Field amplitudes
therefore change with frequency and thickness. Do not interpret the loss curves
as a fixed-flux material-characterisation test. The contour levels in the two
pictures are chosen independently. Microscopic sheet currents are not drawn by
|femm eddy currents|: they have been homogenised and are reported as a loss,
not as a fictitious resolved current density in the drawing plane.

\subsection{Material keys and their scope}
All keys below belong to |/tikz/femm/materials| and may appear inside ordinary
material styles. Material styles are evaluated afresh for each captured region.
Optional empty values inherit a supported catalogue property, then the stated
native default. They are not region keys. Declare all materials and regions
before a solve, measurement, field pic or cached solution is requested.

\Needspace{8\baselineskip}
\begin{key}{/tikz/femm/materials/loss angle=\meta{degrees} (initially empty)}
Common intrinsic magnetic loss angle. The native default is zero. For a linear
material, its relative permeability is
$\mu_r^*=|\mu_r|\exp(-\mathrm i\theta)$: the flux density lags the field.
Accepts PGF expressions and finite values $0\leq\theta<90^\circ$.
|mu r|, or the principal |mu x| and |mu y|, supplies the \emph{magnitude},
not the real part. A negative angle would represent an active medium and is
rejected. This parameter does not reconstruct a time-domain hysteresis loop.
\end{key}
\Needspace{8\baselineskip}
\begin{key}{/tikz/femm/materials/loss angle x=\meta{degrees} (initially empty)}
\begin{key}{/tikz/femm/materials/loss angle y=\meta{degrees} (initially empty)}
Principal-axis loss angles. An unspecified angle inherits |loss angle|, whose
native default is zero. Explicit principal angles take precedence over a common
angle. An explicit common angle also replaces catalogue principal defaults;
explicit principal overrides remain effective. |material angle| rotates both
real and imaginary constitutive tensors together in the field plane.
In axisymmetric geometry, the unrotated axes are $r$ and $z$.
\end{key}
\end{key}
\Needspace{8\baselineskip}
\begin{key}{/tikz/femm/materials/lamination=\meta{choice} (initially inherit)}
A choice handler with |inherit|, |none|, |stacked|, |parallel x| and |parallel y|.
|inherit| uses the catalogue interpretation when present, otherwise |none|.
|none| describes a bulk material: its conductivity participates in resolved
eddy-current diffusion unless the region is a stranded winding.

|stacked| describes insulated sheets parallel to the two-component magnetic
field plane, with stacking normal to that plane. Both in-plane principal
permeabilities receive the effective-sheet correction. It suppresses bulk
transport conductivity in the finite-element equation, because the sheet
currents are already included in the effective law. At zero thickness it
means the ideal infinitely fine insulated-sheet limit. It never means a solid
conducting block of thickness zero.

The existing |parallel x| and |parallel y| choices retain their static mixing
interpretation and remain rejected in harmonic problems. |material angle|
rotates principal axes within the sheets; it does not tilt the sheet normal.
In axisymmetry the sheet approximation is local to the meridional field plane;
check that it is appropriate for the physical construction.
\end{key}
\Needspace{7\baselineskip}
\begin{key}{/tikz/femm/materials/lamination thickness=\meta{millimetres} (initially empty)}
Nonnegative thickness of one conducting sheet, excluding insulation. The native
default is zero. Supply a number in millimetres, not a TeX dimension.
A positive value requires |lamination=stacked|. The value is
independent of |mesh size|, picture scale and model length units. A zero value
or zero conductivity removes microscopic sheet-current loss, while intrinsic
magnetic loss may remain. No individual insulation layer is added to the TikZ
geometry or mesh.
\end{key}
\Needspace{7\baselineskip}
\begin{key}{/tikz/femm/materials/lamination fill=\meta{fraction} (initially empty)}
The native default is one; supported catalogue values may supply another value.
Require $0<f\leq1$. For |stacked|, $f$ weights the magnetic sheets and $1-f$
the nonmagnetic insulation. It is unrelated to |winding fill|, which describes
copper occupancy in a stranded winding. A value other than one still requires
a lamination choice. Static |stacked| materials use the zero-frequency arithmetic
mixture along both in-plane axes.
\end{key}

The model is linear. Native stacked sheets require linear intrinsic permeability.
Loss models on captured nonlinear operating points and in the separate axial
$H_z$ diffusion formulation are rejected. A homogenised sheet region cannot
also prescribe an impressed current, belong to a transport circuit, use a
stranded conductor interpretation, or request an isolated-current constraint.
Drive it magnetically from a separate winding or from boundary data.

At frequency zero, loss angles and sheet thickness do not create losses; the
static solve uses the real intrinsic permeability and static fill mixture.
Native cache signatures retain these properties. Static FEM export rejects
new loss metadata instead of silently stripping it. Public harmonic FEM/ANS
interchange is still outside the supported subset.

\subsection{Which loss should be plotted?}
Harmonic regional measurements expose the following SI components. All are
cycle averages in watts and work with selections, saved |name| values,
|\femmvalue| and |\femmresultplot|. No additional integration key is needed.
\begin{center}
\begin{tabular}{lp{.63\linewidth}}
\toprule
Component & Meaning\\
\midrule
|joule_loss| & Resolved bulk-current loss and stranded-winding copper loss.\\
|eddy_loss| & The existing resolved induced-current or zero-net redistribution loss; it is a part of the electrical loss, not an extra addend.\\
|hysteresis_loss| & Dissipation of the intrinsic complex permeability, integrated through the microscopic sheet field when applicable.\\
|lamination_loss| & Microscopic Joule loss inside insulated conducting sheets.\\
|magnetic_loss| & |hysteresis_loss| plus |lamination_loss|.\\
|total_loss| & |joule_loss| plus |magnetic_loss|.\\
\bottomrule
\end{tabular}
\end{center}
For an ordinary lossless material, the new magnetic components are zero.
For a homogenised sheet with no transport source, |J|, |Jeddy| and |joule_loss|
are zero at the resolved scale, even when |lamination_loss| is positive.
Do not add |eddy_loss| to |total_loss|, and do not infer that a zero resolved
current means zero heating. The intrinsic term is called |hysteresis_loss|
for consistency with magnetic engineering terminology; its constant-angle
linear law is not a microscopic hysteresis model.

The existing energy observable is
$\tfrac14\int\operatorname{Re}(\widehat{\mathbf B}\cdot
\widehat{\mathbf H}^{\,*})\,dV$.
The historical |coenergy| component remains an alias of this average; it is
not a conservative potential from which to differentiate harmonic forces.
In a source-free exterior with zero Dirichlet potential, driven circuit ports
satisfy
\[
 \sum_b\tfrac12\operatorname{Re}(\widehat V_b\widehat I_b^*)
      =P_{\rm total},\qquad
 \sum_b\tfrac12\operatorname{Im}(\widehat V_b\widehat I_b^*)
      =2\omega W.
\]
Boundary excitation or impressed sources outside the measured ports contribute
additional work. These identities check assembly and integration; they do not
replace mesh/domain convergence or material calibration. The energy observable
is not a general thermodynamic energy formula for an arbitrary dispersive medium.
Profiles of |Hx| and |Hy|, including phase, RMS and instantaneous views, use the
full complex law rather than applying a real permeability to each phasor part.

\subsection{Effective law, microscopic losses and numerical checks}
For one principal direction let $\mu^*=\mu_0\mu_r^*$ be the intrinsic absolute
permeability and let $d$ be the sheet thickness in metres. Define
\[
 k=\sqrt{\mathrm i\omega\sigma\mu^*},\quad
 z=kd/2,\quad F(z)=\frac{\tanh z}{z},\quad
 \mu_{r,\mathrm{eff}}^*=1-f+f\mu_r^*F(z).
\]
The passive square-root branch has nonnegative real part. For $z=0$ the
analytic limit is $F(0)=1$. The solver assembles the inverse effective tensor,
including its imaginary stiffness, and the mass term only for resolved bulk
conductivity. Axisymmetric integration retains the usual radial weights.

The intrinsic loss is not obtained by subtracting a second finite-element
solve. The exact sheet field
$H(s)=H_s\cosh(ks)/\cosh(kd/2)$ supplies its mean square. With
$Q=d^{-1}\int_{-d/2}^{d/2}|H(s)/H_s|^2\,ds$, the two physical loss densities are
\[
 p_h=\frac{f\omega\mu_0(-\operatorname{Im}\mu_r^*)}{2}Q|H_s|^2,
 \qquad
 p_e=\frac{f}{2\sigma d}\int_{-d/2}^{d/2}
                    \left|\frac{dH}{ds}\right|^2ds.
\]
Their sum equals
$\tfrac\omega2\operatorname{Im}(\widehat{\mathbf H}\cdot
\widehat{\mathbf B}^{\,*})$ per homogenised volume. The implementation obtains
the sheet term from this power identity, with an explicit passivity check.
It uses a small-argument series and decaying exponentials to avoid cancellation
and overflow in the thin-sheet and deep-skin limits.

\nolinkurl{tests/test_losses.lua} checks affine complex anisotropy and rotation,
independent numerical integration of the exact microscopic field, extreme
skin-depth limits, phase covariance, coupled planar/axisymmetric power balance,
cache reconstruction and invalid combinations. The optional
\nolinkurl{scripts/losses-interop.py} compares independently meshed xfemm solutions
for intrinsic loss, ordinary and deep-skin sheets, axisymmetry, and a core
excited by neighbouring currents. This follows the linear |LamType=0| and
|Theta_hx|/|Theta_hy| conventions in xfemm's harmonic solver; unsupported
in-plane stacking choices are not guessed.

The full catalogue remains available. Supported linear |LamType=0| records
map principal permeability, phase angles, conductivity, sheet thickness and
fill into this law. Catalogue B/H curves retain their nonlinear static meaning:
adding this linear loss model does not turn them into nonlinear harmonic laws.
The supplied catalogue currently has no linear record with nonzero sheet/loss
metadata; the example therefore defines its assumptions explicitly. Detailed
wire/strand models and nonlinear harmonic approximations remain unsupported models.

Cache revision 13 fingerprints the intrinsic axes and their orientation, loss angles, fill and
thickness, in addition to frequency. Older snapshots must be recomputed.
Compatible meshes can be reused after a constitutive change, and frozen caches
reconstruct the effective tensors without solving again. Generic cold/warm/shared
frozen tests exercise these properties in plain LuaTeX, LuaLaTeX and ConTeXt MkIV.

% SPDX-License-Identifier: LPPL-1.3c
% Copyright (C) 2026 Christophe Jorssen
% Author and Current Maintainer: Christophe Jorssen <christophe.jorssen@gmail.com>
% LPPL maintenance status: maintained
% This file is part of luafemm. See LICENSE and LICENSES.md for its terms.
\section{A small excitation around a saturated operating point}
\label{sec:operating}

\subsection{Kamil asks which inductance his students should use}
Kamil has plotted the flux linkage of a coil as its current increases.
The curve bends as the iron saturates. He would now like to connect that
large-signal curve to the inductance seen by a small additional excitation.
Dividing flux linkage by current gives a secant slope; a small perturbation
sees the tangent slope. These are different questions, even when the geometry
and the winding have not changed.

An \emph{operating point} captures a converged static solution and its exact
ordered mesh. A second picture inherits the geometry and material declarations,
but uses a linear constitutive law evaluated at that saved state. Two modes
are available: |frozen| uses the local secant permeability, while
|incremental| uses the differential tensor. Both planar and axisymmetric
static models are supported. Section~\ref{sec:biased-harmonic} extends incremental
mode to a small harmonic excitation. Neither mode advances time or moves the rotor.

\subsection{First solve, then capture, then perturb}
The order matters. Kamil declares his materials, boundaries and winding in
an ordinary |femm/problem|. At picture level he adds:
\begin{codeexample}[code only]
\tikzset{biased coil/.style={
  femm/problem={xmin=-30,xmax=30,ymin=-25,ymax=25,depth=20},
  femm/circuit={name=drive,current=5},
  femm/operating point={name=bias},
  femm/cache={file=coil-bias,mode=auto}}}
\end{codeexample}
The |femm/operating point| key schedules the capture at the end of that
picture, after all its physical paths have been declared. It solves if
necessary, then takes an independent copy. The snapshot is available only
after the picture has ended. Alternatively,
|\femmoperatingpoint[name=bias]| performs the capture immediately at its
position in the input; place it after the last physical declaration.
A capture does not replace the current model.

The next picture selects the saved name. It does \emph{not} repeat the
physical paths, material definitions or circuit declarations:
\begin{codeexample}[code only]
\begin{tikzpicture}[
  femm/problem={operating point=bias,permeability=incremental,
    circuit currents={drive=1}},
  femm/cache={file=coil-increment,mode=auto}]
  \pic {femm geometry};
  \pic[femm/field style={orange!75!black}] {femm field};
  \femmmeasure[name=small,quantity=circuit,circuit=drive]
\end{tikzpicture}
\end{codeexample}
Here the current is $\delta I=1\,\mathrm A$. This is a normalisation of a
linear response: the resulting linkage is numerically the differential
inductance in henries. It is not a claim that the nonlinear circuit stays
linear over a finite one-ampere excursion. For an actual small change $\epsilon$,
$\mathbf B(I_0+\epsilon)\simeq\mathbf B_0+
\epsilon\,\partial\mathbf B/\partial I$.

In incremental mode every independent bias source is zeroed: circuit currents,
impressed current densities, coercive sources and imposed boundary loads.
The specified |circuit currents| then supply the perturbation. Unlisted
circuits have zero perturbation current. Robin coefficients, boundary kinds
and periodic ties are inherited; their additive loads are homogeneous.
The resulting lines and profiles describe $\delta\mathbf B$ and
$\delta\mathbf H$, not the total field. Ordinary |femm/profile|,
|\femmplot|, |\femmmeasure| and |\femmresultplot| keep their usual syntax.

For a reconstruction of the bias, Kamil instead writes:
\begin{codeexample}[code only]
\begin{tikzpicture}[
  femm/problem={operating point=bias,permeability=frozen}]
  \pic {femm geometry};
  \pic {femm field};
\end{tikzpicture}
\end{codeexample}
Frozen mode retains every bias source by default. With unchanged sources its
solution recovers the nonlinear bias to solver tolerance. An explicit circuit
current replaces that circuit's current; it is not added to it. This permits
source decomposition when all other sources are accounted for. Retained
magnets or nonzero boundary loads must not be mistaken for zero excitation.

\subsection{One core, two slopes}
The complete example below compares the bias and incremental fields of a
closed rectangular core, then plots its flux linkage and both inductances.
The B--H table is illustrative, not a measured steel specification. The field
panels deliberately use separate contour levels; counting lines across the
two panels does not measure their flux ratio.

Start with |\Mesh=3| while arranging the figures, then try |1.5| and a finer
value before interpreting the curves quantitatively. Each bias solve has its
own mesh, but its frozen and incremental responses reuse that mesh exactly.
Within a bias-response pair there is no remeshing or transfer error. Across a
current sweep, geometry-driven meshing is deterministic for unchanged inputs.
Convergence of the tangent solve on one mesh does not establish convergence
of the spatial discretisation.

\luafemmexample{operating-point}

The result component |inductance| is $\lambda/I$ for the ordinary static
model under its existing eligibility rules. On an incremental response with
one nonzero circuit excitation it is $\delta\lambda/\delta I$; the explicit
alias |incremental_inductance| names the same value. Kamil uses the alias in
his plot to make its interpretation unambiguous. Its SI unit is the henry:
\begin{codeexample}[code only]
\femmvalue[component=incremental_inductance]{small}
\femmresultplot[results={small-a,small-b,small-c},
  component=incremental_inductance,abscissa=parameter,
  plot options={orange!75!black,thick,mark=square*}]
\end{codeexample}
Record |parameter=| as the bias current, as the complete example does.
The result's |current| is the perturbation current, not the bias current.

For several windings, excite one port at a time and measure the linkage of
all ports. With a unit excitation at port $j$, the linkage at port $i$ gives
$L_{ij}$. Unexcited ports still report linkage but cannot report a self
inductance by division by zero. When several circuits are excited together,
the response does not report an |inductance| ratio: use the individual
linkages and the declared excitation vector. A matrix is assembled by the
caller; there is no automatic matrix-valued result or circuit-network solver.
For frozen responses the ratio is also omitted in the presence of other
independent sources. In either mode a planar winding must be closed in signed
turns for its linkage to be gauge independent.

\subsection{Key reference}
\begin{key}{/tikz/femm/operating point=\marg{options}}
Schedule a capture at the end of the originating picture. Options belong to
|/femm/operating point|. Use once for each named snapshot; names are
unique across the document. Reusing a name is an error, which prevents a
later picture from silently changing the meaning of an earlier response.
\end{key}
\begin{command}{\femmoperatingpoint\oarg{options}}
Capture immediately, solving first if required. This command and the picture
key have the same options. A capture owns its data and does not borrow a
mutable potential array from the original model.
\end{command}
\begin{key}{/femm/operating point/name=\meta{name} (initially empty)}
Required nonempty snapshot identifier. Prefer letters, digits and hyphens in
examples. A snapshot name is separate from a measurement or profile name.
\end{key}
\begin{key}{/femm/operating point/model=\meta{model name} (initially empty)}
Select a registered |/femm/problem/name| for immediate capture; empty means
the current model. The picture-level key supplies its own originating model
identifier automatically, so another picture cannot redirect that capture.
\end{key}
\begin{key}{/femm/problem/operating point=\meta{name} (initially empty)}
Start a response from a previously captured operating point. Empty starts an
ordinary problem. Geometry, units, problem type, material laws, solver settings,
depth and mesh all come from the snapshot. Explicit geometry or solver options
in the same problem declaration are rejected, rather than silently ignored.
Set those options on the bias problem. |name| may name the new response model;
|femm/cache| and presentation keys remain independent. The one physical override
is |frequency|, for the incremental harmonic extension described next.
\end{key}
\begin{key}{/femm/problem/permeability=\meta{choice} (initially incremental)}
A PGF choice, either |incremental| or |frozen|. It selects the linear material
law and the source defaults explained above. It is distinct from
|/femm/cache/mode=frozen|, which only requires an existing disk cache.
\end{key}
\begin{key}{/femm/problem/circuit currents=\marg{name=value list} (initially empty)}
Replace selected circuit currents, in amperes, in the inherited problem.
For example |circuit currents={drive=1,sense=0}|. PGF evaluates the numerical
expressions, so |drive=1/2| is allowed. Circuit names must already exist in
the bias. An empty list retains the mode's defaults; in incremental mode
this gives a zero response. Use this key instead of redeclaring inherited
circuits with |femm/circuit|. The enclosing problem remains styleable:
\begin{codeexample}[code only]
\tikzset{unit response/.style={femm/problem={
  operating point=#1,permeability=incremental,
  circuit currents={drive=1}}}}
\end{codeexample}
\end{key}

\subsection{What is actually linearised?}
Let the isotropic material table describe the noncoercive law
$\mathbf H_{\mathrm{tab}}=\nu(B)\mathbf B$, where $\nu(B)=H(B)/B$ in SI.
At the static bias $\mathbf B_0$, define $B_0=|\mathbf B_0|$ and
$\mathbf n=\mathbf B_0/B_0$. Frozen permeability uses the scalar reluctivity
$\nu_0=H(B_0)/B_0$. The incremental law is
\[
 \delta\mathbf H=\boldsymbol\nu_{\mathrm{inc}}\,\delta\mathbf B,
 \qquad
 \boldsymbol\nu_{\mathrm{inc}}=
 \nu_0\mathbf I+
 \left(\left.\frac{\mathrm dH}{\mathrm dB}\right|_{B_0}-\nu_0\right)
 \mathbf n\mathbf n^{\mathsf T}.
\]
Thus parallel and perpendicular perturbations see different slopes. At zero
bias the initial table slope is used in both directions. An already linear
tensor retains its tensor in both modes. Permanent-magnet coercivity belongs
to the source term, not to this derivative; incremental mode removes that
source while preserving the recoil law.

The implementation uses the same B--H interpolation and material tangent as
the nonlinear solver. At a corner of a piecewise-linear B--H table the
classical derivative is not unique: the selected interval's tangent is used.
Avoid interpreting a centred finite difference crossing such a corner as an
exact derivative check. Smooth FEMM-style interpolation is available, but
its smoothing changes the effective input curve as described in the material
chapter.

For planar P1 elements the bias curl and tangent are constant in each
triangle. For axisymmetry, $B_r=-\partial_zA_\phi$ and
$B_z=\partial_rA_\phi+A_\phi/r$ vary within an element. The bias and material
tangent are therefore evaluated at each volume quadrature point, with the
same radial weights as the nonlinear solve. Pointwise H profiles use the
local tensor at the sampled point. The static weak operator and its signed
boundary constraints are otherwise unchanged.

Regional |energy| and |coenergy| on an incremental response both mean the
quadratic response energy $\tfrac12\int\delta\mathbf B\cdot
\delta\mathbf H\,\mathrm dV$. They exclude the bias energy and the first-order
cross term in a total-energy expansion. For a single unit-current response,
$2W=L_{\mathrm{inc}}$. In frozen mode the energy refers to the frozen linear
material and retained coercive source; it is not the original nonlinear
stored energy even when the field reconstruction agrees.

\subsection{Persistence, checks and current limits}
Named operating points live for the TeX run. Across compilations, keep the
bias declarations and use their ordinary native solution cache: capture can
then copy the restored solution without repeating Newton iterations. Give
each response its own cache filename. Passive cache format 13 includes the
ordered bias mesh, nodal bias field, constitutive mode, source provenance and
response excitations. A changed bias cannot reuse an old response silently.
Older caches must be recomputed. There is no transfer between unrelated meshes
and no reader for a native FEMM |.ans| operating point yet.

The operating-point model accepts native planar/axisymmetric static models with
single-valued isotropic nonlinear or positive linear tensor laws, including
supported permanent magnets. Named circuits must use prescribed series
currents. Source/boundary callbacks, imported operating points, ideal confined
models, voltage/parallel circuit drives and capture of an already linearised
response are rejected. Geometry and physics are locked after creating a
response: draw labels or untagged decorative paths, but do not add regions,
change materials or redeclare boundaries. Create a new nonlinear bias for a
changed geometry. Small harmonic responses are described in
Section~\ref{sec:biased-harmonic}. Force/torque derivatives, combination with
sliding-position models and FEM/ANS export of response models remain unsupported
and fail explicitly. Independent static rotor positions are described in
Section~\ref{sec:sliding}.
The nonlinear bias itself retains its ordinary supported export capabilities.

The regression suite checks frozen reconstruction, the analytic oblique
material tensor, permanent-magnet source removal, zero/linear limits,
reciprocity, response energy and finite differences of planar and
axisymmetric nonlinear solves. The optional xfemm study compares against
independent ordinary nonlinear solves; its pinned command-line solver rejects
DC previous-field modes. These checks validate the stated static response
model, not a complete FEMM previous-solution workflow. Detailed tolerances
and measured errors are recorded in \nolinkurl{docs/validation.md}.

\subsection{Lua entry points}
The numerical module |require("luafemm-operating")| does not access TeX globals.
|capture(model)| returns an owned snapshot of a converged static problem;
|new(snapshot,options)| constructs an unsolved response. Options are |mode|,
|currents| (a table from circuit names to amperes), |frequency| (hertz),
|phases| (a table from circuit names to degrees) and |cache| (the ordinary
native cache options). They default to |incremental|, empty current/phase tables,
zero frequency and caching off. Unknown options are rejected. For example:
\begin{codeexample}[code only]
local O = require("luafemm-operating")
local bias = O.capture(model) -- model has already converged
local response = O.new(bias, {
  mode = "incremental", currents = {drive = 1},
  cache = {file = "unit-response", mode = "auto"}
})
require("luafemm").solve(response)
local r = require("luafemm-static").result(response, "drive")
-- r.incremental_inductance is in henries when available.
\end{codeexample}
A snapshot may also come from a compatible restored native cache. The snapshot
itself is an in-memory object; |capture| does not create an extra disk format.

% SPDX-License-Identifier: LPPL-1.3c
% Copyright (C) 2026 Christophe Jorssen
% Author and Current Maintainer: Christophe Jorssen <christophe.jorssen@gmail.com>
% LPPL maintenance status: maintained
% This file is part of luafemm. See LICENSE and LICENSES.md for its terms.
\section{A small alternating excitation on a DC bias}
\label{sec:biased-harmonic}

\subsection{Kamil adds a ripple to his solenoid}
The static tangent in Section~\ref{sec:operating} tells Kamil how much extra
flux a small extra current produces. His conducting iron rod presents another
question: does that flux still reach the centre when the excitation oscillates?
The changing flux induces azimuthal eddy currents, which oppose penetration
and dissipate energy. Moreover, a strongly biased rod has a different tangent
permeability from a weakly biased one. Frequency and saturation must therefore
be considered together.

Kamil keeps his nonlinear \emph{DC} calculation and adds a small sinusoidal
current. The package solves the complex linear response about that fixed DC
state. It does not repeatedly solve the nonlinear law along a finite hysteresis
loop. This distinction determines both the input and the interpretation of the
plots: all response fields exclude the DC field.

\subsection{Reuse the bias, choose a frequency}
The geometry and the capture mechanism are exactly those of the preceding
chapter. Kamil first declares the conducting material, the coil and its DC
current, then captures the converged solution at the end of its picture.
In a second picture he selects the captured name and adds |frequency|:
\begin{codeexample}[code only]
\tikzset{ripple/.style 2 args={femm/problem={
  operating point=#1,permeability=incremental,frequency=#2,
  circuit currents={drive=1},circuit phases={drive=0}}}}
% The preceding DC picture has captured a point named rod-low.
\begin{tikzpicture}[ripple={rod-low}{200},
  femm/cache={file=rod-low-200,mode=auto}]
  \pic {femm geometry};
  \pic[femm/phase=0] {femm field};
  \pic[femm/currents={view=symbols,phase=0}] {femm eddy currents};
  \femmmeasure[name=terminal,quantity=circuit,circuit=drive,parameter=200]
  \femmmeasure[name=losses,parameter=200]
\end{tikzpicture}
\end{codeexample}
The |drive| circuit already belongs to the bias. Its current here is the
\emph{peak AC perturbation}, not the sum of DC and AC currents. A unit peak
current conveniently normalises the linear transfer functions. An actual
small peak amplitude $\epsilon$ multiplies the fields by $\epsilon$ and the
reported AC powers by $\epsilon^2$; it does not change impedance or inductance.
Unlisted circuits have zero AC current.

The two phase settings have different jobs. |circuit phases| changes the
physical source phasors and therefore the solve/cache signature.
|femm/phase|, or |phase| inside |femm/currents|, chooses the instant drawn from
an existing complex solution. It does not require another solve. All physical
phasors use the peak convention $\operatorname{Re}(\widehat u e^{i\omega t})$;
phases are in degrees.

\subsection{Two bias currents, one frequency sweep}
The complete example uses an axisymmetric rod of radius $5\,\mathrm{mm}$ and
length $24\,\mathrm{mm}$ inside a 100-turn stranded solenoid. The illustrated
half-section is rotated around $r=0$; its eddy currents close around that axis.
The conductivity is $2\,\mathrm{MS/m}$ and the B--H table is an illustrative
single-valued law, not a measured steel specification. The coil carries either
$5\,\mathrm A$ or $50\,\mathrm A$ DC. Each captured bias then supplies six
AC responses between 5 and $200\,\mathrm{Hz}$.

While arranging the drawing, try |\CoreMesh=.6| and only two frequencies.
Then use |.25|, followed by finer values, before interpreting small differences
in the curves. Refinement belongs to the \emph{bias} geometry: every response
inherits that exact mesh. A mesh sufficient for the DC field need not resolve
AC penetration. An approximate local skin depth is
$\delta=\sqrt{2/(\omega\mu_{\mathrm{diff}}\sigma)}$; in a nonuniform tensor
material it is only a guide. Check the smallest relevant penetration length
at the highest frequency and weakest differential reluctivity, as well as the
outer air boundary. Increasing contour counts or profile samples does not
refine the finite elements.

\luafemmexample{biased-harmonic}

Kamil should read the line panels as instantaneous \emph{response} fields.
Their levels are chosen independently. A denser-looking panel alone does not
prove a larger flux. The dot/cross symbols indicate current direction normal
to the meridional drawing, with the convention of Section~\ref{sec:currents}.
The inductance and loss curves provide quantitative comparisons with explicit
units. The |eddy_loss| component includes the resolved rod currents and excludes
stranded-coil copper loss. Neither curve includes DC heating.
For these parameters the stronger bias lowers the real differential inductance,
but its rod loss becomes larger as frequency rises. A lower tangent permeability
also changes penetration and the current distribution: saturation alone does
not imply a monotone decrease of eddy-current loss.

\subsection{Complex inductance, impedance and power}
For a single excited port define
\[
 L_{\mathrm{inc}}^*=\frac{\widehat\lambda}{\widehat I},\qquad
 Z=R+i\omega L_{\mathrm{inc}}^*.
\]
The superscript star here labels a complex inductance; it is not a conjugation.
The linkage includes the response of the conducting body. Consequently
\[
 \operatorname{Re}Z-R=-\omega\operatorname{Im}L_{\mathrm{inc}}^*,
 \qquad \operatorname{Im}Z=\omega\operatorname{Re}L_{\mathrm{inc}}^*.
\]
With the package's sign convention a passive eddy-current load normally has
negative imaginary inductance. At low frequency the real component approaches
the static differential inductance, provided the same bias and mesh are used.
Do not replace it by $\lambda_0/I_0$, which is the static secant ratio.

% Keep the key heading with its description after the full-page figure.
\newpage
\begin{key}{/femm/problem/circuit phases=\marg{name=value list} (initially empty)}
Set selected inherited circuit phases in degrees. PGF evaluates each numerical
expression. For example, with both circuits already declared in the bias:
\begin{codeexample}[code only]
circuit phases={drive=0,sense=180/2}
\end{codeexample}
Names must exist in the captured bias. This key requires |operating point|;
a nonzero phase requires positive |frequency|. Unlisted circuits retain their
DC phase of zero. An empty list changes no phase. Source amplitudes are set
separately by |circuit currents|. Both lists are ordinary styleable PGF keys.
\end{key}

For an inherited problem, |/femm/problem/frequency| is the one allowed physical
problem override. Its default remains zero (the static response); a positive
value in hertz selects the harmonic response and requires
|permeability=incremental|. |permeability=frozen| with positive frequency is
rejected: its retained DC sources must not silently become AC sources. Other
problem settings, material laws, geometry, boundaries and solver tolerances
come from the bias. Set them before capture.

The circuit result adds two components, both in henries:
\begin{itemize}
\item |incremental_inductance_real| gives the real part;
\item |incremental_inductance_imaginary| gives the imaginary part.
\end{itemize}
They are available only on
an incremental harmonic model with exactly one nonzero circuit excitation,
and only for that excited circuit. A zero-current sensing port still reports
complex linkage. With several simultaneously excited ports, use those linkages
and the declared current vector; no scalar self-inductance ratio is reported.
To build a small-signal matrix, solve one excitation at a time. The existing
|impedance_real|, |impedance_imaginary|, voltage, resistance and complex linkage
components retain their meanings and eligibility rules.
\begin{codeexample}[code only]
\femmvalue[component=incremental_inductance_real]{terminal}
\femmresultplot[results={port-low-5,port-low-10,port-low-20},
  component=incremental_inductance_real,
  plot options={blue!65!black,thick,mark=*}]
\end{codeexample}
Record the frequency with |parameter| when measuring, as in the complete
example. Profile paths retain the usual syntax; use |representation=amplitude|
for $|\widehat B_z|$, |real| or |imaginary| for phasor components, and |snapshot|
with |phase| for instantaneous fields. A harmonic component amplitude is
nonnegative; it loses the phase information needed to reconstruct a waveform.

The cycle-average response energy and terminal powers are
\[
 \overline W=\frac14\int
 \widehat{\mathbf B}^{\dagger}\boldsymbol\nu_{\mathrm{inc}}
 \widehat{\mathbf B}\,\mathrm dV,\quad
 P=\frac12\operatorname{Re}(\widehat V\overline{\widehat I}),\quad
 Q=\frac12\operatorname{Im}(\widehat V\overline{\widehat I}).
\]
The dagger denotes complex conjugate transpose. For the example's single
driven port and homogeneous Dirichlet/natural boundaries,
$Q=2\omega\overline W$, while $P$ equals the total resolved Joule loss plus
stranded-coil copper loss. With several driven ports, sum their complex powers.
Regional energy excludes any energy represented by a Robin boundary term;
include that term separately when checking reactive power for such a boundary.
Likewise, a prescribed conductor-length override can include resistance of wire
outside the model volume, whose loss is absent from the regional integral.
These are AC response quantities, excluding DC
energy, DC winding heating and bias-response cross terms. The software does
not provide force derivatives or mixed DC/AC force harmonics in this mode.

\subsection{Constitutive model, persistence and limits}
The differential tensor is evaluated from the captured $\mathbf B_0$, using
the formula in Section~\ref{sec:operating}. It is fixed throughout the complex
solve. Both real and imaginary response fields use that same tensor; the
solver never evaluates a new B--H tangent from the AC amplitude. In axisymmetry
it evaluates the local bias curl and tensor at each volume quadrature point,
including the $A_\phi/r$ term. The stiffness is the static differential
operator; the conductive term is $i\omega\sigma A$. In this formulation
$\widehat J_{\rm eddy}=-i\omega\sigma\widehat A$.

Permanent magnets may establish the DC bias. Their coercive source is removed
from the incremental excitation, while their supported recoil law remains.
A real tangent does not model hysteresis, demagnetisation history, complex
permeability or frequency-dependent material losses. Nor is this a nonlinear
large-amplitude harmonic or transient solver.

The captured operating-point interface retains prescribed series-current
windings, real positive material tensors, supported ordinary boundaries, and
no floating-conductor constraints or voltage/parallel drives in these responses,
harmonic homogenised laminations, specialised periodic/open boundaries or
public harmonic FEM/ANS interchange. Unbiased harmonic circuits support the
additional constraints in Section~\ref{sec:ac-circuits}. Here planar solid conductors use zero imposed
longitudinal electric-potential gradient; axisymmetric currents close naturally.
Imported operating points and response-model FEM/ANS export remain unsupported.

Use a distinct cache file for each bias/frequency/excitation. Native cache
format 13 includes the exact bias mesh and potential, mode, frequency and
source amplitudes/phases. Presentation-only phase changes reuse it. The named
bias is still captured in each TeX run, optionally from its own restored DC
cache. There is no new persistent format for the captured name itself.

\subsection{Lua entry point and validation}
The same numerical API accepts frequency and phase overrides:
\begin{codeexample}[code only]
local response = require("luafemm-operating").new(bias, {
  mode = "incremental", frequency = 200,
  currents = {drive = 1}, phases = {drive = 90},
  cache = {file = "biased-200", mode = "auto"}
})
require("luafemm").solve(response)
local port = require("luafemm-harmonic").circuit(response, "drive")
\end{codeexample}
The bias must be a native converged unchanged DC model or its owned capture.
Frequency is a finite nonnegative number in hertz; phases are finite numbers
in degrees. Frequency and phases default to zero and an empty table.

Regression checks cover an analytic biased conducting slab with refinement,
low-frequency limits, phase covariance, permanent-magnet source removal,
power/energy balances and immutable bias/cache data. An optional xfemm oracle
solves independent linear AC problems with analytically derived tangent tensors
at a manufactured uniform bias. It checks the differential constitutive and
complex diffusion operators without claiming support for xfemm's previous-field
file workflow. The generic format matrix exercises the same phasor caches in
plain LuaTeX, LuaLaTeX and ConTeXt MkIV. Quantitative results and tolerances are
recorded in \nolinkurl{docs/validation.md}.

% SPDX-License-Identifier: LPPL-1.3c
% Copyright (C) 2026 Christophe Jorssen
% Author and Current Maintainer: Christophe Jorssen <christophe.jorssen@gmail.com>
% LPPL maintenance status: maintained
% This file is part of luafemm. See LICENSE and LICENSES.md for its terms.
\section{A rotor moves inside a circular air band}
\label{sec:sliding}

\subsection{Kamil turns the rotor}
Kamil has drawn his machine. He now wants to turn its rotor while keeping the
stator current fixed, plot the magnetic field at a few positions, and compare
torque with flux linkage. Redrawing and remeshing the entire machine at every
angle would mix physical changes with changes in the finite-element mesh.
A \emph{moving air band} keeps the bulk triangles attached to their respective
bodies. Only a narrow annulus of air receives new triangles at each position.

This first implementation treats independent \emph{planar static} positions.
The complete nonlinear field is solved afresh at each angle. It introduces
neither angular velocity nor motion-induced voltage. Its triangles are ordinary
P1 finite elements, not FEMM's specialised analytical air-gap elements (AGE).
There can still be discretisation ripple when the band diagonals change.

\Needspace{10\baselineskip}
\subsection{Declare once, capture, then choose positions}
The order matters. First declare the problem, materials, circuits and all
physical paths or nodes. Then capture a mesh-only template. Finally create
one new problem per requested angle and perform the measurements there.
\begin{enumerate}
\item Draw every rotor domain strictly inside an inner circle. Draw every
  stator domain outside an outer circle. Leave a source-free air collar on
  both sides of the circles: neither trace may coincide with a material edge.
\item Attach |femm/sliding band| to the declaration picture. The key captures
  at the end of that picture, after all its paths. Do not request a field,
  mesh, profile value, measurement or operating point before this capture.
\item In a new picture, set |femm/problem={sliding band=machine,rotor angle=30}|.
  This copies the physical data, turns the rotor and connects the band.
\item Draw the inherited geometry/field and measure torque, circuit linkage
  or ordinary profiles. Repeat step 3 for another angle.
\end{enumerate}

\begin{codeexample}[code only]
\tikzset{machine position/.style={
  femm/problem={sliding band=machine,rotor angle=#1},
  femm/cache={file=machine-angle-#1,mode=auto}}}
% A preceding declaration picture captured the template named machine.
\begin{tikzpicture}[machine position=30]
  \pic {femm geometry};
  \pic {femm field};
  \femmmeasure[name=at-30,quantity=force,method=band,parameter=30]
  \femmmeasure[name=port-30,quantity=circuit,circuit=drive,parameter=30]
  \path[femm/profile={name=gap-30,samples=181}]
    (4.8,0) arc[start angle=0,end angle=360,radius=4.8];
\end{tikzpicture}
\end{codeexample}
The picture transformation controls presentation. It does not turn the physical
rotor: that is the job of |rotor angle|. In particular, an untagged rotated
ellipse is just a drawing overlay. Tagged physical paths in the declaration
picture retain their full TikZ expressiveness; the usual flattening tolerance
still controls their computational polygons. The band traces themselves are
regular polygonal approximations of circles.

\Needspace{10\baselineskip}
\subsection{A complete reluctance-machine experiment}
Kamil uses an elliptical rotor and an annular stator with two prescribed
100-turn winding legs. The signs of their turns close the winding, so its
linkage is independent of the potential gauge. The steel table is illustrative
and single-valued. The two field panels use independently chosen contour levels;
compare the measured linkage, rather than visible line counts, for magnitude.
The dashed circles show the artificial band, not a new material interface.

The source starts with a relatively coarse mesh. While adjusting colours and
labels, Kamil uses only two or three angles and |mesh size=1.2|. He then decreases
the bulk spacing, increases |segments| and checks several |layers| values.
For a final quantitative plot he also moves the exterior box farther away.
A smooth-looking interpolated plot is not evidence of converged torque.

\luafemmexample{sliding-band}

At fixed current, the rotor tends towards increasing magnetic coenergy. For
this example it aligns approximately at $90^\circ$ and the torque reverses
there. The geometry repeats after $180^\circ$; the API itself normalises angles
only modulo $360^\circ$. Asymmetry in a discretised template can cause small
differences between nominally equivalent half-turn positions.

The loop in this example is ordinary |\foreach|. Each named result remains
available to |\femmresultplot| after its picture has ended. Store the intended
abscissa explicitly with |parameter|. Changing |femm/lines|, colour, scale or
profile sample count does not change the solved physical position.

\Needspace{10\baselineskip}
\subsection{Capture keys and geometric contract}
\Needspace{7\baselineskip}
\begin{key}{/tikz/femm/sliding band=\marg{key list}}
Schedule capture at the end of this picture, binding it to this picture's model.
The keys belong to |/femm/sliding band|. There is no default template name.
Only one band is supported per template. All physical declarations must precede
the capture and no mesh or solve may already exist. A mesh cache restored
\emph{by the capture itself} is allowed.
\end{key}
\Needspace{7\baselineskip}
\begin{command}{\femmslidingband\oarg{key list}}
Capture immediately. Place this command after the last physical declaration.
It accepts the same key family as the picture key. It does not draw or solve.
Each invocation resets the capture options to the defaults below before applying
its list; package styles can supply recurring settings.
\end{command}
\Needspace{7\baselineskip}
\begin{key}{/femm/sliding band/name=\meta{name} (initially empty)}
Required document-wide template identifier. Duplicate names are errors. The
name remains available across TeX groups for the current run; it is not itself
a saved file or a persistent object identifier.
\end{key}
\Needspace{7\baselineskip}
\begin{key}{/femm/sliding band/model=\meta{name} (initially empty)}
For the immediate command, select an existing named model instead of the current
one. It must still be unmeshed. The picture-level key supplies its own originating
model name after the user's list, so later pictures cannot redirect its capture.
\end{key}
\Needspace{7\baselineskip}
\begin{key}{/femm/sliding band/inner radius=\meta{number} (initially 0)}
Inner trace radius in model units. Required positive value. Every rotating path
must lie strictly inside the inscribed polygon, not merely inside the circle.
\end{key}
\Needspace{7\baselineskip}
\begin{key}{/femm/sliding band/outer radius=\meta{number} (initially 0)}
Outer trace radius in model units, strictly larger than the inner radius.
Every stator path must lie outside this circle; compound stator paths may
surround it but their material must not fill the band. The entire circle must
lie strictly inside the rectangular problem box.
\end{key}
\Needspace{7\baselineskip}
\begin{key}{/femm/sliding band/center x=\meta{number} (initially 0)}
Horizontal coordinate of the rotation centre in model units.
\end{key}
\Needspace{7\baselineskip}
\begin{key}{/femm/sliding band/center y=\meta{number} (initially 0)}
Vertical coordinate of the same centre. These are model coordinates, unaffected
by later drawing transformations. They do not set the torque origin.
\end{key}
\Needspace{7\baselineskip}
\begin{key}{/femm/sliding band/segments=\meta{integer} (initially 96)}
Number of vertices on each initially declared trace and each intermediate ring;
allowed range 16--1024. The template mesher may subdivide the trace chords,
so their final vertex counts can differ. Raising this value refines the angular
resolution and the polygonal geometry; it does not refine the bulk interiors.
\end{key}
\Needspace{7\baselineskip}
\begin{key}{/femm/sliding band/layers=\meta{integer} (initially 2)}
Number of radial strips, from 1 to 16. Intermediate rings are equally spaced
in radius and fixed in the stator frame. With inner radius $r_i$, outer radius
$r_o$, $L$ layers and $N$ segments, the nesting check requires
\[
 \bigl(r_i+(r_o-r_i)/L\bigr)\cos(\pi/N)>r_i.
\]
If it fails, increase |segments| or reduce |layers|. Excessive radial refinement
with insufficient angular refinement creates poor triangles. The ordinary
minimum-angle and maximum-area targets, when nonzero, are checked on every
new position; an unmet target is an error, not a silently accepted mesh.
\end{key}
All six numerical capture keys accept ordinary PGF mathematical expressions.
No Lua source is evaluated from key values. A path crossing or filling the band
is rejected, including a large air refinement rectangle. Use local spacing
on rotor/stator regions instead. Both traces must be in air with vacuum
permeability, zero coercivity and no current. The region order retains its
usual meaning inside each bulk domain.

\Needspace{10\baselineskip}
\subsection{Position and measurement keys}
\Needspace{7\baselineskip}
\begin{key}{/femm/problem/sliding band=\meta{name} (initially empty)}
Create an independent position from the named template. An empty value selects
the ordinary problem workflow. The template owns geometry, materials, circuits,
exterior conditions, mesh targets, units, depth and solver settings. Do not
redeclare them in a position picture. Graphics, profiles, selections and
measurements may still be added. Cache options belong to the new picture.
\end{key}
\Needspace{7\baselineskip}
\begin{key}{/femm/problem/rotor angle=\meta{number} (initially 0)}
Mechanical rotation in degrees, counterclockwise about the band centre, relative
to the captured geometry. PGF expressions are allowed, for example |90/3|.
Negative values are allowed and finite angles are reduced modulo $360^\circ$.
A nonzero angle without |sliding band| is an error. It is not an electrical
angle or a display phase. Ordinary TeX grouping controls problem-key inheritance.
\end{key}
The existing |/femm/problem/circuit currents| list can override prescribed
series currents in amperes. Unlisted currents retain their template values.
This differs from the zero-excitation default of an incremental response.
Unknown circuit names and phase overrides are rejected. The new current and
angle enter the physical cache signature. No other physical problem override
is accepted, even if it would numerically equal the inherited setting.

\Needspace{7\baselineskip}
\begin{key}{/femm/measure/method=band}
Additional choice for |quantity=force| on a sliding position. Integrate Maxwell
stress across the entire air band to obtain the total |Fx|, |Fy| and |torque|
on everything inside it. No body selection or profile is needed or accepted.
The origin remains |origin x=0,origin y=0| in model units; explicitly set it to
the band centre for torque about an offset rotation axis. Forces use newtons,
torque uses newton metres, and positive torque is counterclockwise. The problem
must have positive |depth|. Existing integral components cover the whole model.
\end{key}
The general measurement default remains |method=weighted|. Set |method=band|
explicitly, or define a style:
\begin{codeexample}[code only]
\pgfkeys{/femm/measure/rotor torque/.style={quantity=force,method=band}}
\femmmeasure[rotor torque,name=at-60,parameter=60]
\end{codeexample}
|weighting| and |tolerance| configure the ordinary Laplace-weighted body method;
they do not change the band's prescribed weights. For an independent diagnostic,
register a closed contour strictly inside source-free air and use the existing
|method=maxwell,profile=...|. A single contour can have more radial bias than
the volume integral. Compare several contours and refined meshes.

Ordinary profiles and circuit measurements retain their syntax. For example,
|\femmplot[profile=gap-30,component=By,abscissa=angle]| inside a PGFPlots axis
plots the angular gap field. Set the profile centre explicitly for an offset
machine. Material selections use the original material names even when a
rotated anisotropic rotor internally needs its own tensor copy. Such a material
filter also selects stator regions with that material; add a geometric selection
when isolating a body for the ordinary weighted-force method.

\Needspace{10\baselineskip}
\subsection{Numerical construction and accuracy}
The template is meshed once with two protected polygonal traces. Band triangles
are discarded; retained bulk nodes and triangle connectivity are copied into
each position. Rotor coordinates, material axes and permanent-magnet directions
rotate physically. Uniform, radial, tangential and angular magnetisation laws
follow their material frame, including the law centre. Stator coordinates stay
unchanged. Each new annular strip is triangulated by connecting its ordered
inner and outer vertices. Orientation predicates, edge incidence, coverage,
Euler characteristic and retained constraints are checked before solving.

For |method=band|, the nodal weight $w$ is one on the rotor bulk, zero on the
stator bulk, and $1-\ell/L$ on intermediate ring $\ell$. P1 interpolation makes
its gradient constant in each band triangle. With planar Maxwell stress
\[
 \mathbf T=\frac1{\mu_0}\left(\mathbf B\mathbf B^{\mathsf T}
             -\frac12|\mathbf B|^2\mathbf I\right),\qquad
 \mathbf F=-\int_{\rm band}\mathbf T\nabla w\,\mathrm dV,
\]
the torque is the moment of the same force density about the chosen origin.
The piecewise-constant field and linear moment arm are integrated exactly for
the discrete P1 field. This exact integration does not eliminate finite-element
error. The band's geometric weights avoid an additional Laplace solve and
average stress over its thickness. They are not FEMM's AGE torque formula.

Useful checks include preservation of a uniform field for arbitrary angles,
zero net stress in that uniform field, a magnet in an imposed transverse field,
and the fixed-current identity
$T=\partial W'/\partial\theta$ with $\theta$ in radians. Compare the energy
finite difference at several angular steps; making the step extremely small
can amplify reconnection ripple. Refine bulk spacing, angular segmentation,
radial layers and exterior distance independently. The implementation does not
promise noise-free cogging torque or an error bound from a small solver residual.

\Needspace{10\baselineskip}
\subsection{Persistence, interchange and current limits}
Use separate cache files for the template mesh and each angle/current state.
The template can reuse a mesh-only cache; each position can reuse its solution.
Frozen mode checks the geometry, sources and exact inherited mesh just as it
does for other models. Native format 13 gains a moving-band signature carrying
method identity, template provenance, normalised angle, retained/reconnected
mesh and force weights. Ordinary models keep their previous fingerprints and
stored schema. Template names must be recreated in each TeX run; they are not
serialised as a document-wide object registry.

A supported position exports as an ordinary static |.fem| or |.ans| snapshot
at its current angle. Its material paths, magnet directions and mesh are already
rotated. A FEMM user can solve that one |.fem| normally; the file does not contain
a live moving-band controller or native AGE records. The answer export preserves
the computed P1 mesh and potentials. Existing representability checks still
apply: a rotated off-diagonal material tensor cannot be encoded in FEMM's
global diagonal material records, and nonlinear ANS export needs
|interpolation=femm|. Angular coercive laws retain their documented export rules.
Importing a snapshot creates an ordinary static model, not a sliding template.

The initial scope is one complete circular band in a rectangular native planar
DC problem with |mesher=delaunay| and |field model=open|. Custom domains,
tagged/paired boundaries, periodic sectors, imported templates, callbacks,
parallel/voltage circuits, harmonic excitation, ideal circuits and operating-point
responses are excluded. A native machine node may supply rotor/stator domains
when all its paths satisfy the air-collar contract. Existing TikZ node anchors
are drawing objects and are not recreated or rebound by a later position.

\Needspace{10\baselineskip}
\subsection{Lua entry points and validation}
\begin{codeexample}[code only]
local S = require("luafemm-sliding")
local template = S.capture(model, {
  inner_radius=4.3, outer_radius=5.3,
  center_x=0, center_y=0, segments=96, layers=3
})
local position = S.new(template, {
  angle=30, currents={drive=10},
  cache={file="machine-30",mode="auto"}
})
require("luafemm").solve(position)
local force = require("luafemm-post").measure(position, {
  name="at-30",quantity="force",method="band"
})
\end{codeexample}
|capture| meshes the supplied unmeshed model and returns independently owned
template data; treat those returned data as read-only. |new| returns an independent
meshed, unsolved model. Its angle defaults to zero, currents to an empty override
list and cache mode to |off|. Unknown options are errors. Physical edits to a
position after creation are rejected; create another position from the template.

The portable tests cover bulk identity, angle wrapping, affine fields, analytic
magnet torque, nonlinear coenergy differences, tensor and angular-law rotation,
source overrides and passive caches. The format matrix exercises shared cold,
warm and frozen data in plain LuaTeX, LuaLaTeX and ConTeXt MkIV. The optional
|scripts/sliding-interop.py| exports static angle snapshots and compares torque
and linkage with independent xfemm remeshes and solves. Numerical evidence and
its tolerances are recorded in \nolinkurl{docs/validation.md}.

% SPDX-License-Identifier: LPPL-1.3c
% Copyright (C) 2026 Christophe Jorssen
% Author and Current Maintainer: Christophe Jorssen <christophe.jorssen@gmail.com>
% LPPL maintenance status: maintained
% This file is part of luafemm. See LICENSE and LICENSES.md for its terms.
\section{Drawing eddy-current loops and splitting a conductor}
\label{sec:currents}

\subsection{Which plane contains the current?}
Kamil can now calculate a conducting body's losses. He would also like to show
his students where the induced currents flow. Before adding arrows, he checks
what the drawing plane represents. In an ordinary planar magnetic problem,
$\mathbf B$ lies in the page and $J_z$ points through it. In an axisymmetric
$(r,z)$ section, the current is azimuthal: its circles also leave that page.
Neither picture contains in-plane current streamlines.

The |femm eddy currents| pic chooses a physical view of the current:
\begin{itemize}
\item |view=symbols| draws dots and crosses for currents out of and into the
page. This is the automatic view of planar and axisymmetric models.
\item |view=loops| draws an axisymmetric body's actual circular currents in a
separate transverse $(x,y)$ section, at a specified axial position.
\item |view=lines| traces computed transverse streamlines in the new
|type=axial| diffusion model. Here $B_z$ points through the page, and
$\mathbf J$ circulates within it. This is that model's automatic view.
\end{itemize}
The last model is especially useful for comparing a solid cylinder with an
insulated bundle under the axial field of a long solenoid. It is a different
reduction of Maxwell's equations, not a rotated axisymmetric drawing.
No general three-dimensional streamline tracer is implied.

\subsection{A cylinder, nineteen rods and one exciting coil}
Kamil starts with a copper cylinder of radius $R$ inside a long coil of radius
$R+\varepsilon$. The coil imposes an axial sinusoidal field. He looks along the
cylinder's axis, so its section is a circle. Instead of placing an out-of-page
current source in that circle, he declares the imposed \emph{magnetic} field:
\begin{codeexample}[code only]
\tikzset{induction experiment/.style={
  femm/problem={type=axial,frequency=50,applied field=.01,
    depth=100,xmin=-8,xmax=8,ymin=-8,ymax=8,mesh size=.65}}}
% In a picture with this style:
\filldraw[fill=gray!25,
  femm/region={material=copper,mesh size=.24,curve tolerance=.002}]
  (0,0) circle[radius=6];
\pic[femm/currents={phase=90,lines=6}] {femm eddy currents};
\end{codeexample}
The peak field in the surrounding insulator is $B_0=0.01\,\mathrm T$ and
$f=50\,\mathrm{Hz}$. A depth of $100$ means $100\,\mathrm{mm}$ with the default
unit. It converts per-length energy and loss into totals; it does not introduce
finite-length end effects. The long-solenoid excitation corresponds to
$n\widehat I=B_0/\mu_0$, with $n$ the number of turns per unit length.
The drawn coil circle identifies its radius. It is not a meshed circuit or a
computed finite coil. This approximation prescribes $H_z$ in the connected
insulator and neglects displacement current.

Kamil then replaces the solid section by $N$ circular conductors of radius
$a=R/\sqrt N$. Their total copper area remains
$N\pi a^2=\pi R^2$. The rods are electrically separated by insulation, so each
has its own closed current loops. The illustrated $N=19$ packing has one
central rod, six on the first ring and twelve on the second, as in the supplied
lamination sketch. The example exposes |\Radius|, |\Strands|, |\Gap|,
|\Epsilon|, |\Frequency|, |\AppliedB| and |\Length| as ordinary TeX parameters.
Its explicit packing accepts integer $1\leq N\leq19$; other packings can be
written with ordinary TikZ paths and loops without changing the solver.

There is a geometric catch: the bundle's enclosing radius is \emph{larger}
than $R$. Insulating gaps and packing voids occupy additional area. One cannot
retain both the same metal area and the same radius-$R$ envelope. The example
uses the same coil radius in both pictures and checks that the bundle fits.
For the nineteen-rod packing, its outer ring has radius
\[
 d=2a+g,\qquad \rho_2=\frac{d}{2\sin15^\circ},\qquad
 R_{\rm envelope}=\rho_2+a,
\]
where $g$ is the minimum gap within each ring. The alternating angular offset
keeps the two rings separated. With $R=6\,\mathrm{mm}$, $g=0.12\,\mathrm{mm}$
and $\varepsilon=1.5\,\mathrm{mm}$, both specimens fit in the
$7.5\,\mathrm{mm}$ coil. Increasing $N$ or the gap may require a larger coil.
The computational rectangle must also contain every physical region.

\subsection{Measure the effect instead of counting arrows}
Each picture chooses its contour levels independently. Six current lines per
rod do not mean that the rods carry the same current density as the solid
cylinder! Kamil registers a radial profile and measures the copper losses:
\begin{codeexample}[code only]
% Declare this profile while the solid-cylinder model is current:
\path[femm/profile={name=solid-radius,samples=81}] (0,0)--(5.994,0);
\femmselect[name=solid-metal,material=copper]
\femmmeasure[name=solid-loss,selection=solid-metal]
% Later, in a PGFPlots axis:
\femmplot[profile=solid-radius,component=Jn,abscissa=t,
  representation=rms,plot options={blue,thick}]
\end{codeexample}
On a profile running along the positive $x$ radius, its left normal is the
positive azimuthal direction, so |Jn| is $J_\theta$. The complete example stops
at $0.999a$, just inside the conductor, and rescales the plotted abscissa to
$r/a$ with a PGFPlots |x filter|. This avoids the ambiguous triangle choice
exactly on a metal--air interface, where current density is discontinuous.
The derivative of a P1 field is piecewise constant; steps in a current profile
are a visible discretisation effect. More sample points alone do not remove
them. Start with a coarse mesh when arranging the figure, then reduce local
|mesh size| relative to both the rod radius and the skin depth, and refine
|curve tolerance| before comparing losses.

In the weak-skin limit, a nonmagnetic circular conductor has
\[
 \widehat J_\theta(r)\simeq-\frac{\mathrm i\omega\sigma B_0r}{2},\qquad
 \frac{P(a)}{L}\simeq\frac{\pi\sigma\omega^2B_0^2a^4}{16}.
\]
Consequently $P_N/P_1\simeq1/N$ at equal total metal area. This is a limiting
law, not an imposed numerical scaling: the solver computes each geometry's
shielding and currents. As skin effects grow, the ratio need not be $1/N$.
For the stated copper parameters the computed ratio is about $0.053$, close
to $1/19\simeq0.0526$. The complete source prints its own solved losses and
ratio, and plots the two RMS current profiles on a common scale.
\luafemmexample{eddy-current-bundle}

\subsection{Showing the circles of an axisymmetric current}
A conducting tube inside a finite solenoid is a different useful example.
Keep |type=axisymmetric|, where the finite coil is an actual winding source.
A first pic marks the induced current through the meridional page. A second
pic draws its circular loops at |section=0|, displaced to the right for clarity:
\begin{codeexample}[code only]
\pic[femm/currents={view=symbols,spacing=1.5,phase=90}]
  {femm eddy currents};
\pic[femm/currents={view=loops,section=0,center x=55,
  spacing=.6,phase=90},femm/arrow spacing=10mm]
  {femm eddy currents};
\end{codeexample}
Here $+J_\varphi$ points \emph{into} the $(r,z)$ page, so it is a cross there;
it runs anticlockwise in the transverse $(x,y)$ view. A positive planar $J_z$
is a dot. Both views use the same solved phasors and the same time phase.
The winding is stranded and is not included in the eddy-current glyphs.
\luafemmexample{eddy-current-views}

\subsection{Problem keys and the order of operations}
Declare |femm/problem| first, followed by materials and all physical regions.
Register profiles and selections. Then request a current pic, a measurement or
a plot; each solves on demand. After meshing, do not add or alter physical
regions. Use a new problem for another specimen or frequency. Any number of
current pics may choose independent phases and styles without recomputation.

\begin{key}{/femm/problem/type=axial}
Select the transverse long-conductor $H_z$ formulation. |planar| remains the
default; |axisymmetric| means a meridional $A_\varphi$ model instead.
Axial mode requires positive |frequency| and |mesher=delaunay| (the default
mesher). Arbitrary simply connected solid-conductor cross-sections and several
separated conductors are supported. Their surrounding nonconducting material
must have relative permeability one and remain connected to the exterior.
An insulating cavity enclosed by metal has an unknown coupled axial field and
is rejected. Excluded holes are also rejected; draw insulating regions when
appropriate. Real isotropic linear conductive materials are accepted; nonlinear
laws, coercivity, anisotropy and catalogue loss/lamination/wire models are not.

This mode has passive conductors only: impressed current densities, stranded
regions, named winding circuits, specialised boundaries, ideal circuits and
planar component nodes are unavailable. The scalar conductor-surface condition
is imposed automatically. Standard vector-potential boundary values must stay
at their defaults. Public FEM/ANS interchange does not represent this model;
use its native cache and the TikZ source. Ordinary TikZ curve capture, material
styles and local mesh controls remain available.
\end{key}
\begin{key}{/femm/problem/applied field=\meta{tesla} (initially .01)}
Signed peak axial flux-density amplitude in the surrounding insulator, for
|type=axial|. It sets $\widehat H_0=(B_0/\mu_0)e^{\mathrm i\alpha}$.
Zero is allowed and yields zero field and current. This is a physical source,
not a plotting scale. It has no effect on planar or axisymmetric models.
\end{key}
\begin{key}{/femm/problem/applied phase=\meta{degrees} (initially 0)}
Phase $\alpha$ of the imposed axial field, with the harmonic
$e^{+\mathrm i\omega t}$ convention. Source amplitude and phase enter the
physical cache fingerprint. They are distinct from the current pic's time phase.
\end{key}

\subsection{Current-pic keys and styles}
\begin{key}{/tikz/femm/currents=\marg{options}}
Set the following |/femm/currents| keys. Apply the key to a |tikzpicture| or
scope for shared defaults, or to a pic for local overrides. It is a normal
PGF key family and can be included in reusable |.style| definitions.
\end{key}
\begin{key}{/femm/currents/view=\meta{choice} (initially auto)}
A |.is choice| handler accepting |auto|, |lines|, |symbols| and |loops|.
Automatic selection follows the formulation as explained above. Incompatible
views raise an error; planar scalar currents are never reinterpreted as
in-plane vectors. Only induced currents in solid conductors are drawn.
\end{key}
\begin{key}{/femm/currents/phase=\meta{degrees} (initially 90)}
Instantaneous time phase $\omega t$. A quarter-period default makes weak-skin
currents readily visible under a phase-zero imposed cosine field. It is
independent of |/tikz/femm/phase| and the profile's |phase| (both initially zero).
\end{key}
\begin{key}{/femm/currents/material=\meta{name} (initially empty)}
Restrict the view to a declared material. Empty means all conductive materials;
unknown names raise an error. Use distinct material names when specimens need
independent selection even though their properties coincide.
\end{key}
\begin{key}{/femm/currents/lines=\meta{integer} (initially 7)}
For |view=lines|, choose 1 to 200 equally spaced interior streamfunction levels
between the selected conductors' extrema. A level may produce several closed
paths, one in each rod. This does not specify a line count per connected body
or an absolute current increment. Zero current produces no paths.
\end{key}
\begin{key}{/femm/currents/spacing=\meta{model units} (initially 2)}
Positive target spacing of symbol centres or sampled loop radii. The bounds
are divided into equally sized cells no larger than this spacing, with samples
at their centres. Only centres in conducting material are drawn. Thin regions
can be missed by a coarse grid; reduce spacing to resolve them. This controls
the display only, not the field mesh. Symbol grids above 10000 cells and radial
loop grids above 1000 cells are rejected. |view=lines| ignores this key.
\end{key}
\begin{key}{/femm/currents/section=\meta{model units} (initially 0)}
Axial coordinate $z$ of a transverse |view=loops| section, within the model's
vertical bounds. Each radius is sampled in the original meridional solution.
\end{key}
\begin{key}{/femm/currents/center x=\meta{model units} (initially 0)}
\end{key}
\begin{key}{/femm/currents/center y=\meta{model units} (initially 0)}
Drawing centre of the transverse loop view, expressed in the problem's root
coordinate frame. These keys move its illustration, not the physical solution
or the sampled section. Other views ignore them. Like field pics, current pics
use the root problem transform; a local pic shift does not move solved geometry.
\end{key}
\begin{key}{/femm/currents/reference current=\meta{amperes per square metre} (initially 0)}
For symbols and loops, zero chooses the largest sampled absolute instantaneous
current density. A positive value sets an explicit common opacity reference
for comparison across pics. Relative strength $q=\min(1,|J|/J_{\rm ref})$
gives opacity $0.2+0.8q$. Sign controls direction; opacity encodes magnitude.
The key does not affect |view=lines|, whose paths have unit opacity.
\end{key}
\begin{key}{/femm/currents/minimum fraction=\meta{number} (initially .02)}
Suppress symbols and loops with $q$ at or below this threshold. The accepted
range is $0\leq q_{\min}<1$. This avoids drawing numerical near-zero currents
with arbitrary direction. |view=lines| ignores the threshold.
\end{key}
\begin{key}{/femm/currents/symbol radius=\meta{TeX dimension} (initially 1.3pt)}
Radius of each dot/cross circle, in physical page units. Model scaling affects
its position, not this radius. Keep the dimension positive.
\end{key}
\begin{key}{/tikz/femm/eddy style=\marg{TikZ options}}
Append local drawing options after the base |femm eddy line| or
|femm eddy symbol| style. The defaults use a fine dark-red stroke and the
existing oriented-field arrows. |femm/arrows|, |femm/arrow spacing|,
|femm/arrow length| and |femm/arrow width| also apply to current paths.
Set colours and line patterns with ordinary TikZ styles:
\begin{codeexample}[code only]
\tikzset{induced loops/.style={femm/currents={phase=90,lines=8},
  femm/eddy style={blue!65!black,line width=.4pt},femm/arrow spacing=8mm}}
\pic[induced loops] {femm eddy currents};
\end{codeexample}
\end{key}

\subsection{Profiles, losses and the discrete formulation}
Axial profiles add |Bz| and |Hz|, and the transverse current components
|Jx|, |Jy|, |Jt| and |Jn|. The last two project onto the oriented profile's
unit tangent and left normal. Current units are $\mathrm{A/m^2}$.
|J| and |Jeddy| coincide and mean the vector norm in axial mode; they remain
signed out-of-page scalars in planar/axisymmetric mode. A vector norm has no
single phase: request |phase| on a signed component instead. All usual harmonic
representations work. In-plane magnetic components are zero; |B| and |H|
are the norms of the axial field. Vector-potential and meridional-flux
components are unavailable. The |femm field| pic is rejected because its
in-plane magnetic contours would have no physical meaning in this model.

Regional measurements supply |area|, |volume|, |energy|, |coenergy|,
|joule_loss| and |eddy_loss|. The two losses coincide; energy and coenergy
are equal cycle averages. Supply positive |depth| for totals. Current integrals
and forces are unavailable in this reduction. Losses use
$\int(|\widehat J_x|^2+|\widehat J_y|^2)/(2\sigma)\,dV$ and need no graphical
contour or glyph sampling.

In a conductive domain, the weak form is
\[
 \int\frac1\sigma\nabla\widehat H_z\mathbin{\cdot}\nabla v\,dS
 +\mathrm i\omega\int\mu\widehat H_zv\,dS=0,
 \qquad \widehat H_z|_{\partial\Omega}=\widehat H_0,
 \qquad \widehat{\mathbf J}=(\partial_y\widehat H_z,-\partial_x\widehat H_z).
\]
The implementation stores the dimensionless P1 perturbation
$u=\widehat H_z/\widehat H_0-1$ and solves its unit-drive equation even when
$B_0=0$, then multiplies physical fields by the requested source. The same
complex ILU(0)--GMRES solver and consistent mass used for harmonic magnetics
apply. Internally named |A| arrays contain this perturbation in axial mode;
they must never be exported or interpreted as vector potential. Use
|luafemm-harmonic.sample| for phasors and the public profile/measurement keys.
Cache format 13 distinguishes the formulation, source and both scalar arrays.

At a fixed phase, the $H_z$ contours are current streamlines, with orientation
from $(\partial_yH_z,-\partial_xH_z)$. Constant surface $H_z$ makes the normal
current vanish on conductor boundaries. P1 interpolation yields continuous
streamfunction contours and piecewise constant currents. The cylinder benchmark
compares against
\[
 \widehat H_z(r)=\widehat H_0\frac{I_0(kr)}{I_0(ka)},\qquad
 \widehat J_\theta(r)=-\widehat H_0 k\frac{I_1(kr)}{I_0(ka)},\qquad
 k=\frac{1+\mathrm i}{\delta}.
\]
The tests evaluate the modified Bessel functions independently by power series,
check mesh convergence and Joule losses, then test equal-area insulated disks,
closed loop orientation, cache restoration and all three generic TeX formats.
They do not validate finite-cylinder end effects or inter-strand contact.

% SPDX-License-Identifier: LPPL-1.3c
% Copyright (C) 2026 Christophe Jorssen
% Author and Current Maintainer: Christophe Jorssen <christophe.jorssen@gmail.com>
% LPPL maintenance status: maintained
% This file is part of luafemm. See LICENSE and LICENSES.md for its terms.
\section{Parametric components as nodes}
\label{sec:components}
A component is a genuine PGF node with a fixed physical geometry and named
anchors. One node can declare several material regions. Its text and labels
are annotations: their font, width and padding never resize its magnetic
parts. Additional E and toroidal shapes, independent section widths and ideal-mode
mean contours are described in Section~\ref{sec:ideal}.
Components use the same region, mesh, solver and cache infrastructure
as explicitly drawn paths; the chosen field model determines their computational boundary conditions.

\begin{key}{/tikz/femm/u core=\marg{options} (default \normalfont empty)}
Select a U-shaped core open towards local $+y$. The node registers one
simple closed iron contour. It has no armature, winding or gap refinement
regions. The default dimensions match the core in the opening tutorial.
\end{key}
\begin{key}{/tikz/femm/u electromagnet=\marg{options} (default \normalfont empty)}
Select a U core, an armature above its two poles, and two winding sections
around the left leg. A single node registers four material regions and,
by default, two air regions for gap refinement.
The armature has the same width as the core; both gaps have the same height.
\end{key}

Put either key on a node inside a picture with |femm/problem|. It selects
the internal component shape and enables |transform shape|. Declare every
component and additional physical path before requesting a mesh or a field.
Components do not solve during construction. They may be named or anonymous.
Ordinary PGF grouping applies; a reusable TikZ style can set their options.

\begin{codeexample}[width=6cm]
\begin{tikzpicture}[femm/problem={
  xmin=-60,xmax=60,ymin=-35,ymax=35},
  scale=.5]
  \node[femm/u core={
    width=30,height=25,
    leg thickness=5}] (A) at (-20,0) {};
  \node[femm/u core={
    width=30,height=25,
    leg thickness=5,
    core style={fill=blue!15}},
    rotate=15] (B) at (20,0) {};
  \draw[red,thick]
    (A.left pole)--(A.right pole);
  \draw[blue,thick]
    (B.left pole)--(B.right pole);
\end{tikzpicture}
\end{codeexample}

\subsection{Geometry, materials and excitation}
The following subkeys belong to |/femm/component| and are supplied inside
either component key. All lengths are numeric literals in model millimetres,
not TeX dimensions; no |mm| suffix is needed. Values are finite numbers,
not PGF expressions. Defaults are inherited within the usual TeX scope.
Electromagnet-only options have no effect on |femm/u core|.

\begin{key}{/femm/component/width=\meta{number} (initially 80)}
Outer width of the U core and armature. It must exceed twice the leg thickness.
The winding extends outside this width on the left.
\end{key}
\begin{key}{/femm/component/height=\meta{number} (initially 60)}
Outer height of the U core alone, from its underside to its pole faces.
It must exceed the leg thickness.
\end{key}
\begin{key}{/femm/component/leg thickness=\meta{positive number} (initially 15)}
Common thickness of both legs and the bottom yoke. The window has width
|width-2*leg thickness| and height |height-leg thickness|.
\end{key}
\begin{key}{/femm/component/gap=\meta{positive number} (initially 3)}
Height of each air gap between a pole and the armature; electromagnet only.
Zero is rejected in this initial component interface.
\end{key}
\begin{key}{/femm/component/armature thickness=\meta{positive number} (initially 15)}
Height of the separate armature; electromagnet only. This is an independent
parameter: changing the leg thickness does not change this default.
\end{key}
\begin{key}{/femm/component/coil width=\meta{positive number} (initially 6)}
Width of each rectangular winding section; electromagnet only.
\end{key}
\begin{key}{/femm/component/coil height=\meta{positive number} (initially 26)}
Height of each winding section; electromagnet only. Both sections must fit
between the top of the yoke and the pole faces.
\end{key}
\begin{key}{/femm/component/coil clearance=\meta{nonnegative number} (initially 2)}
Distance from each winding section to its adjacent face of the left leg.
The inner section must fit inside the U window; electromagnet only.
\end{key}
\begin{key}{/femm/component/coil center y=\meta{number} (initially 3)}
Vertical position of the two section centres relative to |core center|,
positive towards the poles. It must keep the entire winding above the yoke
and below the pole faces; electromagnet only.
\end{key}
\begin{key}{/femm/component/ampere turns=\meta{number} (initially 6000)}
Signed $NI$ in A\,turns, for the complete winding. A positive value produces
positive $J_z$ in the outer (left) section and negative $J_z$ in the inner
section. The integral of $J_z$ over each section is respectively $NI$ and
$-NI$; their net current is zero. Negative values reverse the excitation;
zero leaves a passive winding.

The package computes $J_z=NI/S$ using the area of the actual captured
polygon in square metres. Scaling or deforming the node therefore preserves
ampere-turns and adjusts current density. This convention differs from a raw
|femm/current density| on a path, which holds density fixed.
The winding is homogenised, not drawn or solved turn by turn.
\end{key}
\begin{key}{/femm/component/core material=\meta{material} (initially pure-iron)}
Material of the U core, resolved through the existing region material
selector. Catalogue identifiers and custom |femm/materials/| styles are accepted.
\end{key}
\begin{key}{/femm/component/armature material=\meta{material} (initially pure-iron)}
Material of the armature; electromagnet only.
\end{key}
\begin{key}{/femm/component/coil material=\meta{material} (initially copper)}
Material of both winding sections; electromagnet only. Excitation is determined
by |ampere turns|, not by a current density stored in the material catalogue.
The core, armature and gap regions have zero imposed current density.
\end{key}
\begin{key}{/femm/component/magnetization angle=\meta{number} (initially 0)}
Direction of any coercive material in degrees from the component's local
$+x$ axis. It applies to material parts and has no effect on passive materials.
The direction is transformed into the model frame when the node is placed;
rotating the component rotates its prescribed magnetisation. This differs
from |femm/region/magnetization angle|, which is already a model-frame angle.
\end{key}
\begin{key}{/femm/component/mesh size=\meta{nonnegative number} (initially 0)}
Local spacing for the core, armature and winding. Zero inherits the problem
spacing. As with ordinary regions, positive values refine rather than coarsen.
Mesh lengths are model-frame values; node scaling does not scale these keys.
\end{key}
\begin{key}{/femm/component/gap mesh size=\meta{nonnegative number} (initially 1)}
Local spacing in both air gaps; electromagnet only. A positive value adds
two air refinement rectangles. In open mode zero omits these two declarations and uses
the ordinary global mesh in the gaps. In ideal mode the two gaps remain
part of the confined domain, using the global mesh spacing.
\end{key}

\subsection{Anchors and placement}
The |center| anchor is the centre of the rectangular envelope of all parts,
including the exterior winding. In an electromagnet it differs from
|core center|. Use |anchor=core center| to place the U at the same coordinates
as the tutorial's explicit polygon. Text is centred at |window center|
independently of the anchor used for placement.

\begin{center}
\begin{tabular}{p{.42\linewidth}p{.49\linewidth}}
\toprule Anchor & Meaning in local component coordinates \\
\midrule
|center| & Centre of the component envelope \\
|north|, |south|, |east|, |west| & Midpoints of envelope sides \\
|north east|, |north west|,\newline |south east|, |south west|
  & Corners of the envelope \\
|core center| & Centre of the U's outer rectangle \\
|window center| & Centre of the open window; text location \\
|yoke center| & Centre of the bottom iron bar \\
|left pole|, |right pole| & Centres of the two pole faces \\
|left pole inner|, |left pole outer| & Inner/outer endpoints of the left pole \\
|right pole inner|, |right pole outer| & Inner/outer endpoints of the right pole \\
|left armature|, |right armature| & Points facing the pole centres across the gaps \\
|left gap|, |right gap| & Centres of the gaps \\
|armature center| & Centre of the armature rectangle \\
|coil positive|, |coil negative| & Centres of the outer/inner winding sections \\
\bottomrule
\end{tabular}
\end{center}
The last four rows are available only on an electromagnet; requesting one on
a bare core is an error. The coil anchor names describe the sign for positive
|ampere turns|; their geometric locations do not swap when excitation reverses.
An automatic connection such as |(M)--(N)| meets the rectangular envelope,
which may enclose air. Use physical anchors for probes or exact pole contacts.

\begin{codeexample}[code only]
\node[femm/u electromagnet={gap=3,ampere turns=6000},
  anchor=left pole,rotate=20] (M) at (0,0) {};
\draw[femm/profile={name=left-gap,samples=101},red,dashed]
  (M.left pole)--(M.left armature);
\end{codeexample}

Nodes support |at|, |anchor|, rotations and local affine transformations,
including inherited transformations because |transform shape| is enabled.
Keep that option enabled. Whole-picture transforms remain display transforms,
subject to the capture precision described in Section~\ref{sec:model}.
Changing a physical transform or component dimension invalidates the cache.
Changing text or painting styles does not. Very small rounding differences
under display transformations may conservatively cause a cache miss.

Dimensions use the original picture coordinate basis. Changing local |x|
or |y| vectors does not resize an existing shape: use component dimensions or
node transformations. |minimum width|, |minimum height|, |inner sep| and
|outer sep| do not alter its fixed physical outline or anchors.
Positioning-library operations such as |right=10mm of M| use the envelope.
Delayed placement, including matrix-managed nodes, is rejected: physical
regions must be captured at their final position, once.

\subsection{Drawing styles and a complete example}
\begin{keylist}[/femm/component]{core style=\marg{TikZ options},
  armature style=\marg{TikZ options},coil style=\marg{TikZ options},
  gap style=\marg{TikZ options}}
Append painting options to the corresponding part style within the current
scope; no extra options are appended initially. Use colours, line widths,
opacity or dash patterns here. Geometry-changing keys, path decorations
that replace contours, and physical keys inside these styles are unsupported:
use the explicit component parameters for physical changes.
\end{keylist}
\begin{stylekey}{/tikz/femm component core}
Base core style: |draw=black!70,fill=gray!15|.
\end{stylekey}
\begin{stylekey}{/tikz/femm component armature}
Base armature style: |femm component core|. Thus it inherits core-style changes
unless explicitly overridden.
\end{stylekey}
\begin{stylekey}{/tikz/femm component coil}
Base winding style: |draw=orange!70!black,fill=orange!25|.
\end{stylekey}
\begin{stylekey}{/tikz/femm component gap}
Empty initially: refinement rectangles are invisible. A |gap style| has
no effect when |gap mesh size=0| because no gap paths are emitted.
\end{stylekey}
These styles paint the individual parts. Node-level |draw| or |fill| does not
replace their styles. Ordinary node text options and |label| remain available.
Each component is registered once, in declaration order: later physical
regions may still override it according to the ordinary overlap rule.

The file \nolinkurl{examples/u-electromagnet-node.tex} builds the tutorial
electromagnet with one node, locates its measurement path with gap anchors,
and plots $|\mathbf B|$ and $B_y$. It enables the persistent cache and starts
with a coarse mesh. The core is a single concave polygon; all other regions
are separate simple polygons. Compound material paths are available
independently (Section~\ref{sec:domains}); these node templates do not need
them. Rotor and stator templates are described in Section~\ref{sec:machines}.

The complete source below uses node anchors for the gap scan and labels.
Changing a component dimension moves these locations with the physical parts.
The result includes the field and the PGFPlots scan produced by that source.
\luafemmexample{u-electromagnet-node}

% SPDX-License-Identifier: LPPL-1.3c
% Copyright (C) 2026 Christophe Jorssen
% Author and Current Maintainer: Christophe Jorssen <christophe.jorssen@gmail.com>
% LPPL maintenance status: maintained
% This file is part of luafemm. See LICENSE and LICENSES.md for its terms.
\section{Rotors, stators and windings}
\label{sec:machines}

The machine nodes are genuine PGF shapes, using the same capture, materials,
mesh, solver and persistent-cache machinery as ordinary physical paths.
They describe a constant-depth planar cross-section with axial current density.
They require |field model=open|. Declaring a machine does not start a solve.
A rotor may precede or follow its stator: compound core contours preserve the
bore rather than painting an air disk over earlier regions.

\begin{keylist}[/tikz]{femm/rotor=\marg{options},femm/stator=\marg{options}}
Select a machine component and apply its options in |/femm/machine|.
The default argument is empty. Both keys enable |transform shape| and fix
node padding to zero. Physical dimensions are independent of node text.
Standard placement and transformation conventions are those of
Section~\ref{sec:components}. A circular gap scan additionally requires
concentric components with circular, unsheared transformed surfaces.
\end{keylist}

\subsection{Geometry and materials}
All lengths below are finite numerical literals in model millimetres; angles
are mechanical degrees, positive counterclockwise from local $+x$.
As in the existing component interface, dimension values are not unevaluated
PGF mathematical expressions. Evaluate expressions before supplying the keys.
Mesh sizes are measured in the model frame and do not scale with a node.

\begin{key}{/femm/machine/radius=\meta{positive number} (initially 29)}
Nominal rotor air-gap radius, before cutting its slots. Rotor only.
\end{key}
\begin{key}{/femm/machine/bore radius=\meta{positive number} (initially 30)}
Nominal stator bore radius, before cutting its slots. Stator only.
\end{key}
\begin{key}{/femm/machine/outer shape=\meta{circle or square} (initially circle)}
PGF choice selecting the stator core's outer outline. The bore remains circular.
\end{key}
\begin{keylist}[/femm/machine]{outer radius=\meta{number} (initially 45),
  outer width=\meta{number} (initially 90)}
Outer core radius for a circular stator, or full side length for a square
stator. The inactive dimension has no effect. The core must contain its bore,
slots and the requested minimum yoke.
\end{keylist}
\begin{key}{/femm/machine/shaft radius=\meta{nonnegative number} (initially 0)}
Rotor shaft radius. Zero omits the shaft region. A positive value cuts a
circular hole from the rotor core and fills it with |shaft material|.
\end{key}
\begin{key}{/femm/machine/housing thickness=\meta{nonnegative number} (initially 0)}
Separate stator casing outside the core, circular or square as selected above.
Zero omits it. It has its own material and drawing style.
\end{key}
\begin{key}{/femm/machine/minimum yoke=\meta{positive number} (initially 1)}
Minimum geometric clearance from slots to the shaft or outer core surface.
This is a construction check, not a saturation criterion. The angular-envelope
check also rejects overlapping adjacent slots conservatively; deeply staggered
slots with overlapping angular envelopes are not admitted.
\end{key}
\begin{keylist}[/femm/machine]{core material=\meta{material} (initially pure-iron),
 shaft material=\meta{material} (initially air),
 housing material=\meta{material} (initially air)}
Select catalogue identifiers or custom |femm/materials/| styles. The shaft and
housing have zero prescribed current. Conductors have independent selectors.
The default air housing is physically passive even when its style paints it.
\end{keylist}
\begin{key}{/femm/machine/magnetization angle=\meta{number} (initially 0)}
Local coercive-field direction for any coercive material used in a part.
It is transformed with the component. No automatic radial magnet segmentation
is implied by selecting a permanent-magnet material.
\end{key}
\begin{keylist}[/femm/machine]{mesh size=\meta{nonnegative number} (initially 0),
 slot mesh size=\meta{nonnegative number} (initially 0),
 conductor mesh size=\meta{nonnegative number} (initially 0)}
Local refinement for core/shaft/housing, slot cavities and conductors,
respectively. Zero inherits the problem spacing. Slot records can override
the latter two values. Refinement remains bounded by the ordinary mesh budget.
\end{keylist}
\begin{key}{/femm/machine/curve tolerance=\meta{positive number} (initially .03)}
Maximum chord sagitta used to discretise circular interfaces and notch fillets.
The arc step is also capped at ten degrees (fifteen for small fillets).
All machine material paths are emitted as polygons and are exported to FEMM
as segments. The nominal circles, their polygonal approximation and the mesh
are distinct: refine both tolerance and spacing in a convergence study.
\end{key}

\subsection{Slot placement and shape}
\begin{key}{/femm/machine/slots=\marg{options}}
Apply shared slot parameters in |/femm/machine/slots|. Its argument has no
implicit reset: PGF styles and grouping work normally. The defaults below
produce no slots until a count or explicit position list is provided.
\end{key}
\begin{key}{/femm/machine/slots/placement=\meta{uniform or explicit} (initially uniform)}
Choice between a generated angular sequence and the supplied |angles| list.
\end{key}
\begin{key}{/femm/machine/slots/count=\meta{integer} (initially 0)}
Number of slots for uniform placement, from 0 to 512. Zero makes a smooth
component. Explicit placement uses the number of entries in |angles| instead.
\end{key}
\begin{keylist}[/femm/machine/slots]{start angle=\meta{number} (initially 0),
 pitch=\meta{number} (initially 0)}
The $i$th uniform slot is at |start angle| plus $(i-1)$ times the pitch.
Zero pitch means $360/\texttt{count}$. Nonzero pitch allows a partial angular
array; angles are reduced modulo 360 for geometry. Coincident slots fail.
\end{keylist}
\begin{key}{/femm/machine/slots/angles=\marg{comma-separated numbers} (initially empty)}
Mechanical slot-centre angles for explicit placement, up to 512 entries.
Indices follow list order, even when the list is not sorted angularly.
\end{key}
\begin{key}{/femm/machine/slots/offset=\meta{number} (initially 0)}
Additional angular displacement, especially useful in an individual slot
record. It changes the geometry and the automatic winding's angular weight.
\end{key}
\begin{key}{/femm/machine/slots/shape=\meta{choice} (initially trapezoid)}
Choices are |rectangle|, |trapezoid| and |radial|. These describe the body
behind the narrow mouth. Rectangle has equal body widths; trapezoid uses the
independent bottom width; radial makes the body's two flanks pass through the
machine centre. The mouth itself has parallel sides.
\end{key}
\begin{keylist}[/femm/machine/slots]{depth=\meta{positive number} (initially 8),
 neck depth=\meta{positive number} (initially 1)}
Depth to the bottom, and depth to the shoulder where the mouth meets the
body. Both are measured along the slot's centre ray from the nominal surface,
into the iron: outward for the stator and inward for the rotor.
The total depth must exceed the neck depth.
\end{keylist}
\begin{keylist}[/femm/machine/slots]{opening width=\meta{positive number} (initially 2),
 opening angle=\meta{number} (initially 0)}
Width is the chord between the two mouth endpoints on the nominal circle;
it must be smaller than the nominal radius.
A positive opening angle overrides that width and specifies their subtended
angle; it must be less than 90 degrees. Zero selects the width convention.
This explicit precedence also applies inside styles and slot overrides.
\end{keylist}
\begin{keylist}[/femm/machine/slots]{width=\meta{positive number} (initially 4),
 bottom width=\meta{positive number} (initially 5)}
Tangential body width at the shoulder and bottom. The shoulder width must
contain the mouth. |bottom width| applies only to trapezoids; rectangles use
|width| and radial slots derive the bottom width from the two radii.
\end{keylist}
\begin{key}{/femm/machine/slots/corner radius=\meta{nonnegative number} (initially 0)}
Circular fillets on internal notch corners; zero leaves sharp corners.
The radius must fit adjacent edges and must not exceed conductor clearance.
Core and cavity share the same fillet vertices. The surface-mouth endpoints
remain on the nominal circle.
\end{key}
\begin{key}{/femm/machine/slots/material=\meta{material} (initially air)}
Material filling the entire cut slot before its conductors are inserted.
Use a custom nonmagnetic material to represent insulation. This region carries
zero imposed current. An insulating liner and residual air can be represented
by their common magnetic permeability in this magnetostatic model.
\end{key}
\begin{keylist}[/femm/machine/slots]{mesh size=\meta{nonnegative number} (initially 0),
 conductor mesh size=\meta{nonnegative number} (initially 0)}
Positive values override the machine's respective slot/conductor mesh keys.
Zero uses those machine values, which may themselves inherit the problem.
\end{keylist}
\begin{key}{/femm/machine/slot=\marg{options}}
Append an individual slot override. Supply |index|, a one-based slot number,
and any shared slot keys that describe its geometry, filling or conductors.
Unspecified fields retain the shared values. Placement, count, start angle, pitch and angles are global
array settings and are rejected in an individual slot record.
A slot index may appear only once. All records are compiled when the node is
constructed, after shared parameters are resolved, independently of the order
in which |slots| and |slot| appeared in that node's options.
\end{key}
\begin{key}{/femm/machine/slot/index=\meta{integer}}
Required existing slot index. Slot override fields live in |/femm/machine/slot|
and start empty; their names and meanings match the shared slot keys.
\end{key}

\subsection{Conductor sections and excitation}
\begin{key}{/femm/machine/slots/conductor material=\meta{material} (initially copper)}
Material of every conductor layer in the slot. This is independent of the
imposed current: changing copper to aluminium does not change a prescribed
$NI$. The present solver does not compute conductivity-driven eddy currents.
\end{key}
\begin{key}{/femm/machine/slots/clearance=\meta{positive number} (initially .4)}
Minimum inset from the straight body walls. Conductors occupy only the body,
not the mouth. The sloping side-wall inset is measured perpendicular to the
wall; the bottom and shoulder insets are radial. An excessive inset is an error.
\end{key}
\begin{keylist}[/femm/machine/slots]{layers=\meta{integer} (initially 1),
 layer gap=\meta{nonnegative number} (initially .3)}
One to eight equal-depth conductor layers, ordered from the shoulder towards
the bottom, with the stated radial separation. Each layer retains the local
trapezoidal width and has its own area. Unassigned layers are passive.
\end{keylist}
\begin{key}{/femm/machine/slots/current density=\meta{number} (initially 0)}
Uniform prescribed $J_z$ in A/m$^2$ in every conductor layer. Positive means
out of the page. A nonzero value conflicts with a coil or automatic winding
on the same layer. Raw density stays fixed under node scaling, unlike $NI$.
\end{key}
\begin{key}{/femm/machine/conductor=\marg{options}}
Append a conductor override, applied after slot geometry and before winding
assignment. Subkeys belong to |/femm/machine/conductor|: |slot=1|, |layer=1|,
|material|, |current density| and |style|. The last three start empty and leave
the corresponding slot setting unchanged. Repeated records apply in order.
Geometry is controlled by the slot's widths, depth, clearance and layer keys.
\end{key}
\begin{key}{/femm/machine/coil=\marg{options}}
Append a balanced two-sided coil. Subkeys are in |/femm/machine/coil|.
Multiple coils add their signed ampere-turns if they share a conductor layer.
The records do not draw end connections: the model is a transverse section.
\end{key}
\begin{key}{/femm/machine/coil/slots=\marg{first,second}}
Required two distinct existing slot indices. The first receives $+NI$ and
the second $-NI$. Their angular separation is unrestricted by the winding
interface; geometric collision checks still apply.
\end{key}
\begin{key}{/femm/machine/coil/layer=\meta{integer} (initially 1)}
Conductor layer used in both slots. It must exist in both.
\end{key}
\begin{keylist}[/femm/machine/coil]{turns=\meta{nonnegative number},
 current=\meta{number},ampere turns=\meta{number}}
Choose the product |turns*current| or the direct |ampere turns| value.
Defaults for the product are one turn and zero current. Specifying a direct
value together with either product key is rejected. Negative current or
ampere-turns reverses the coil. Fractional homogenised turns are permitted.
Actual polygon area determines $J_z=NI/S$ after physical transformations.
\end{keylist}
\begin{key}{/femm/machine/coil/material=\meta{material} (initially empty)}
Optional material override for both assigned conductor layers. Empty retains
their individual materials. If several coils override a shared layer's
material, the last such record wins.
\end{key}
\begin{key}{/femm/machine/winding=\marg{options}}
Configure an automatic distributed winding. Its subkeys are listed below.
Automatic distribution and explicit |coil| records cannot be combined.
\end{key}
\begin{key}{/femm/machine/winding/distribution=\meta{choice} (initially none)}
|none| leaves excitation to coil records or raw densities. |uniform| uses
weights $\operatorname{sign}\sin(p(\theta_i-\theta_0))$; |sine| uses
$\sin(p(\theta_i-\theta_0))$. Weights at numerical zeros are zero. Positions
must give a balanced sum. See Section~\ref{sec:tutorial-machines} for the
normalisation and its physical interpretation.
\end{key}
\begin{key}{/femm/machine/winding/ampere turns=\meta{number} (initially 1000)}
Total excitation assigned to the positive-weight slots, distributed across
all their layers. Negative values reverse the entire winding. This is not an
excitation per slot. The negative-weight slots carry the opposite total.
\end{key}
\begin{keylist}[/femm/machine/winding]{pole pairs=\meta{positive integer} (initially 1),
 axis=\meta{number} (initially 0)}
Spatial order $p$ and mechanical angular origin $\theta_0$ of the current
weights. Rotation of the whole node rotates its slots and sources together.
\end{keylist}

\subsection{Styles, anchors and a complete construction}
\begin{keylist}[/femm/machine]{core style=\marg{options},shaft style=\marg{options},
 housing style=\marg{options},slot style=\marg{options},conductor style=\marg{options}}
Append ordinary TikZ painting options to the respective part style. There are
no appended options initially. Physical changes belong in physical keys;
path replacements, transformations and physical-region keys in these painting
styles are unsupported.
\end{keylist}
\begin{keylist}[/femm/machine/slots]{style=\marg{options},conductor style=\marg{options}}
Extra painting options for a slot cavity or its conductor layers, empty
initially. Individual slot records can override them; |conductor| records
can override one layer's style. The global part style is applied first.
\end{keylist}
\begin{stylekey}{/tikz/femm machine core}
Initially |draw=black!70,fill=gray!15|.
\end{stylekey}
\begin{stylekey}{/tikz/femm machine conductor}
Initially |draw=orange!70!black,fill=orange!25|.
\end{stylekey}
\begin{stylekey}{/tikz/femm machine slot}
Initially empty; cavity material remains present even when invisible.
\end{stylekey}
\begin{stylekey}{/tikz/femm machine shaft}
Initially |draw=black!50|.
\end{stylekey}
\begin{stylekey}{/tikz/femm machine housing}
Initially |draw=black!70,fill=gray!30|.
\end{stylekey}

The standard compass anchors refer to the bounding box. |center|,
|core center| and |window center| coincide at the machine axis. The latter
also determines text placement. Physical anchors |gap east|, |gap north|,
|gap west| and |gap south| lie on the nominal air-gap circle. For each slot
$i$, the node supplies |slot i| (body centre), |slot i opening| (nominal surface),
|slot i bottom| and |slot i conductor j| (layer centre). Indices are literal
integers in the anchor name, for example |(S.slot 3 conductor 2)|.
An unavailable index raises an error when queried.

\luafemmexample{machine-nodes}

\subsection{Refinement and circular profiles between nodes}
\begin{key}{/tikz/femm/air gap=\marg{options}}
Set |/femm/air gap| options: required |rotor| and |stator| node names, and
|mesh size| (initially .5 mm, strictly positive). The |femm air gap| pic adds
an air refinement annulus after both nodes exist and before any solve.
It checks concentricity, circular transformed surfaces and positive clearance.
The annulus stays two curve tolerances from each nominal surface so that
polygonal approximation cannot paint air over iron. It refines the gap's
interior, not a prescribed number of radial element layers.
\end{key}
\begin{stylekey}{/tikz/femm air gap}
Initially empty. Paint the refinement band by redefining this style if desired.
\end{stylekey}
\begin{key}{/tikz/femm/gap profile=\marg{options}}
Set |/femm/gap profile| options for the |femm gap profile| pic:
|rotor| and |stator| are required names; |name=gap| names the registered
profile; |position=.5| is its fractional radial location from rotor to stator;
|samples=361| includes the closing endpoint; |angle origin=0| selects its
starting model-frame angle; |pole pairs=1| scales the electrical abscissa.
The position must lie strictly inside the gap, with clearance from the
polygonal walls. The pic registers an exact circular scan without solving.
Its visible annotation uses short straight segments. Its drawing transform is
the source model's frame; additional pic placement does not move the scan.
\end{key}
\begin{stylekey}{/tikz/femm gap profile}
Initially |draw=red!70!black,dashed,thin|. An empty drawing style can hide the
annotation without removing the profile.
\end{stylekey}

The gap must exceed four transformed curve tolerances. Uniform scaling and
rotation preserve circularity; unequal scaling and shear are rejected by the
gap helpers, though the individual nodes can still be transformed that way.
A gap-profile centre is retained with its source model, so subsequent figures
do not change its interpretation. For circulation, use an ordinary captured path rather than this exact circular
scan. Direct circular profiles are also available
as |require('luafemm-profiles').circle(model,name,x,y,radius,options)| in Lua.

\subsection{Polar plots and spatial harmonics}
\begin{keylist}[/femm/profile]{center x=\meta{number} (initially 0),
 center y=\meta{number} (initially 0),angle origin=\meta{number} (initially 0),
 pole pairs=\meta{positive integer} (initially 1)}
Polar reference for an ordinary captured profile. The centre is specified in
model coordinates. Components |Br| and |Btheta| use the outward radial and
counterclockwise angular unit vectors; |Hr| and |Htheta| follow the same rule.
At the centre, polar components are undefined and plotted as missing values.
Cartesian components retain their existing meaning.
\end{keylist}
The plot |component| choice accepts these four additional values. Its
|abscissa| choice additionally accepts |angle| and |electrical angle|.
Both are in degrees, relative to |angle origin|; the electrical angle is
|pole pairs| times the mechanical angle. Ordinary profiles unwrap successive
angles by the shortest increment, so sample densely enough to resolve the
winding around the centre. Exact circular profiles run from 0 to 360 degrees.

\begin{command}{\femmharmonicplot\oarg{options}}
Inside a PGFPlots axis, emit an ordinary |\addplot+| of harmonic amplitudes
against integer mechanical spatial order. No trailing semicolon is needed.
Requires a full circular gap profile (or Lua circular profile), with finite
samples. The repeated closing endpoint is omitted from Fourier sums.
\end{command}
\begin{command}{\femmharmonic\oarg{options}}
Print one spectrum statistic using PGF's number formatter. Both commands
parse their options locally in |/femm/harmonic| and can share family styles.
The calculation is pure Lua and reuses the profile's cached field samples.
\end{command}
\begin{key}{/femm/harmonic/profile=\meta{identifier}}
Required circular profile name; initially empty.
\end{key}
\begin{key}{/femm/harmonic/component=\meta{component} (initially Br)}
Field component to analyse: B, H, Az, their Cartesian, path-relative or polar
components. Circulation is not a periodic field component and is rejected.
\end{key}
\begin{key}{/femm/harmonic/maximum order=\meta{positive integer} (initially 15)}
Largest mechanical spatial order included. It must be strictly below half the
number of nonduplicate angular samples. This is a sampling limit, not a claim
that the FEM mesh resolves that harmonic. Increase both mesh resolution and
sampling and compare results when high-order amplitudes matter.
\end{key}
\begin{key}{/femm/harmonic/fundamental=\meta{positive integer} (initially 1)}
Reference order for normalisation and THD. For a $p$-pole-pair machine set it
to $p$. THD includes all other orders up to |maximum order|, including any
subharmonics, and excludes DC. It is a ratio, not a percentage.
\end{key}
\begin{key}{/femm/harmonic/normalize=\meta{true or false} (initially false)}
For spectrum plots, divide amplitudes by the selected fundamental. A zero
fundamental cannot be normalised. This key does not change scalar output.
\end{key}
\begin{key}{/femm/harmonic/order=\meta{positive integer} (initially 1)}
Order selected by scalar amplitude or phase output. It must not exceed
|maximum order|. Ignored for mean and THD.
\end{key}
\begin{key}{/femm/harmonic/quantity=\meta{choice} (initially amplitude)}
Scalar |amplitude|, |phase|, |mean| or |thd|. Phase is in degrees in the
convention $A\cos(n\theta+\phi)$, relative to the profile's angular origin.
A zero-amplitude phase is conventional, not physically meaningful.
Mean is the DC term; THD is undefined for a zero fundamental and raises an error.
\end{key}
\begin{key}{/femm/harmonic/plot options=\marg{PGFPlots options} (initially empty)}
Options passed to the emitted |\addplot+|. Axes, legends, colours, bar plots
and logarithmic scales remain standard PGFPlots settings.
\end{key}

Finite-element B is piecewise constant; plots deliberately retain those
jumps. Harmonic amplitudes are estimates from sampled FEM values, without
smoothing or forced sinusoidal fitting. Rotor-angle sweeps require distinct
static models and remeshing. Sliding air-gap elements, sector periodicity and
transient circuit coupling are separate extensions.

% SPDX-License-Identifier: LPPL-1.3c
% Copyright (C) 2026 Christophe Jorssen
% Author and Current Maintainer: Christophe Jorssen <christophe.jorssen@gmail.com>
% LPPL maintenance status: maintained
% This file is part of luafemm. See LICENSE and LICENSES.md for its terms.
\section{Ideal magnetic circuits and mean contours}
\label{sec:ideal}
An ideal circuit confines magnetic flux to a prescribed union of iron,
permanent magnets and air gaps. No flux leaves the lateral walls, including
the lateral walls of the gaps. This is a different physical boundary-value
problem from the ordinary open-air model: clipping the ordinary field picture
would not impose this condition.

Flux is conserved along every unbranched tube. At a junction the signed
incoming and outgoing fluxes balance. Neither statement requires a uniform
local field. A narrowing section increases its average normal induction,
and the local field also varies around corners and inside curved material.
The material laws remain finite and may be nonlinear. The word \emph{ideal}
does not replace iron by infinite permeability or remove its magnetomotive
force drop.

\subsection{A first confined circuit}
Let us begin with a complete U electromagnet. The node already knows where
its iron, winding sections and two gaps are located. In ideal mode it includes
the gaps in the computational domain, even if |gap mesh size=0|. It excludes
the winding sections and uses their enclosed currents as excitation.

\begin{codeexample}[width=6cm]
\begin{tikzpicture}[
  femm/problem={field model=ideal,
    depth=10,mesh size=4},scale=.55]
  \node[femm/u electromagnet={
    ampere turns=1200,
    gap mesh size=2}] (M) {};
  \pic[femm/lines=9] {femm field};
  \pic[femm/mean={component=M}]
    {femm mean path};
\end{tikzpicture}
\end{codeexample}

The red contour is a geometric reference for applying Ampere's law. It follows
the centres of the component sections and closes through the gaps. It need
not coincide with one of the blue field lines. Start with a coarse mesh while
arranging the figure; refine the mesh and curved geometry when examining local
field values or preparing the final result. Conservation of flux is a discrete
identity, but it does not by itself measure the accuracy of the predicted flux.

\begin{key}{/femm/problem/field model=\meta{open or ideal} (initially open)}
Choose the ordinary exterior-air model or the confined-domain model through
PGF's choice handler. Set it inside |femm/problem| before the picture begins.
|ideal| requires |mesher=delaunay| and at least one included physical region.
The bounding rectangle still encloses every declared vertex and bounds the
staging mesh, but its exterior air is removed before assembly.
Default rectangular boundary data are unused; nondefault exterior data,
Lua source/boundary callbacks and FEM import are rejected in this mode.
The option follows ordinary TeX scoping and participates in cache identity.
\end{key}

\begin{key}{/femm/region/ideal domain=\meta{true or false} (initially false)}
Mark a region as included in the confined computational domain. Supply it
inside |femm/region|; it inherits within the surrounding TeX scope. The last
covering region determines both the material and this membership flag.
Thus an air gap or a later material overlay needs |ideal domain=true| as well.
Compound-path holes preserve earlier declarations, as for ordinary materials.
The flag has no effect in |field model=open|.

Component nodes set the flag for their core, armature, magnet and gap parts,
and clear it for winding parts. A current density inside the active ideal
domain is rejected: windings must lie in excluded windows or exterior air.
All physical paths still accept the normal TikZ curves and transformations.
\end{key}

For a native path model, give the iron and the confined gaps their own styles.
The following complete example uses a rounded compound path, an explicit
window excitation and a user-drawn mean contour. Its gap section is also a
named profile, allowing the flux to be printed without numerical quadrature.
\luafemmexample{ideal-paths}

\subsection{Excitation and declaration order}
\begin{key}{/tikz/femm/excitation=\marg{options}}
Add a signed enclosed-current seed for ideal mode. Each declaration resets its
|x|, |y| and |ampere turns| options to zero before processing the supplied list.
Picture options and outer-scope styles queue their declarations until the
problem exists, independently of their position relative to |femm/problem|.
Inside a started picture the declaration is immediate. Several declarations
add; they do not replace one another. The seed is passive data, not a Lua
expression or a request to draw a winding.
\end{key}
\begin{keylist}[/femm/excitation]{x=\meta{number},y=\meta{number}}
Seed coordinates in model millimetres, initially zero for each declaration.
These are numerical model coordinates, unaffected by a local drawing transform.
At meshing time the point must lie strictly inside an excluded bounded window;
points on interfaces, inside magnetic material or in unbounded exterior air
are errors.
\end{keylist}
\begin{key}{/femm/excitation/ampere turns=\meta{number} (initially 0)}
Signed enclosed $NI$ in amperes. Positive values point along $+z$, and produce
positive counterclockwise circulation of $\mathbf H$ around that window.
Negative and zero values are accepted. Seeds add to currents already described
by winding regions. Do not add a seed for a winding that a component has
already declared.
\end{key}

The order is: select the problem and cache policy; declare all components,
physical paths and excitations; then request a field, a mean contour, a mesh,
a profile plot or a flux value. Every one of those requests can freeze the
physical model. Merely registering a |femm/profile| path still does not solve.
Named profiles retain their originating model after the picture ends.
Generic component nodes have exactly the same placement rules in both modes.

\subsection{Mean contours and transported-flux lines}
\begin{key}{/tikz/pics/femm mean path}
Resolve the component and cycle selected by |femm/mean|, solve on demand and
draw the resulting contour in the model's original picture frame. Local pic
translations do not detach it from the component. An unknown component,
unavailable cycle or contour outside the computational domain is an error.
\end{key}
\begin{key}{/tikz/femm/mean=\marg{options}}
Set the following options in the |/femm/mean| family. They inherit according
to normal PGF scoping; a pic's local options do not leak into the next pic.
\end{key}
\begin{key}{/femm/mean/component=\meta{node name} (initially empty)}
Required name of an already declared component node in the current model.
Use the full node name when applying a TikZ name prefix or suffix.
\end{key}
\begin{key}{/femm/mean/cycle=\meta{name} (initially main)}
Select |main| on a U electromagnet or either toroid; select |left| or |right|
on an E electromagnet. Bare U and E cores have no closed mean cycle.
The two E cycles share the central branch. They are contours for Ampere's law;
their shared segment is not a bifurcating field line.
\end{key}
\begin{key}{/femm/mean/kind=\meta{geometric or flux} (initially geometric)}
|geometric| follows section centres, orienting the contour with its net
circulation of $\mathbf B$. At zero field it keeps its construction direction.
Changing the excitation sign can therefore reverse its profile's starting
direction. No implication of uniform $\mathbf H$ along this contour is made.

|flux| selects an actual computed isoline halfway between the selected window
potential and its outer reference potential. It requires ideal mode, one
enclosed window and one unambiguous closed isoline surrounding that window
alone. It rejects a zero-flux window or a branched situation without such an
isolated line; it does not silently substitute a geometric contour. In a
single tube this line divides the transported flux equally, rather than
necessarily dividing the geometric section equally. Its starting point follows
the contour extractor's mesh ordering and can change after remeshing.
\end{key}
\begin{key}{/femm/mean/profile=\meta{identifier} (initially empty)}
Optionally register the chosen contour for |\femmplot|. Empty means no profile.
The identifier follows the existing profile rules; reusing it replaces the
previous profile. The normal offset is zero and outside-domain samples error.
\end{key}
\begin{key}{/femm/mean/samples=\meta{integer} (initially 201)}
Number of equally spaced arc-length profile samples, from 2 to 100000.
It changes only sampling, not the field mesh or the shape of the drawn contour.
\end{key}
\begin{stylekey}{/tikz/femm mean path}
The base style uses |draw=red!75!black|, |thick|, |dashed| and
|femm oriented field|. Use |/.append style| to retain the direction markers. The existing
arrow spacing, dimensions and enable/disable keys apply.
\end{stylekey}
\begin{key}{/tikz/femm/mean style=\marg{TikZ options}}
Set additional options applied after the base mean style, initially empty.
This replaces the extra options locally. |femm current mean| is internal
storage, not a physical declaration.
\end{key}

For a custom contour, use an ordinary closed |\draw| path with the
|femm mean path| style and |femm/profile={name=...}|. This keeps arbitrary
Bezier curves, rounded corners and named coordinates available. Its direction
is the direction in which the user writes the path; it is not automatically
reoriented. Automatic skeleton extraction from an arbitrary material domain
is deliberately not assumed.

\subsection{Flux values and Ampere profiles}
\begin{command}{\femmflux\oarg{options}}
Solve the source model of a named profile and print its signed section flux
in webers using PGF's standard number printer. Options belong to |/femm/flux|.
The command works in all three formats and requires a positive problem |depth|.
PGF's |/pgf/number format| styles control numerical presentation.
\end{command}
\begin{key}{/femm/flux/profile=\meta{identifier}}
Required named section profile; empty at the start of each command invocation.
Its path must remain in the computational domain and have |offset=0|. Declare
a displaced section explicitly if needed. For a path from $P$ to $Q$, with
right normal, the returned value is
\[
 \Phi_{P\to Q}=d\,[A_z(Q)-A_z(P)],
\]
where $d$ is the physical extrusion depth in metres. Reversing the path reverses
the sign. This right-normal convention is the opposite of the left normal
used by the profile component |Bn|. The result depends on endpoint potentials,
not on the number of samples. A closed section path returns zero.
\end{key}

\begin{key}{/femm/flux/snap tolerance=\meta{auto or nonnegative number} (initially auto)}
The default |auto| selects .003 model units for planar profiles and zero
for axisymmetric profiles, including imported models.
Endpoint projection tolerance in model millimetres. Before integrating, project
each endpoint to the nearest computational boundary if it is closer than this
distance. This accommodates small PGF coordinate differences between rotated
node anchors and captured part paths. Intermediate vertices and the physical
mesh are unchanged; the adjusted section must still remain in the domain.
Use zero to disable projection, or a smaller tolerance for very thin features.
The Lua |flux| API defaults to zero and accepts this tolerance as its third
argument. Like other option values, the TeX tolerance inherits in its scope.
\end{key}

To compare branch fluxes, place each section all the way between its two walls
and use consistent orientations. Use actual component pole anchors or captured
path endpoints: curved boundaries are polygonalised, and a point on the ideal
analytic circle can lie just outside the computational polygon.

\begin{key}{/femm/plot/component=circulation}
Additional PGFPlots component, in amperes:
$C(s)=\int_0^s\mathbf H\cdot d\boldsymbol\ell$ along the named profile.
The integral is evaluated by splitting each polyline segment at triangle
edges, then integrating the piecewise-constant element field exactly.
Increasing |samples| adds plotted abscissas; it does not improve the integral.
An offset profile or a path leaving the domain is rejected for this component,
even with |outside=nan|. At an exact material/element interface the usual first
containing triangle convention applies.

A closed mean contour should approach the linked $NI$ as the mesh is refined.
Finite-element $\mathbf H$ is not exactly tangentially continuous, so an
arbitrary contour's circulation need not equal $NI$ exactly on a coarse mesh.
This is distinct from the exact discrete conservation of magnetic flux.
\end{key}

Here are complete U examples with constant and unequal sections. Both show
the field, the geometric contour and its profile; the last plot displays
the magnetomotive force accumulated along the unequal-section circuit.
\luafemmexample{ideal-u}

\subsection{Additional component shapes}
All keys below extend |/femm/component|. They use numeric model millimetres
unless described as angles. They inherit within the surrounding TeX scope;
options on a node are local to that node. The existing material selectors,
mesh sizes, painting styles, affine placement and fixed node envelopes apply.
The E and toroidal shapes work in both field models. In open mode an air-gap
refinement region is optional; in ideal mode it is always included.

\begin{keylist}[/tikz/femm]{e core=\marg{options},e electromagnet=\marg{options}}
Select an E core open towards local $+y$, with a central leg at $x=0$.
The electromagnet adds an armature, two winding sections around the central
leg and up to three gaps. Positive |ampere turns| puts positive $J_z$ in
the left winding section and negative $J_z$ in the right. The core alone
adds no winding, armature, excitation or mean cycle.
\end{keylist}
\begin{keylist}[/femm/component]{left thickness=\meta{number},
 right thickness=\meta{number},yoke thickness=\meta{number},
 center thickness=\meta{number}}
Initially zero: inherit |leg thickness|. Positive overrides specify the
individual section thicknesses. Left, right and yoke apply to U and E shapes;
centre applies only to E. Negative values are errors. The resulting windows
and free leg heights must remain positive, and the winding must fit.
Changing one section does not change the independently specified armature.
\end{keylist}
\begin{keylist}[/femm/component]{left gap=\meta{number},
 center gap=\meta{number},right gap=\meta{number}}
E-only gap lengths, initially $-1$ to inherit |gap|. Otherwise require
nonnegative values; zero joins that pole directly to the armature.
The armature underside is at |height/2+max(left gap,center gap,right gap)|,
after resolving inherited values. Each leg stops its own gap distance below
this plane. A gap cannot consume the full leg above the yoke. These options
do not change the equal-gap U component.
\end{keylist}

The E shape provides the usual envelope and |core center| anchors, plus
|left pole|, |center pole|, |right pole|, |left window|, |right window|,
|yoke center| and |window center| (the left window's text position).
The outer legs also provide |left pole inner|, |left pole outer|,
|right pole inner| and |right pole outer|. The electromagnet additionally
provides |left gap|, |center gap|, |right gap|, the corresponding three
|... armature| anchors, |armature center| and both |coil positive| and
|coil negative|. Zero-length gaps still have a centre anchor.
An unavailable anchor reports an error rather than a guessed location.

\luafemmexample{ideal-e}

\begin{keylist}[/tikz/femm]{toroid=\marg{options},tapered toroid=\marg{options}}
Select an annular core centred at |core center|, with its gap towards local
$+x$. The tapered shape adds two straight-sided transitions adjacent to the
gap. A positive |magnet span| replaces the opposite arc with permanent magnet.
These are planar cross-sections with a constant extrusion depth. The taper
can illustrate the meridional outline of truncated-conical pole pieces, but
the calculation uses planar sections, not circular areas $\pi r^2$ and not
an axisymmetric or three-dimensional field model.
\end{keylist}
\begin{key}{/femm/component/radius=\meta{number} (initially 30)}
Mean radius. Must exceed half the core thickness.
\end{key}
\begin{key}{/femm/component/core thickness=\meta{positive number} (initially 10)}
Radial width of the regular toroidal section. Its physical area is this width
times the specified constant depth; it is not a diameter of a round section.
\end{key}
\begin{key}{/femm/component/pole thickness=\meta{positive number} (initially 6)}
Width at the gap on a tapered toroid, no larger than |core thickness|.
It is unused on an ordinary toroid or when |taper angle=0|.
\end{key}
\begin{key}{/femm/component/gap angle=\meta{number} (initially 8)}
Angular opening in degrees, from zero inclusive to 90 exclusive. Pole centres
are at angles $\pm\texttt{gap angle}/2$. The gap is the quadrilateral joining
their inner and outer endpoints. Zero closes the toroid. Unlike a U gap,
this is an angle: the pole-centre separation is
$2R\sin(\texttt{gap angle}/2)$.
\end{key}
\begin{key}{/femm/component/taper angle=\meta{number} (initially 25)}
Angular extent in degrees of each transition next to a tapered toroid's gap,
from zero inclusive to 90 exclusive. Straight sides join the narrow pole
section to the full-width annulus. The geometric mean follows their centre
chord. Zero disables tapering. It is unused by an ordinary toroid.
\end{key}
\begin{key}{/femm/component/magnet span=\meta{number} (initially 0)}
Angular extent of the magnet arc, centred at $180^\circ$, from zero inclusive
to 180 exclusive. Zero omits the magnet. The magnet must not overlap a taper.
The arc is subdivided into slices with tangent magnetisation; the existing
|magnetization angle| is an additional local rotation of that direction.
Set it to 180 to reverse the magnet. This prescribed direction is transformed
with the node. Hysteresis history and irreversible demagnetisation are not added.
\end{key}
\begin{key}{/femm/component/magnet material=\meta{material} (initially cast-alnico-5-lng37)}
Catalogue or user material style for the toroidal magnet arc. Set
|ampere turns=0| for a magnet without a winding; otherwise both excitations
are present. The example below selects a linear-recoil NdFeB catalogue entry.
\end{key}
\begin{key}{/femm/component/arc step=\meta{positive number} (initially 5)}
Maximum angular span in degrees of polygon edges and magnet slices, at most
30, subject to a maximum of 4096 subdivisions per full revolution.
An excessively small value is rejected before allocating geometry.
Refine this independently of |mesh size| to improve the geometric and
magnetisation-direction approximation. It does not change a native TikZ path's
|curve tolerance|. Two same-material arc pieces can leave visible drawing seams.
\end{key}
\begin{stylekey}{/tikz/femm component magnet}
Base magnet painting style: |draw=red!70!black,fill=red!20|.
\end{stylekey}
\begin{key}{/femm/component/magnet style=\marg{TikZ options}}
Append local painting options to the magnet style, initially empty.
Use it as for the existing core, armature, coil and gap styles.
\end{key}

The toroids provide |core center| and |window center| at the annulus centre,
|upper pole| and |lower pole| at the two section centres, their four
|... pole inner|/|... pole outer| endpoints, |gap center| and |magnet center|.
The latter marks the opposite arc position even if no magnet is present.
A nonzero |ampere turns| adds one positive winding section in the window and
one negative return section outside the core, using the existing coil width,
height and clearance. The inner rectangle must fit strictly inside the window;
clearance must be positive. Zero omits both rectangles. The two coil anchors
exist only when these rectangles exist. |coil center y| is not used by toroids.
The component's envelope expands to include the exterior winding.

\luafemmexample{ideal-toroids}

\subsection{Discrete conservation, accuracy and interchange}
The confined domain is carved from an audited constrained mesh. Each outer
boundary has $A_z=0$ as a reference; all vertices of each window boundary
share one floating degree of freedom. The corresponding equation prescribes
the enclosed current, including explicit seeds and excluded winding regions.
Permanent magnets contribute through the usual constitutive weak form.
Independent material components have independent outer references.

Continuous P1 $A_z$ gives $\mathbf B=(\partial_y A_z,-\partial_x A_z)$.
Its normal component agrees across a shared triangle edge, and is zero on
an equipotential wall. Consequently any two consistently oriented complete
sections of one branch give the same endpoint potential difference, to
floating-point precision. Across a junction these differences telescope,
giving the signed flux balance. This remains true for nonlinear materials
and unequal sections. Local magnitudes, tangential continuity of $\mathbf H$
and contour circulations still need mesh-convergence checks.

Ideal field pictures space their levels over the boundary-potential range,
so the level spacing measures transported flux. A finite number of lines
does not encode every branch ratio as an exact ratio of integer line counts.
Permanent magnets can also produce local recirculation; confinement constrains
the net branch flux, not uniform pointwise induction. Explicit Lua contour
levels can inspect such additional potential extrema if needed.

Caches store the confined mesh and validate its boundary-potential ties on
restoration. Cache revision 13 rejects incompatible records. FEMM import
and FEM/ANS export are not offered for this mode: the floating-window/current
constraints are not represented by the existing interchange subset, and an
export must not silently turn the problem back into an open-air calculation.
The ordinary mode retains its existing interchange support.

% SPDX-License-Identifier: LPPL-1.3c
% Copyright (C) 2026 Christophe Jorssen
% Author and Current Maintainer: Christophe Jorssen <christophe.jorssen@gmail.com>
% LPPL maintenance status: maintained
% This file is part of luafemm. See LICENSE and LICENSES.md for its terms.
\section{Material library and constitutive laws}
\label{sec:materials}
The embedded catalogue contains all 246 records from the supplied FEMM
magnetic material library: air, permanent magnets, soft magnetic materials,
nonmagnetic metals and winding wires. A material is selected by an ASCII
identifier, such as |pure-iron|, |copper|, |aluminum-6061-t6| or
|cast-alnico-5-lng37|. See Appendix~\ref{sec:catalogue} for every choice.
This is the magnetic solver's catalogue; dielectric and thermal material
databases belong to other FEMM solvers and are not imported.

The names, folders, every numeric field and every B--H point are preserved.
The two records named Supermalloy have distinct identifiers:
|supermalloy| and
|supermalloy-metals-handbook-dc-magnetization-curves|. No record is
overwritten by name matching. Installed documents need neither the FEMM
clone nor its executable.

\subsection{Custom and derived materials}
Define a material as a style under |/tikz/femm/materials|. The style name
then becomes a valid |material| choice. Define styles before the first
use, and use a name that is not reserved by the catalogue.
\begin{codeexample}[code only]
\tikzset{
  femm/materials/my-iron/.style={
    bh={0/0,.5/80,1/180,1.3/350,1.5/900,1.8/7000,2/20000}},
  femm/materials/my-metal/.style={library=copper,mu r=1.01},
  femm/materials/recoil-magnet/.style={mu r=1.05,remanence=1.2},
  core/.style={femm/region={material=my-iron},draw,fill=gray!20}
}
\end{codeexample}

The following keys are evaluated while resolving a material style. Empty
internal slots mean ``inherit from the library, if any''. The effective
defaults below apply when there is no library base. Explicit properties
override the base, independently of their order in the style.

\begin{key}{/tikz/femm/materials/library=\meta{choice}}
Use a catalogue record as the base of a custom material. This is a PGF
|.is choice| handler, with no initial base. A private copy is created;
overrides never change the catalogue itself. A material name is model-wide:
do not redefine its physical law between regions of the same model.
\end{key}
\begin{key}{/tikz/femm/materials/mu r=\meta{positive number}}
Relative permeability of a linear isotropic material; effective default 1.
A nonempty |bh| table takes precedence. Setting |mu r| on a derived
nonlinear material does not remove its B--H table.
\end{key}
\begin{key}{/tikz/femm/materials/bh=\marg{B/H pairs}}
Comma-separated numeric pairs: B in tesla, H in A/m. The table must have
at least two points, begin with |0/0|, and have strictly increasing B and H.
There is no curve by default. A nonempty table replaces the base table;
an empty value inherits it. The package interpolates H(B) linearly. Beyond
the final point, the slope is $1/\mu_0$. This is deliberately not FEMM's
spline interpolation, so imported data do not imply bitwise FEMM results.
\end{key}
\begin{key}{/tikz/femm/materials/coercivity=\meta{nonnegative number}}
Prescribed source strength $H_c$ in A/m; effective default zero.
It can accompany either |mu r| or a \emph{shifted} B--H table. Its direction
is chosen by each region's |magnetization angle|.
\end{key}
\begin{key}{/tikz/femm/materials/remanence=\meta{nonnegative number}}
Remanent induction $B_r$ in tesla, for a linear recoil law; effective
default zero. A positive value is incompatible with a B--H table or
nonzero |coercivity|, including inherited ones. It gives
$H_c=B_r/(\mu_0\mu_r)$. This extrapolated intercept is not necessarily the
measured coercivity of a nonlinear hard magnet.
\end{key}

\subsection{Permanent magnets}
The constitutive convention is
\[
 \mathbf H=\nu(|\mathbf B|)\mathbf B-H_c\mathbf m,
 \qquad \mathbf m=(\cos\alpha,\sin\alpha),\qquad
 \nu(B)=H_{\rm table}(B)/B.
\]
Along the positive magnetisation direction, the physical field is
$H_{\rm physical}(B)=H_{\rm table}(B)-H_c$. Thus the table starts at
|0/0| while the physical demagnetisation branch starts at $H=-H_c$.
Remanence occurs where $H_{\rm table}(B_r)=H_c$. Do not paste a measured
negative-H demagnetisation table directly into |bh|.

The |cast-alnico-5-lng37| record has $H_c=48600$ A/m and a table endpoint
at $B=1.204$ T, $H_{\rm table}=48600$ A/m. The example
|examples/alnico-u.tex| uses that table with illustrative soft iron;
|examples/materials-profiles.tex| uses the catalogue's |pure-iron| instead
(complete source and result in Section~\ref{ex:materials-profiles}).
Both produce a field without imposed current. The magnetisation direction
is prescribed, not an internal history variable: hysteresis, irreversible
demagnetisation and measured transverse anisotropy are not modelled.

\subsection{Complete example: explicit material tables}
The following U magnet shows the original illustrative soft-iron curve
and the shifted alnico table explicitly in reusable material styles.
The local air-gap refinement regions are declared after the solids and
occupy only the gaps. The output includes field lines, geometry and the
computed diagnostics. These illustrative iron data are not a commercial
steel specification.
\luafemmexample{alnico-u}

\subsection{Static interpretation of the catalogue}
All records in this revision have equal |Mu_x| and |Mu_y|. For |LamType=0|,
a fill fraction $f$ is applied as
\[
 \mu_{r,\rm eff}=f\mu_r+(1-f),\qquad
 B_{\rm eff}(H)=f B(H)+(1-f)\mu_0H.
\]
This accounts for the ten records with fill fraction 0.98. The winding
types 3 and 8 have unit relative permeability in the static model, as in
FEMM's static assembly. Custom linear anisotropy and directional sheet mixtures are described in
Section~\ref{sec:static}. The catalogue adapter retains its audited isotropic
interpretation and rejects unsupported database combinations.

Conductivity is used by the named DC winding model in
Section~\ref{sec:static} and by the harmonic diffusion solver in
Section~\ref{sec:harmonic}. Section~\ref{sec:losses} adds supported linear principal loss angles and
insulated-sheet thickness/fill. Nonlinear catalogue B/H curves retain their
static meaning; selecting them does not create a nonlinear harmonic law.
Wire diameter and strand count remain metadata for unsupported wire models. The real material source |J_re| is converted from A/mm$^2$ to A/m$^2$.
\emph{Sixty-nine copper-clad aluminium wire records specify 1 A/mm$^2$.}
Use |current density=0| to make one of those regions passive. This value
is also shown in the machine-readable and Markdown catalogues.

\subsection{Provenance and data terms}
The source is FEMM revision
\nolinkurl{7d9e8ed90772ddb939a3256f976e231bcb5cca97},
file |bin/matlib.dat|. The repository keeps the original bytes, the
upstream notice and a deterministic Lua importer. The imported data are
identified separately from the original LPPL code of this package; FEMM's
Aladdin Free Public License notice is included in |data/femm/LICENSE.txt|.
The release does not relabel the FEMM data as LPPL data.

% SPDX-License-Identifier: LPPL-1.3c
% Copyright (C) 2026 Christophe Jorssen
% Author and Current Maintainer: Christophe Jorssen <christophe.jorssen@gmail.com>
% LPPL maintenance status: maintained
% This file is part of luafemm. See LICENSE and LICENSES.md for its terms.
\section{Solving, field lines, mesh and numeric results}
\label{sec:field}
\begin{key}{/tikz/femm/solve}
Resolve the current problem when this path executes, for example
|\path[femm/solve];|. This is optional: field pictures and requested
profile values solve on demand. Declare all regions before the first
mesh or solve request. A mesh-only picture does not solve the system.
Section~\ref{sec:workflow} gives the complete declaration and computation order.
\end{key}
\begin{command}{\femmsolve}
Command form of the same solve operation. A current model must exist.
\end{command}

\begin{key}{/tikz/pics/femm field}
Used as |\pic {femm field};|. Solve if necessary and draw oriented
isolines of $A_z$. Each level can have several connected components.
Because $\mathbf B=(\partial_y A_z,-\partial_x A_z)$, these are field lines
of the computed P1 solution. They are polygonal, without graphical smoothing.
\end{key}
\begin{key}{/tikz/pics/femm mesh}
Used as |\pic {femm mesh};|. Generate the mesh if necessary, then draw each
edge once. It is useful for inspecting interfaces and local refinement.
\end{key}
Both pictures are anchored to the original model frame. Moving the pic
inside a local scope does not detach the field from the media. Use a clip
to display a smaller part of the computed domain.

\begin{key}{/tikz/femm/lines=\meta{positive integer} (initially 29)}
Number of evenly spaced potential levels. Levels are placed at the
midpoints of equal bins between the minimum and maximum potential.
In ideal mode the range uses boundary potentials, representing transported
flux rather than additional local recirculation extrema (Section~\ref{sec:ideal}).
This counts levels, not necessarily visible connected curves. Zero-field
models produce no contours.
\end{key}
\begin{key}{/tikz/femm/arrows=\meta{true or false} (initially true)}
Enable or disable the direction markers. The direction follows B, not an
arbitrary contour traversal. With no value, PGF's boolean handler selects true.
\end{key}
\begin{key}{/tikz/femm/arrow spacing=\meta{dimension} (initially 22mm)}
Distance between repeated markers along the rendered path. Markers occupy
the middle part of each contour (from relative positions .12 to .95).
Short paths need not contain a marker. This is a graphic dimension, not
a mesh length or physical parameter.
\end{key}
\begin{key}{/tikz/femm/arrow length=\meta{dimension} (initially 1.3mm)}
Length of the Stealth marker, in the displayed figure.
\end{key}
\begin{key}{/tikz/femm/arrow width=\meta{dimension} (initially .85mm)}
Width of the Stealth marker.
\end{key}
\begin{key}{/tikz/femm/field style=\marg{TikZ options}}
Extra graphics options for this field picture; empty initially. For
example |{teal!70!black,line width=.5pt}|. The base |femm field| style is
applied first, and these options second. This key sets rather than appends
the extra options within the current TeX scope.
\end{key}
\begin{key}{/tikz/femm/mesh style=\marg{TikZ options}}
Extra mesh graphics options, empty initially. The base |femm mesh| style
is applied first. No material or mesh-size property is changed by a
graphic style.
\end{key}
\begin{stylekey}{/tikz/femm field}
Base contour style: |blue!65!black,line width=.32pt,femm oriented field|.
Extend with |/.append style| if the standard arrow handling is to remain.
\end{stylekey}
\begin{stylekey}{/tikz/femm mesh}
Base mesh style: |draw=black!25,line width=.12pt|.
\end{stylekey}
\begin{key}{/tikz/femm oriented field}
Reusable marker setup used by |femm field|. It respects |femm/arrows|
and the marker dimensions. PGF's floating-point vector length is enabled
locally for decoration: very short contour segments would otherwise
overflow PGF's fixed-point reciprocal. The contour vertices are unchanged.
\end{key}
The helper styles |femm current field| and |femm current mesh| are internal
storage for the corresponding extra-style keys; use those keys rather
than depending on the helper implementation.

\subsection{Numeric output}
\begin{command}{\femmB\marg{x}\marg{y}}
Print $|\mathbf B|$ in tesla, with three digits after the decimal point.
Coordinates are model-unit numeric literals. The current model must
already be solved; the first containing triangle is used at interfaces.
For more precision or components, use a profile or the Lua API.
\end{command}
\begin{command}{\femmstat\marg{name}}
Print a solved-model statistic with up to eight significant digits.
Unknown or unavailable values are errors.
\end{command}
\begin{center}
\begin{tabular}{ll}
\toprule Name & Meaning \\
\midrule
|nodes|, |triangles| & Mesh vertex and element counts \\
|newton| & Accepted Newton corrections \\
|cg| & Total conjugate-gradient iterations \\
|residual| & Final relative nonlinear residual \\
|seconds| & CPU solution time, excluding mesh generation \\
|peak| & Largest element value of $|\mathbf B|$ in T \\
|net_current| & Integral of $J_z$ over the section, in A \\
\bottomrule
\end{tabular}
\end{center}
\begin{command}{\femmmeshstat\marg{name}}
Print a constrained-mesh diagnostic with up to six significant digits.
These values require the Delaunay mesher; they do not exist for |grid|.
\end{command}
\begin{center}
\begin{tabular}{ll}
\toprule Name & Meaning \\
\midrule
|min_angle| & Smallest triangle angle in degrees \\
|max_area| & Largest triangle area in model units squared \\
|mesh_tolerance| & Actual vertex-merging distance in model units \\
|seconds| & CPU mesh generation and audit time \\
|flips| & Constraint and refinement legalisation flips \\
|steiner_points| & Additional quality-refinement vertices \\
|segment_splits| & Quality-refinement segment bisections \\
|circumcenters| & Accepted quality-refinement circumcentres \\
|exact_orientation|, |exact_incircle| & Exact-arithmetic fallback counts \\
|exact_diametral| & Exact diametral-disk fallback count \\
\bottomrule
\end{tabular}
\end{center}
Statistics remain accessible after the source picture's TeX group closes.
Starting another model changes the current model for these commands;
named profiles, in contrast, retain their own source model.

% SPDX-License-Identifier: LPPL-1.3c
% Copyright (C) 2026 Christophe Jorssen
% Author and Current Maintainer: Christophe Jorssen <christophe.jorssen@gmail.com>
% LPPL maintenance status: maintained
% This file is part of luafemm. See LICENSE and LICENSES.md for its terms.
\section{Saving and reusing computations}
\label{sec:cache}
The mesh and converged solution can be reused across compilations. This is
particularly useful while editing labels, field-line styles or a PGFPlots
axis. The document still declares its physical paths: they identify the
requested calculation and let luafemm detect outdated results.

\begin{codeexample}[width=5cm]
\begin{tikzpicture}[
  femm/problem={xmin=-20,xmax=20,ymin=-20,ymax=20,mesh size=2},
  femm/cache={file=magnet-field,mode=auto}]
  \filldraw[femm/region={material=cast-alnico-5-lng37,
    magnetization angle=90},fill=red!20]
    (-4,-6) rectangle (4,6);
  \draw[femm/profile={name=probe,samples=101},dashed]
    (-12,9)--(12,9);
  \pic[femm/lines=17] {femm field};
\end{tikzpicture}
\end{codeexample}

On the first compilation the field request meshes and solves the problem,
then saves the result. Later compilations restore the nodal potential and
derive the element fields from the current material declarations. Profiles
and field lines are regenerated on demand, so their graphical and sampling
options remain independent of the saved calculation. For |type=axial|,
the saved scalar is the dimensionless axial H perturbation, not a vector
potential; current loops and profiles are regenerated from its complex values.
The current passive cache format is 11; older snapshots must be recomputed.
Operating-point responses additionally fingerprint the bias field, its exact
ordered mesh and the response law; see Section~\ref{sec:operating}. The cache
choice |mode=frozen| and the constitutive choice |permeability=frozen| are
independent.

\subsection{Picture options}
\begin{key}{/tikz/femm/cache=\marg{options} (default \normalfont empty)}
Configure persistent reuse for a model. Like |femm/problem|, this key
belongs in the picture options, before the model's begin-picture hook
runs. It may precede or follow |femm/problem|. Without this key, caching
is off. Invoking it sets |mode=auto| before processing its subkeys.
Thus |femm/cache| alone enables automatic reuse with the default filename.
Normal PGF grouping and inherited styles apply.
\end{key}
\begin{key}{/femm/cache/file=\meta{filename stem} (initially \string\jobname-femm)}
The package appends |.lfc|. For example, |file=results/u-magnet| uses
\nolinkurl{results/u-magnet.lfc}. Macros are expanded and passed as escaped
strings. Spaces in filenames are supported.

Relative names are resolved under TeX's output directory when one is set,
otherwise under the compilation's working directory. With
|--output-directory=build|, |file=magnet| therefore writes
\nolinkurl{build/magnet.lfc}. Absolute names are used as supplied.
Any parent directory must already exist; luafemm does not create directories.
Use a distinct stem for each physical problem and for concurrent writers.
The default is convenient for a document containing one magnetic model.
\end{key}
\begin{key}{/femm/cache/mode=\meta{choice} (initially off)}
A PGF |.is choice| selector. Unknown choices are errors. The enclosing
|femm/cache| key selects |auto| unless explicitly overridden.
\end{key}
\begin{key}{/femm/cache/mode/auto}
Restore a matching solution. If only the mesh matches, restore and
reclassify it with the current materials and excitations, then solve.
Otherwise generate a new mesh and solve. Missing, damaged or incompatible
records produce a logged miss and are replaced after successful preparation.
\end{key}
\begin{key}{/femm/cache/mode/refresh}
Ignore disk contents, generate a new mesh and solution when requested,
and replace the cache. Repeated requests within the same model reuse its
in-memory data; refresh does not mean solving for every field picture.
\end{key}
\begin{key}{/femm/cache/mode/off}
Neither read nor write cache files. Existing files are left untouched.
This is the behaviour of documents that do not enable |femm/cache|.
\end{key}
\begin{key}{/femm/cache/mode/frozen}
Require compatible saved data and never write a file. A mesh picture
requires a matching mesh; a field picture, explicit solve or sampled
profile requires a matching converged solution. A missing, damaged,
incompatible or outdated record raises an error before numerical work.
It never silently displays an outdated field or falls back to a solve.
Changing presentation options remains allowed.
\end{key}

\subsection{What invalidates the result?}
\begin{center}
\begin{tabular}{p{.57\linewidth}p{.34\linewidth}}
\toprule Change & Work in |auto| mode \\
\midrule
Domain, contours, fill rule, units, mesher, spacing, mesh tolerance or quality controls
  & Generate mesh and solve \\
Material law, coercivity, region current, magnetisation or boundary data
  & Reuse mesh; solve again \\
Solver tolerance or iteration limit & Reuse mesh; solve again \\
Colour, stroke, arrows, field-line count, clip, labels
  & Reuse solution \\
Profile path, sample count, component, offset or PGFPlots options
  & Reuse solution; sample again \\
Incompatible package or cache-format version & Generate mesh and solve \\
\bottomrule
\end{tabular}
\end{center}

The geometry descriptor uses the evaluated, flattened and rounded model
coordinates, not the source text. A curve-tolerance change invalidates the
mesh when it changes those coordinates. Whole-picture transformations
normally affect only presentation; transformations inside a physical path
affect the model. The current picture frame is always retained. Precision
effects that change captured coordinates can cause a conservative cache miss.
Region order matters because later filled regions take material precedence.
Every subcontour and the native fill rule are included, even when a hole
contains only background air. A fill-rule change can alter local refinement
and therefore invalidates the mesh as well as the solution.

The solution descriptor contains the actual constitutive data, including
B--H points and effective coercivity, rather than just catalogue identifiers.
Unused materials, log callbacks, frames and profile definitions are excluded.
Descriptors are compared in full; a short checksum collision cannot cause a
stale model to match. This is a single-record cache, not a history of parameter
sets: changing a model back may require another solve.

\subsection{Diagnostics and saved statistics}
\begin{command}{\femmcachestatus}
Print the current model's cache outcome without starting a calculation:
|off|, |pending|, |miss|, |mesh|, |solution| or |refresh|.
|pending| means caching is enabled but no mesh or solve has been requested.
|mesh| means only the mesh was restored; after an automatic solve it still
records that reuse outcome. |solution| means a converged solution was restored.
|miss| and |refresh| describe how the current calculation was obtained, even
after it has been saved successfully. A current model must exist.
\end{command}

The log gives the resolved filename and the reason for a miss or partial
reuse. Lua scripts can inspect |model.cache_info|. A restored solution
retains the original Newton/CG counts, residual and solution time in
|\femmstat|; the mesh diagnostics likewise describe its original generation.
In particular, |seconds| is \emph{not} the current cache-loading time.
All material and region declarations must precede the first mesh or solve.

\subsection{File handling and scope}
The private |.lfc| format stores passive, versioned data with lossless
binary64 round trips: geometry and physical descriptors, SI mesh vertices,
connectivity, exterior tags, diagnostics and, when available, nodal potential.
It is parsed without |load|, |loadfile| or |dofile|. Size, integrity and
structural checks run before restoration. A checksum detects accidental
damage; it is not an authentication mechanism.

Writes use a neighbouring temporary file and a final rename. A failed write
or replacement leaves the previous target intact and raises an error. If
the filesystem cannot replace an existing target through rename, the
operation fails explicitly. File-access errors on writing are not silently
ignored. Only a converged solution is saved; an unsuccessful solve may
leave a valid mesh-only cache. Interrupted processes can leave temporary
files, which are not considered cache records.

The current format supports both luafemm meshers and can be shared among
plain LuaTeX, LuaLaTeX and ConTeXt MkIV using the same package version.
Lua-only source and boundary closures require caching to be off, because
their captured state cannot be fingerprinted reliably.

These files are disposable computational caches, not FEMM |.fem| or |.ans|
files and not a replacement for the document source. FEMM exchange and
loading a foreign model without TeX declarations are outside this version.
Keep the source, material data and package version to reproduce a result;
use |frozen| to prevent an unintended recomputation while editing.

% Ideal boundary constraints are rebuilt from physical declarations.
Ideal models also include domain membership and window excitations in their
cache descriptors, and validate the restored boundary-potential ties.
Their mean-contour styles and profile names do not affect the physical cache.

\subsection{Moving-band positions}
A sliding template reuses a mesh-only cache; each angle/current state reuses a
separate solution cache. Format 11 includes the new method, angle, exact assembled
mesh, force weights and template provenance in these positions' signatures.
Ordinary fingerprints and stored fields are unchanged. Named templates are
recaptured each run. See Section~\ref{sec:sliding}.

% SPDX-License-Identifier: LPPL-1.3c
% Copyright (C) 2026 Christophe Jorssen
% Author and Current Maintainer: Christophe Jorssen <christophe.jorssen@gmail.com>
% LPPL maintenance status: maintained
% This file is part of luafemm. See LICENSE and LICENSES.md for its terms.
\section{Measurement paths and the PGFPlots library}
\label{sec:profiles}
Load |\usepgfplotslibrary{femm}| after PGFPlots for the standard library
entry point. It loads the generic TikZ library if necessary. The profile
capture belongs to the TikZ library; only plotting requires PGFPlots.

\begin{key}{/tikz/femm/profile=\marg{options}}
Register an evaluated measurement path, without adding a physical region.
Use |\draw| to display the path or |\path| to keep it invisible. A single
connected open or closed path is required; arcs, circles, Bézier curves,
rounded corners, named coordinates and local transformations are accepted.
Do not combine |femm/region| and |femm/profile| on the same path.
Capture itself does not mesh or solve the model.
It may precede or follow physical declarations; only the first sampling
or plotting request requires those declarations to be complete. See
Section~\ref{sec:workflow}. Each profile remains one connected curve;
compound material paths require separate profiles for their separate contours.
\end{key}

The profile subkeys belong to |/femm/profile|:
\begin{key}{/femm/profile/name=\meta{identifier}}
Required name, with letters, digits, hyphens or underscores. The registry
is document-wide; reusing a name replaces its old geometry and cached
values. A profile keeps its own model even after another picture starts.
\end{key}
\begin{key}{/femm/profile/samples=\meta{integer} (initially 201)}
Number of equally spaced arc-length samples, including both endpoints;
between 2 and 100000. Increasing this number refines the graph's sampling,
not the finite-element solution. A narrow peak may be missed between samples.
\end{key}
\begin{key}{/femm/profile/curve tolerance=\meta{positive number} (initially .03)}
Flattening tolerance in model units. Independent of the material-path
tolerance. Samples follow the resulting polyline, so its arc length is
an approximation to the original smooth curve's length.
\end{key}
\begin{key}{/femm/profile/offset=\meta{number} (initially 0)}
Signed displacement towards the path's left normal, in model units.
Use a small nonzero offset to choose a side of a material interface.
The reported distance $s$ remains measured on the original path; the
sampled $x,y$ coordinates include the displacement. At sharp vertices,
offsets follow the chosen segment normal, not a rounded offset curve.
\end{key}
\begin{key}{/femm/profile/outside=\meta{error or nan} (initially error)}
PGF choice controlling samples outside the rectangular domain. |error|
stops with the offending point. In ideal mode the excluded air and windows
also count as outside. |nan| supplies missing values, creating
gaps in the graph. It does not hide a mesher or solver failure.
\end{key}

\begin{command}{\femmplot\oarg{options}}
Inside a PGFPlots axis, emit a complete |\addplot+ ... coordinates ...;|
command. Do not append another semicolon. The first request solves the
profile's source model if necessary and caches all components. Further
component plots reuse those rows. Values pass in memory; no data file,
Python process or shell escape is involved. Options are parsed locally;
|\addplot+| executes in the enclosing axis so that legends and cycle lists
advance normally.
\end{command}

Every call resets the following plot options, in family |/femm/plot|:
\begin{key}{/femm/plot/profile=\meta{identifier}}
Required name of a previously registered profile.
\end{key}
\begin{key}{/femm/plot/component=\meta{choice} (initially B)}
The additional |circulation| choice integrates $\mathbf H$ along a zero-offset
path; see Section~\ref{sec:ideal}. Standard choices are |B|, |Bx|, |By|, |Bt|, |Bn|, |H|, |Hx|, |Hy|, |Ht|, |Hn| or |Az|.
Unsubscripted B and H mean the magnitudes. B components are in T, H
components in A/m and Az in T\,m. Cartesian components refer to the model
frame. Tangential and normal components follow the path direction.
\end{key}
\begin{key}{/femm/plot/abscissa=\meta{choice} (initially s)}
|s| is distance from the first point in model units; |t| is $s/L$ from
zero to one; |x| and |y| are sampled coordinates. In the TikZ interface,
distances are in millimetres. The parameter |t| is \emph{not} the polynomial
parameter of a Bézier curve. Nonmonotone x/y paths are plotted in traversal
order, not sorted by the chosen abscissa.
\end{key}
\begin{key}{/femm/plot/plot options=\marg{PGFPlots options}}
Ordinary options appended to |\addplot+|, empty by default. For example
|{red,thick,dashed,no markers}|. All axis labels, bounds, grids, legends and
cycle lists remain PGFPlots options. The emitted plot also sets
|unbounded coords=jump| unless explicitly overridden here.
\end{key}

\clearpage
\subsection{A curve and its components}
This example includes both the source picture and the graph and is
executed during the manual build.
\begin{codeexample}[width=6cm]
\begin{tikzpicture}[femm/problem={
  xmin=-15,xmax=15,ymin=-15,ymax=15,
  mesh size=2}]
  \filldraw[femm/region={
    material=cast-alnico-5-lng37,
    magnetization angle=90},
    fill=red!20]
    (-3,-4) rectangle (3,4);
  \draw[femm/profile={name=demo,
    samples=51},blue,dashed]
    (-8,6) .. controls (0,12) .. (8,6);
\end{tikzpicture}
\par\medskip
\begin{tikzpicture}
\begin{axis}[width=5.8cm,height=4.5cm,
  xlabel={$s$ (mm)},ylabel={Field (T)},
  no markers]
  \femmplot[profile=demo,component=B,
    plot options={black,thick}]
  \addlegendentry{$|\mathbf B|$}
  \femmplot[profile=demo,component=By,
    plot options={red}]
  \addlegendentry{$B_y$}
\end{axis}
\end{tikzpicture}
\end{codeexample}

\subsection{Complete example: straight and curved scans}
This example combines catalogue iron, alnico and passive copper. The first
scan traverses both air gaps and plots Cartesian components. The second
follows a cubic curve and plots tangent and normal components. Both graphs
retain the finite-element jumps rather than applying graphical smoothing.
\luafemmexample{materials-profiles}

\subsection{Tangents, interfaces and plotting conventions}
For unit tangent $\hat{\mathbf t}$, the left normal is
$\hat{\mathbf n}=(-\hat t_y,\hat t_x)$.
The signed components are $B_t=\mathbf B\cdot\hat{\mathbf t}$ and
$B_n=\mathbf B\cdot\hat{\mathbf n}$, and likewise for H. Reversing the
path changes these signs at a common point, while Bx and By are unchanged.
At a polygon vertex, the following segment supplies the tangent; the
final point uses the last segment. A closed path repeats its first
position as its final sample, but endpoint tangents can differ at a corner.

The P1 potential is continuous. Its gradient, and therefore B and H,
is constant in each triangle and may jump between elements or materials.
No smoothing is applied. PGFPlots normally joins consecutive samples by
straight segments. |const plot| changes the graphic convention but does
not locate all triangle-edge intersections. At an exact interface the
first containing triangle supplies the value; an offset or a path wholly
inside the desired medium makes the sampling side explicit.

% SPDX-License-Identifier: LPPL-1.3c
% Copyright (C) 2026 Christophe Jorssen
% Author and Current Maintainer: Christophe Jorssen <christophe.jorssen@gmail.com>
% LPPL maintenance status: maintained
% This file is part of luafemm. See LICENSE and LICENSES.md for its terms.
\section{Regional measurements, forces and torques}
\label{sec:postprocessing}
This chapter describes planar measurements. Axisymmetric volume weights,
component names and force restrictions are specified in Section~\ref{sec:axisymmetric}.
The planar postprocessor works on the same unsmoothed, piecewise-constant field that
produces the field lines. A \emph{selection} identifies existing elements;
a \emph{measurement} solves if necessary and records a named numerical snapshot.
The snapshot can be printed, plotted or drawn as a force arrow. There is no
external process, temporary data table or shell escape.

\subsection{From a drawing to a measured attraction}
Declare all physical regions and set a positive problem |depth| first. Then
select the moving body, request its force and draw the result. For example,
after constructing the U electromagnet used below:
\begin{codeexample}[code only]
\path[femm/selection={name=armature,material=pure-iron}]
  (-42,31) rectangle (42,55);
\path[femm/measure={name=pull,selection=armature,quantity=force}];
\pic[femm/force={result=pull}] {femm force};
\end{codeexample}
The selection rectangle is not another material region. It restricts the
|pure-iron| filter to the armature, leaving the core out. The default force
method automatically constructs a weighting field through the surrounding air.
Its orange arrow starts at the selected area's centroid, has an illustrative
length of 12 model units and carries a label in newtons.

The complete example varies both gaps together, retains the three results and
plots the attraction. Start with these draft settings, then reduce mesh size
and gap spacing targets and compare the \emph{force}, not just the appearance
of the field lines. Increasing the number of plotted field lines does not
improve the mechanical result.
\luafemmexample{forces}

\subsection{Selecting existing elements}
\begin{key}{/tikz/femm/selection=\marg{options}}
Capture a closed selection path through a TikZ postaction. Supports curves,
transformations, compound paths and the current nonzero/even-odd fill rule.
Do not combine this capture with |femm/region| or |femm/profile| on one path.
Capture does not mesh or solve. It may occur before or after the physical
declarations, but those declarations must be complete before measurement.
\end{key}
\begin{command}{\femmselect\oarg{options}}
Register a selection without a geometric mask, normally by material.
For example, |name=iron| with |material=pure-iron| selects every element assigned
that material. With an empty material it selects the whole active domain.
\end{command}
\begin{key}{/tikz/femm/select=\marg{options}}
Key form of |\femmselect|, for example |\path[femm/select={name=all}];|.
Use after the problem has started, not among its pre-start picture options.
\end{key}
The selection family is |/femm/selection|:
\begin{key}{/femm/selection/name=\meta{identifier}}
Required document-wide name: letters, digits, underscores and hyphens only.
Reusing a name replaces the selection. Each selection retains its source model;
using it to measure another model is an error.
\end{key}
\begin{key}{/femm/selection/material=\meta{name} (initially empty)}
Exact model material name, including catalogue identifiers. Empty accepts any
material. Combined with a path, both tests must pass. This tests final element
assignments, so later material painting and holes are already accounted for.
\end{key}
\begin{key}{/femm/selection/curve tolerance=\meta{number} (initially .03)}
Positive model-unit tolerance for flattening the selection path. It neither
refines the physical geometry nor changes the mesh. Command selections ignore it.
\end{key}
Elements are selected as a whole by their centroid. A mask cutting through a
material is therefore a mesh-dependent selection, not an exact clipped integral.
For a mechanical body, enclose the whole body and filter by material when
necessary. Empty selections fail explicitly. Existing selection paths cannot
recover a physical interface absent from the solved mesh.

\subsection{Measuring and naming a result}
\begin{command}{\femmmeasure\oarg{options}}
Evaluate on the current model. Requires an explicit positive problem |depth|
(millimetres in the TikZ frontend), even for a result also containing area
averages. It must follow all physical declarations. It solves on demand and
can reuse the physical solution cache. The auxiliary force-weight problem may
still run on a cache hit. It does not modify the magnetic solution.
\end{command}
\begin{key}{/tikz/femm/measure=\marg{options}}
Equivalent key form, intended for a standalone |\path[femm/measure={...}];|
after physical declarations. It executes while the key is processed; it is
not delayed until the end of the picture or until after a region postaction.
Do not put it on the same path as a physical declaration.
\end{key}
Each measurement parses options locally, resetting the following defaults.
Ordinary PGF styles compose within |/femm/measure|:
\begin{codeexample}[code only]
\pgfkeys{/femm/measure/armature pull/.style={
  selection=armature,quantity=force,method=weighted}}
\femmmeasure[armature pull,name=pull,parameter=3]
\end{codeexample}
\begin{key}{/femm/measure/name=\meta{identifier}}
Required result name, with the same syntax as selection names. A result stores
numerical values at measurement time, independently of later pictures or TeX
groups. Repeating the name replaces the snapshot. Changing a selection does not
retroactively change an existing result; measure again explicitly.
\end{key}
\begin{key}{/femm/measure/selection=\meta{name} (initially empty)}
Selected existing elements. Empty means the whole active model for integrals.
Weighted and Lorentz forces require an explicit named selection.
\end{key}
\begin{key}{/femm/measure/quantity=\meta{integrals or force} (initially integrals)}
|integrals| records area, volume, current, field averages, energy and coenergy.
|force| also records signed |Fx|, |Fy|, out-of-plane |torque| and |force| magnitude.
All forces require |field model=open|: artificial ideal-flux confinement does
not define the exterior air stresses needed for physical mechanical forces.
\end{key}
\begin{key}{/femm/measure/method=\meta{weighted, maxwell or lorentz} (initially weighted)}
Only used for |quantity=force|. |weighted| integrates air stress with an
auxiliary weight; |maxwell| integrates a supplied closed air contour;
|lorentz| integrates $\mathbf J\times\mathbf B$ over a nonmagnetic conductor.
Each method returns force and torque together.
\end{key}
\begin{key}{/femm/measure/profile=\meta{name} (initially empty)}
Required by |method=maxwell|. Use a closed, simple TikZ |femm/profile| with
zero offset. It must belong to this model and stay in source-free air, away
from material interfaces. Orientation is normalized automatically: reversing
the path does not reverse the outward force. Integration splits at triangle
crossings, independently of the profile's |samples| count. Specialized exact
gap scans are not accepted; capture a normal TikZ circle instead.
\end{key}
\begin{key}{/femm/measure/weighting=\meta{midgap, area or harmonic} (initially midgap)}
Only used by |weighted|. |harmonic| solves the ordinary Laplace weight.
|area| uses diffusion proportional to the square root of triangle area.
|midgap| uses that area-weighted solution and rounds nodal weights at one half,
emphasizing an interior layer of air elements. Compare methods and refine the
mesh near tight gaps; this choice is not an error estimate.
\end{key}
\begin{key}{/femm/measure/tolerance=\meta{number} (initially 1e-10)}
Relative residual tolerance of the auxiliary SPD weight solve, strictly between
zero and one. Distinct from the magnetic nonlinear tolerance and mesh accuracy.
\end{key}
\begin{key}{/femm/measure/origin x=\meta{number} (initially 0)}
\end{key}
\begin{key}{/femm/measure/origin y=\meta{number} (initially 0)}
Torque origin in model coordinates, not the force arrow's position. Positive
torque is counterclockwise, about the positive out-of-page axis. Force components
are in the model frame, even when the whole picture is rotated for presentation.
\end{key}
\begin{key}{/femm/measure/parameter=\meta{number} (initially 0)}
User-supplied abscissa retained with the snapshot, for instance gap in mm,
excitation in ampere-turns or angle in degrees. It does not change the physical
model; declare the corresponding model change before measuring.
\end{key}

\subsection{Numbers, units and PGFPlots}
\begin{command}{\femmresult\marg{name}\marg{component}}
Expandable unformatted scalar, suitable for numeric PGF calculations. Components
and units are listed below. Requesting a force from an integral-only result is
an error, not a silently returned zero.
\end{command}
\begin{center}\begin{tabular}{ll}\toprule
Component & Meaning and unit\\\midrule
|area|, |volume| & Selected area (m$^2$), volume (m$^3$)\\
|current| & Signed $\int J_z\,dS$ (A), without multiplying by depth\\
|energy|, |coenergy| & Volume integrals (J)\\
|Bx|, |By|, |B| & Area averages of components and magnitude (T)\\
|Hx|, |Hy|, |H| & Area averages of components and magnitude (A/m)\\
|Fx|, |Fy|, |force| & Force components and magnitude (N)\\
|torque| & Torque about the requested origin (N\,m)\\
|x|, |y| & Selected area centroid (model units)\\
|elements|, |parameter| & Element count, user abscissa\\\bottomrule
\end{tabular}\end{center}
The average of |B| is not generally the magnitude of the averaged vector.
Area averages and volume averages coincide because planar depth is constant.

This small affine-field example has an analytical energy reference:
$W=B^2 Sd/(2\mu_0)$. The boundary's coefficient uses SI coordinates, while
the picture dimensions and depth use millimetres.
\begin{codeexample}[width=5cm]
\begin{tikzpicture}[femm/problem={
  xmin=-10,xmax=10,ymin=-5,ymax=5,
  mesh size=2,depth=5},
  femm/boundary={side=all,
    type=dirichlet,y coefficient=.2}]
  \draw[fill=blue!5]
    (-10,-5) rectangle (10,5);
  \pic[femm/lines=7] {femm field};
  \femmmeasure[name=uniform-measure]
  \node[align=center] at (0,-13) {
    $B=$\femmvalue[component=B]
      {uniform-measure}\\
    $W=$\femmvalue[
      number format={sci,precision=3}]
      {uniform-measure}};
\end{tikzpicture}
\end{codeexample}

\begin{command}{\femmvalue\oarg{options}\marg{name}}
Print a scalar using PGF number formatting and its unit, in text or math mode.
For example |\femmvalue[component=Fy]{pull}|. Options belong to |/femm/value|.
\end{command}
\begin{key}{/femm/value/component=\meta{component} (initially energy)}
Select any public numeric component in the table.
\end{key}
\begin{key}{/femm/value/number format=\marg{PGF number options} (initially precision=3)}
Pass options to |/pgf/number format|, locally. Use |sci| or |fixed relative|
when appropriate for the magnitude and desired precision.
\end{key}
\begin{key}{/femm/value/unit=\meta{text} (initially auto)}
|auto| uses the listed SI unit, empty suppresses it, and another value supplies
upright math unit text. Overriding the label does not convert the number.
Centroids and user parameters have no automatic unit label.
\end{key}
\begin{command}{\femmresultplot\oarg{options}}
Emit a complete |\addplot+| in a PGFPlots axis. It retains the normal cycle list
and legend behaviour. No solve occurs: every named snapshot must already exist.
\end{command}
\begin{key}{/femm/result plot/results=\marg{comma-separated names}}
Required list, plotted in the given order. Results can come from different
pictures. No implicit sorting or interpolation of physical solutions occurs.
\end{key}
\begin{key}{/femm/result plot/component=\meta{component} (initially force)}
Ordinate, from the public component table.
\end{key}
\begin{key}{/femm/result plot/abscissa=\meta{component} (initially parameter)}
Abscissa, usually the parameter saved at measurement time.
\end{key}
\begin{key}{/femm/result plot/plot options=\marg{PGFPlots options} (initially empty)}
Ordinary |\addplot+| options. Axis labels, unit scaling, legends, logarithmic
axes and grids remain normal PGFPlots settings. All plot options reset per call.
\end{key}

\subsection{Force glyphs and styles}
\begin{key}{/tikz/pics/femm force}
Draw the force snapshot in its original problem picture. It cannot be attached
to another model. A zero force draws no arrow. The pic uses the saved model frame
so local scope transformations do not detach it from the physical geometry.
\end{key}
\begin{key}{/tikz/femm/force=\marg{options}}
Set glyph options in |/femm/force|, normally on the pic itself.
\end{key}
\begin{key}{/femm/force/result=\meta{name}}
Required force snapshot.
\end{key}
\begin{key}{/femm/force/length=\meta{positive number} (initially 12)}
Illustrative arrow length in model units; it is not proportional to force.
\end{key}
\begin{key}{/femm/force/scale=\meta{nonnegative number} (initially 0)}
Zero uses |length|. A positive value specifies model units per newton and
overrides |length|, allowing comparable arrow lengths across several bodies.
\end{key}
\begin{key}{/femm/force/x=\meta{number} (initially empty)}
\end{key}
\begin{key}{/femm/force/y=\meta{number} (initially empty)}
Arrow origin, defaulting independently to the selected area's centroid.
Changing it does not change the torque origin or force computation.
\end{key}
\begin{key}{/femm/force/label=\meta{magnitude or none} (initially magnitude)}
Print the force magnitude in newtons at the arrow tip, or omit the label node.
\end{key}
\begin{key}{/femm/force/number format=\marg{PGF number options} (initially precision=3)}
Number formatting for the magnitude label. For example |{sci,precision=2}|.
\end{key}
\begin{stylekey}{/tikz/femm force arrow}
Orange-dark arrow with a Stealth tip, 1pt line width. Customize with normal
|.style| or |.append style| handlers.
\end{stylekey}
\begin{stylekey}{/tikz/femm force label}
Small black label at the tip, on a lightly opaque white rounded background.
Anchor, position, font, opacity and colours are normal TikZ node options.
\end{stylekey}

\subsection{Numerical conventions and limitations}
For the isotropic law $\mathbf H=H_0(B)\mathbf B/B-H_c\mathbf m$, we use
\[
 w(\mathbf B)=\int_0^B H_0(b)\,db-H_c\mathbf m\cdot\mathbf B,
 \qquad w'=\mathbf B\cdot\mathbf H-w.
\]
The primitive is integrated exactly on each linear or cubic B--H interval,
including its extrapolation. This is a declared variational convention with
$w(\mathbf0)=0$; permanent-magnet energy can be negative. It is not an absolute
chemical/internal energy of the magnet, and is not a claim of FEMM's magnet
energy reference convention. Nonlinear coenergy is not generally equal to
energy. Circuit flux linkage and secant inductance are defined in the static
circuit chapter; differential responses are defined in Section~\ref{sec:operating}.
Do not infer a general inductance from $2W/I^2$ in a biased nonlinear model.

Maxwell stress uses $\mathbf T=(\mathbf B\mathbf B^{\mathsf T}
-B^2\mathbf I/2)/\mu_0$. Contour force is $d\oint\mathbf T\mathbf n\,dl$,
where $d$ is physical depth and $\mathbf n$ is outward. Torque uses the signed
lever about the requested origin. Integration is exact for the finite-element
field on each straight contour piece; geometric and field errors remain.

Weighted force is $-d\int\mathbf T\nabla w\,dS$, with weight one on selected
elements, zero on exposed domain boundaries and other material/source regions.
The remaining weight is solved in source-free air. A selected body touching an
unselected material/source or the domain boundary is rejected. Lorentz force
is restricted to $\mu_r=1$, nonmagnetic, noncoercive selected elements; it does
not include magnetization forces on iron or permanent magnets.

Validate force by mesh/domain refinement and by comparing independent methods.
The optional |python3 scripts/post-interop.py| compares nonlinear U-core regional
energy/coenergy and weighted force/torque with xfemm reading the same ANS mesh.
Analytical conductor/dipole tests and an independent fixed-current coenergy
finite difference exercise signs and physical consistency. These checks do not
provide a universal force error bound.

The Lua module |luafemm-post| exposes these entry points:
\begin{verbatim}
select(model, name, options)
capture(model, name, softpath, options)
measure(model, options)
value(name, component)
coordinates(names, component, abscissa)
\end{verbatim}
Lua selection contours and torque origins use model units; output units match
the table. Result and selection registries persist for the document lifetime.

\subsection{Total rotor stress in a moving air band}
The additional choice |method=band| integrates a prescribed P1 weight across a
captured circular air band. It requires |quantity=force| and a sliding-position
model, and accepts no selection or profile. Torque retains the ordinary origin
keys. See Section~\ref{sec:sliding} for all defaults, the integral and validation.

% SPDX-License-Identifier: LPPL-1.3c
% Copyright (C) 2026 Christophe Jorssen
% Author and Current Maintainer: Christophe Jorssen <christophe.jorssen@gmail.com>
% LPPL maintenance status: maintained
% This file is part of luafemm. See LICENSE and LICENSES.md for its terms.
\section{FEMM problem and solution interchange}
\label{sec:interchange}
The interchange described here applies to |field model=open|. Ideal-mode
floating-window constraints cannot be represented by this subset; FEM/ANS
export and FEM import in ideal mode report an error.

A drawing can now travel between your TeX document and a magnetic FEMM problem.
There are three distinct operations: read a preprocessor |.fem| file, write a
problem for another solver, and write the solution already computed by Lua as
an |.ans| file. None requires FEMM, xfemm or a native library at document runtime.
The optional xfemm tools used to validate this implementation are developer tools.
This chapter describes the supported planar, real, static subset of magnetic
format~4.0. It does not promise that every physical feature of FEMM is implemented.

\subsection{First, keep drawing curves}

Imagine that one side of a conductor is a Bezier curve and another conductor
is a rotated ellipse. We need not replace these paths by hand-written circular
arcs. PGF evaluates the drawing; luafemm approximates its curves by short straight
segments before generating finite elements. This already happens in ordinary
luafemm calculations. The FEM writer uses those computational interfaces.

There are therefore three geometric representations to keep in mind:
\begin{enumerate}
\item The evaluated source path contains the PGF lines and cubic curves, or
      the segments and circular arcs read from a FEM file. Native captured
      soft-path data and their coordinate frame are retained in the model.
\item The computational geometry is polygonal. Its faces carry materials and
      excitations, and its edges carry boundary conditions. It is independent
      of the interior triangle spacing.
\item The solution consists of the actual triangles and nodal potential $A_z$.
      ANS export preserves these data without remeshing or projection.
\end{enumerate}

The source TeX program, node anchors, text, colours and macros are not encoded
in a FEM file. A round trip preserves the represented physical problem, not the
original drawing program. Similarly, an imported circular arc may be preserved
by the passive Lua document writer, whereas the computational exporter writes
its polygonal approximation. No circular-arc fitting is performed on arbitrary
TikZ curves.

Start with a modest |curve tolerance| and |mesh size|. The first controls curve
approximation, the second triangle spacing. Reducing the triangle size cannot
recover curvature already removed by polygonalisation. As the drawing approaches
its final form, reduce each independently and compare the field on the contour
of interest. A narrow curved air gap deserves particular attention: its width
must be resolved geometrically before refining triangles inside it.
The chord tolerance is not a global error estimate for the magnetic field;
PGF coordinate precision, snapping and constructed intersections also matter.

Here is a complete example. The dotted interfaces expose the polygonal geometry;
the last drawing comes from importing the exported FEM problem and solving again.
The graph still refers to the original named profile.
\luafemmexample{interchange-curves}

\subsection{Reading a preprocessor problem}

\begin{key}{/tikz/femm/import=\marg{options}}
Create a model from a magnetic FEM file at the beginning of the picture, after
picture options have been collected. The option family is |/femm/import|.
This key installs the same millimetre coordinate convention as |femm/problem|.
It can be combined with |femm/problem| to override numerical settings explicitly.
Importing does not itself draw, mesh or solve anything.
\end{key}
\begin{key}{/femm/import/file=\meta{filename} (initially empty)}
Required input filename, including |.fem|. The bridge first uses TeX's normal
file search, then checks the output directory for a file generated earlier in
the same document, and finally uses the supplied path. The passive Lua API uses
ordinary working-directory paths. Imported content is never executed as Lua or TeX.
\end{key}

Only explicitly set problem keys override imported settings. In particular,
unspecified luafemm defaults do not replace the file's depth or solver precision.
The domain bounds come from the imported geometry; explicitly supplying |xmin|,
|xmax|, |ymin| or |ymax| together with import is an error. Source length units are
converted directly into model millimetres without a TikZ drawing/recapture step.
The Lua API can instead request another model unit.

The initial imported mesh size is one twelfth of the largest bounding-box extent.
A block-label mesh control supplies a local spacing hint in converted model units.
FEMM's diameter-like area control, minimum-angle setting and smart-mesh heuristics
are retained in passive metadata but do not imply identical meshing algorithms.
Use the documented luafemm mesh and quality keys explicitly for the Lua solve.

An imported model's physical topology is immutable in this version. Add ordinary
TikZ annotations, measurement contours and plots freely, but edit physical
boundaries/material regions in the FEM source before reimporting. Additional
|femm/region| declarations and rectangular side-condition overrides are rejected
rather than being silently ignored. Native TikZ models retain their existing
paths, components, painting precedence and fill-rule interface.

\subsection{Showing the computational geometry}

\begin{key}{/tikz/pics/femm geometry}
Draw the compiled physical interfaces of the current model. This picture does
not register another material region and does not start a magnetic solve.
For a native TikZ model it displays the polygonal interfaces used by the
computational exporter, including the exterior rectangle. For an imported model
it includes the boundaries of true exclusions.
\end{key}
\begin{stylekey}{/tikz/femm geometry}
Base drawing style, initially |draw=black!65,line width=.3pt|.
\end{stylekey}
\begin{key}{/tikz/femm/geometry style=\marg{TikZ options}}
Set |femm current geometry|, applied after the base style. Use this to inspect
interfaces in another colour or with a dotted line without changing the physics.
\end{key}
\begin{stylekey}{/tikz/femm current geometry}
Additional geometry style, initially empty.
\end{stylekey}
\begin{stylekey}{/tikz/femm group/\meta{integer}}
Optional style for an imported segment group. It is tried after the base and
current geometry styles. For example, |femm group/3/.style={red,thick}| highlights
group~3. Undeclared group styles have no effect. Imported names are never turned
into executable TeX control sequences.
\end{stylekey}

\subsection{Writing a problem or its existing solution}

\begin{key}{/tikz/femm/export=\marg{options}}
Append an export request to the picture's end hook. Repeated keys produce repeated
exports, bound to that picture's model. All physical declarations precede export.
A FEM request compiles geometry but does not mesh or solve. An ANS request requires
a converged solution at that point: use a field picture or |\femmsolve| earlier.
A missing solution is an error, not an implicit request for an expensive solve.
\end{key}
\begin{key}{/femm/export/format=\meta{choice} (initially fem)}
A PGF |.is choice| selector with choices |fem| and |ans|. FEM writes a magnetic
problem, ANS writes that problem followed by the existing solution.
\end{key}
\begin{key}{/femm/export/file=\meta{filename} (initially empty)}
Required destination, including its extension. Relative paths follow TeX's
|--output-directory|, as cache files do. Parent directories must already exist.
The writer validates and serializes first, writes a temporary file beside the
destination, then replaces the destination after success. Use a different filename
from the imported source. Repeated exports to the same destination replace it.
\end{key}
\begin{key}{/femm/export/model=\meta{name} (initially empty)}
Select the model for the explicit export command. Empty means the current Lua
model. Picture-level export supplies its bound model automatically, so the option
is intended for delayed command use. Named models remain available after their
picture has ended.
\end{key}
\begin{command}{\femmexport\opt{\oarg{options}}}
Immediately export using the same |/femm/export| family. This is useful after a
profile outside the source picture has triggered its solve. Use |model| explicitly
when subsequent pictures have created other models. It never starts a solve.
\end{command}
\begin{key}{/femm/problem/name=\meta{identifier} (initially empty)}
Register the picture model under a unique name for delayed export. Duplicate
names are errors. Prefer descriptive, nonnumeric names; internal picture IDs
are reserved. Naming does not change the numerical problem or cache identity.
\end{key}
\begin{key}{/femm/problem/depth=\meta{positive number} (initially unset)}
Planar thickness in model millimetres. A native model must specify it before
FEM or ANS export. Imported depth is converted from the declared file unit.
Depth does not change the planar $A_z$ solve, but it matters to FEMM's integrated
quantities. It is included in the solution export signature.
\end{key}
\begin{key}{/femm/problem/interpolation=\meta{choice} (initially linear)}
Choose |linear|, the existing piecewise-linear $H(B)$ law with vacuum-slope
extrapolation, or |femm|, a natural cubic law with FEMM-compatible smoothing and
endpoint-slope extrapolation. Imported nonlinear materials default to |femm|
unless this key explicitly overrides it. Set it before declaring materials.
Linear materials are unaffected.
\end{key}

Nonlinear ANS export requires |interpolation=femm| for every nonlinear material
used by the mesh. The postprocessor reconstructs $B$ from nodal $A$, then computes
$H$ from its own material law; preserving the B--H sample table alone is insufficient.
The writer therefore rejects a nonlinear solution computed with the old law.
FEM export remains available with either policy, but a subsequent FEMM solve
uses FEMM's constitutive law and can produce a different solution.

The static ANS writer includes zero-based element-to-label references, one current
record per block label, periodic nodal constraints when present, and an empty
air-gap trailer. Current density is encoded in effective material records or
in the associated circuit records. Spatially varying coercive directions use
elementwise blocks to preserve the solved quadrature-averaged source; use the
separate FEM file to edit the original angle law. The output uses the
non-incremental four-column element dialect read by the audited postprocessors.
It contains the actual converged mesh and physical potential in tesla metres.
It is not a replacement for the private |.lfc| cache and cannot currently be
imported as a solved luafemm model.

\subsection{The order of operations}

\begin{enumerate}
\item Collect picture options and create or import exactly one model. Apply only
      explicit overrides to imported numerical settings.
\item Declare the native materials and physical paths, or draw a view of imported
      geometry. Declare native boundary conditions before meshing.
\item Compile geometry and mesh when a solve or mesh picture requires them.
      A FEM-only export needs neither a mesh nor a solution.
\item Solve explicitly or through a field picture/profile. Capture measurement
      contours without turning them into physical interfaces.
\item Run the queued exports at picture end. If solving occurs later, use a
      named |\femmexport| command after that solve.
\end{enumerate}

Changing a source path, geometry tolerance, material or boundary requires a new
model and solve. The writer rejects a stale solution whose input signature has
changed. Cache format~5 validates geometry, physical settings and carved-domain
coverage; imported solutions restored from a matching cache can be exported.
Drawing scales, styles and profile sampling do not invalidate the physical cache.
Neither a FEM nor an ANS file reconstructs the original TeX program.

\subsection{Two U-shaped devices, in both directions}

The coil example uses the familiar node component. We add a depth and select the
FEMM material law before declaring the component, then request both exports at
picture end. The final drawing imports the produced FEM file and solves again.
The two coil regions are exported with their effective current densities. Their
original separate turn count and current cannot be recovered from an $NI$ value.
\luafemmexample{interchange-coil}

The permanent-magnet example follows the same sequence. The alnico region carries
a constant magnetisation direction; the surrounding iron remains nonlinear.
The field profiles belong to the original model even after the final imported
picture creates another model.
\luafemmexample{interchange-alnico}

\subsection{Supported features and explicit limits}

The reader accepts LF/CRLF, case-insensitive section keywords, quoted names,
optional label expressions as passive data, and all six FEMM length units.
Names remain opaque byte strings in the Lua codec: no Windows code-page to UTF-8
conversion is attempted. Numerical geometry rendering does not typeset those
names. The writer rejects quotes/newlines that cannot be represented safely in a
quoted scalar. Malformed input is diagnosed with filename and line number.

The solver adapter supports closed planar graphs, circular arcs up to $180^\circ$
per record, nested and disconnected regions, default labels, diagonal linear
anisotropy, isotropic B--H materials, linear directional laminations, constant
or supported angular coercive directions, prescribed real current density,
series circuits, paired straight periodic/antiperiodic segments,
tagged prescribed-potential and constant real mixed conditions, and natural
boundaries. No Mesh seeds exclude their faces entirely. Each
unanchored planar connected component receives its own gauge and current-balance check.

Axisymmetric static input/output is supported as specified in
Section~\ref{sec:axisymmetric}; axisymmetric periodicity and exterior
transformations remain excluded. The adapter rejects AC, previous-solution modes, nonlinear
anisotropy, unsupported lamination/fill physics, FEMM parallel-circuit input,
point properties, arbitrary executable magnetisation expressions, periodic
circular-arc input, sliding-gap boundaries, negative
Robin coefficients, and mixed conditions on internal interfaces. Open or isolated
constraints that do not bound valid faces also receive diagnostics. Default labels
with circuits, conflicting labels and labels on interfaces are rejected.
Unspecified FEM boundaries retain natural conditions; they do not inherit the
native rectangle's default prescribed potential.

Exact export rejects arbitrary source/boundary callbacks and affine mixed data
that a constant FEMM boundary property cannot represent. Segment groups are
retained for drawing. The computational export uses effective material records
per face. Named DC windings export as series records with resolved currents;
parallel connections and voltage control are reduced to the operating point.
Rotated anisotropic tensors with cross terms cannot be exported. Supported
angular laws export as passive expressions; only that restricted grammar and
the presets |theta| and |theta+90| are accepted on import. Unsupported legacy
ampere-turn inputs still export as current density: the writer does not invent
separate turns and current or recover every preprocessor editing attribute. For an unmodified passive document,
|luafemm-fem.serialize| retains its original property and arc records separately.

\subsection{Lua interface and independent validation}

The principal Lua entry points are:
\begin{itemize}
\item |luafemm.import_fem(file, options)| creates an unsolved model;
\item |luafemm.export_fem(model, file)| writes its computational geometry;
\item |luafemm.export_ans(model, file)| writes its converged solution.
\end{itemize}
The passive codec provides |read|, |parse|, |serialize| and |model|; parsing alone
does not authorize unsupported physics for solving. Imported expressions are
never passed to |load|, |dofile| or a TeX input command. The existing ownership
rule applies: imported documents and material tables are copied into the model.

Original small fixtures test units, indexing, malformed input, arc subdivision,
exclusions, labels, series-current normalization, cubic derivatives and cache
restoration. A separate optional C++ client loads emitted ANS files from disk
through xfemm's FPProc. It compares $A$, unsmoothed $B$ and $H$ with Lua samples.
Independent xfemm solves of the exported U devices are compared at two Lua mesh
resolutions; different meshers are not expected to produce identical triangles.
Run |python3 scripts/interop.py --build --xfemm ../xfemm| to reproduce these checks.
The audited revision and local compiler are recorded with the generated reports.
Windows FEMM GUI acceptance remains unverified on the macOS development host.

\subsection{Snapshots of a moving rotor}
A sliding-position export is an ordinary static problem at the current angle.
FEMM can remesh its |.fem|; the |.ans| preserves the solved P1 mesh. No native AGE
or motion-controller record is emitted. Existing tensor and nonlinear-law
restrictions remain in force. See Section~\ref{sec:sliding}.

% SPDX-License-Identifier: LPPL-1.3c
% Copyright (C) 2026 Christophe Jorssen
% Author and Current Maintainer: Christophe Jorssen <christophe.jorssen@gmail.com>
% LPPL maintenance status: maintained
% This file is part of luafemm. See LICENSE and LICENSES.md for its terms.
\section{Numerical method, verification and limits}
\label{sec:numerics}
The transverse axial-H diffusion formulation and its Bessel benchmarks are
detailed in Section~\ref{sec:currents}; it uses a different scalar unknown
from the vector-potential formulations below.

\subsection{Planar finite elements}
The domain is invariant along $z$. In SI units,
$\mathbf A=A_z\mathbf e_z$ and
$\mathbf B=(\partial_y A_z,-\partial_x A_z)$.
For test functions $v$ vanishing on the Dirichlet boundary, the weak problem is
\[
 \int_\Omega \nu\nabla A_z\cdot\nabla v+\int_{\Gamma_w}(cA_z+g)v
 =\int_\Omega J_zv+
  \int_\Omega H_c\mathbf m\cdot(\partial_yv,-\partial_xv).
\]
Here $\Gamma_w$ contains Neumann and Robin sides, with $c=0$ for Neumann.
The boundary tangent is counterclockwise; $H_t=cA_z+g$. On an edge of
length $\ell$, the boundary mass matrix is
$\ell\bigl(\begin{smallmatrix}2&1\\1&2\end{smallmatrix}\bigr)/6$.
It integrates both $cA_z$ and the affine data $g$ exactly. These terms
are added to the residual and, for $cA_z$, its tangent. Corner constraints
and the all-Neumann compatibility check are described in
Section~\ref{sec:boundaries}.
The potential uses continuous piecewise-affine functions on triangles.
Material laws and sources are constant within each element. With basis
functions $N_i$, the element residual is
\[
 R_i=|T|\left[\nu\nabla A_z\cdot\nabla N_i
 -J_z/3-H_c(m_x\partial_yN_i-m_y\partial_xN_i)\right].
\]
The differential tangent follows from the derivative of the selected
$H(B)$ law: piecewise linear by default, or the FEMM-compatible cubic policy
described in Section~\ref{sec:interchange}. The code scales the entire system by $\mu_0$ for numerical
convenience; this does not change the solution. Dirichlet degrees of
freedom are eliminated. Newton steps use backtracking to decrease the
residual. Linear systems use conjugate gradients and incomplete Cholesky
IC(0); an unsuccessful factorisation is restarted with a positive diagonal
shift of the preconditioner, without changing the system being solved.

The default boundary $A_z=0$ represents a finite truncation of the
exterior. For physical depth $d$, potential levels separated by $\Delta A_z$
enclose flux $d\,\Delta A_z$. Equal level spacing therefore represents equal
flux spacing, not equal geometric separation.

\subsection{Constrained meshing and predicates}
The mesher splits crossings, T junctions and collinear overlaps of the
polygonal constraints. It samples their edges and seeds the interior
with triangular lattices, using the finest local size where requested.
Incremental Delaunay insertion is followed by constrained-segment cavity
retriangulation and Lawson edge flips.

Topological decisions use orientation and incircle predicates. A
floating-point filter certifies well-separated signs; uncertain cases
are recomputed using error-free sums/products and floating-point
expansions in Lua. The result has the exact sign of the determinant of
the supplied binary coordinates, including exact zero, provided arithmetic
neither overflows nor underflows. Normalising coordinates before meshing
keeps the intended working range moderate. This does not make the input
geometry exact before its TeX rounding, nor make intersection construction
exact: intersection coordinates still use floating-point arithmetic.

The native staging-mesh audit checks positive areas, rectangular coverage, preservation
of every constraint and vertex, at most two triangles per edge, exposed
edges only on the tagged exterior, Euler's relation $V-E+F=1$, and local
Delaunay legality of every unconstrained flippable edge. Failure stops
the calculation. The vertex limit is 40000 vertices for Delaunay and
150000 grid vertices.

Triangle was studied as an algorithmic reference; its source and binary
are not embedded. This is not a full Triangle port. Optional segment-first
Delaunay refinement enforces angle and area targets, or raises an error
when its budget or geometric precision is exhausted (Section~\ref{sec:mesh-quality}).
There is no specialised acute-corner protection or guaranteed smooth
transition between mesh sizes. With both quality targets disabled, narrow
features and acute corners can create badly conditioned triangles.
Monitor |min_angle| and perform field-convergence studies; triangle quality
alone is not a physical error estimate.

\subsection{What is tested}
The release checks analytic affine fields for Dirichlet, Neumann, Robin
and mixed sides on both meshers (including nonlinear permanent magnets),
manufactured Poisson and Robin convergence, all-Neumann compatibility,
nonlinear saturation, current and magnetisation reversal, an exact
magnet/air interface patch, and a magnetised disk against the infinite-domain
solution. Geometry tests cover intersections, shared edges, small gaps,
concavity, scaling, translation and compound-region fill rules. Circular-shell
shielding is checked against an analytic cavity field with exact exterior
Dirichlet data. An independent Python rational oracle
checks 4530 orientation, incircle and diametral-disk signs; 49 stress
configurations exercise the base mesher. Quality tests independently check
angles and areas, scaled and translated inputs, compound cavities, refined
boundary tags, affine Robin fields, cache reuse and explicit budget failures.
Python is a development-test dependency only.

All 246 catalogue records are validated. Profiles are tested for Cartesian
and local components, arc length, reversed traversal, normal offsets,
out-of-domain policy and retained model lifetime. The same nonlinear
scene is compiled under plain LuaTeX, LuaLaTeX and ConTeXt MkIV; computed
sample values and node counts must agree within $10^{-10}$. Load guards,
catcode restoration, transformed paths and PGFPlots legends are exercised.

For the original mixed alnico/illustrative-iron U with 1 mm gaps, sample
results document sensitivity rather than absolute accuracy:
\begin{center}\small
\begin{tabular}{rrrrr}
\toprule Global/gap step (mm) & Domain (mm$^2$) & Nodes & Left $B_y$ (T) & Right $B_y$ (T)\\
\midrule
2 / .5 & $200\times185$ & 11376 & .782876 & -.782907\\
1.5 / .25 & $200\times185$ & 20409 & .777974 & -.778597\\
2 / .5 & $280\times265$ & 22293 & .775140 & -.775140\\
\bottomrule
\end{tabular}\end{center}
The refinements change the gap field by about 0.6 percent, and the larger
domain by about one percent. Material-model error is a separate issue.
The interchange suite additionally loads Lua-written ANS files through xfemm's
independent disk postprocessor and compares $A$, unsmoothed $B$ and $H$.
Independent solves of the exported U-coil and alnico models agree to within
2 percent in the tested gap-field samples at the two draft mesh resolutions.
This is a bounded validation result, not a universal accuracy guarantee;
see Section~\ref{sec:interchange} and the reproducible reports in |build/interop|.

Imported domains are carved from a constrained staging mesh. Their audit also
checks active area, protected exposed edges and the Euler characteristic of the
active cell complex; shared child boundaries are counted as a union. Exclusions
are never silently replaced by air. Unanchored disconnected components have
independent gauges and current-balance checks.

\subsection{Supported physics and practical diagnostics}
The static solver is planar or axisymmetric, with a linear symmetric permeability
tensor or a monotone tabulated law and prescribed permanent-magnet direction.
Section~\ref{sec:axisymmetric} gives the separate cylindrical formulation.
The separate linear harmonic solver in Section~\ref{sec:harmonic} adds
single-frequency diffusion and eddy currents. Section~\ref{sec:biased-harmonic}
extends it to a fixed differential tensor around a nonlinear DC bias.
General time stepping, hysteresis
history, irreversible demagnetisation and analytical sliding air-gap elements
are absent. Circular P1 moving bands are described in Section~\ref{sec:sliding}. Forces and circular
open-boundary approximations are available as described in their chapters.

If Newton fails, inspect the material data, extreme element aspect ratios
and current scale before raising |max iterations|. For a lost or collapsed
geometric feature, compare its size with PGF capture precision and the
merging tolerance. For a graph with jumps, inspect the elementwise field
and sampling side before requesting graphical smoothing. The log identifies
the failing stage; a solver failure is not turned into a plausible plot.

\subsection{Primary references}
The development was based on studying the source code of FEMM/xfemm
and Triangle, then reimplementing selected formulations and algorithms
in Lua/LuaTeX. No native solver or mesher is embedded. The following primary references
also describe the surrounding algorithms and interfaces:
\begin{itemize}
\item \href{https://www.femm.info/wiki/Documentation}{FEMM documentation},
  planar magnetostatics and material conventions; FEMM revision
  \nolinkurl{7d9e8ed90772ddb939a3256f976e231bcb5cca97}.
\item xfemm revision \nolinkurl{34ac6cc586a346066984526620dfb67383c8e3c5},
  \nolinkurl{cfemm/fsolver/static2d.cpp} and
  \nolinkurl{cfemm/libfemm/CMaterialProp.cpp}.
\item Jonathan Shewchuk,
  \href{https://www.cs.cmu.edu/~quake/robust.html}{Robust Adaptive Floating-Point Geometric Predicates}
  and \href{https://www.cs.cmu.edu/~quake/triangle.research.html}{Triangle algorithm references}.
\item The \href{https://tikz.dev}{PGF/TikZ manual} and the installed
  PGFPlots manual for path transformations, key handlers and axes.
\end{itemize}

\subsection{Circular moving-band operator}
Section~\ref{sec:sliding} describes retained bulk meshes, annular reconnection
and prescribed-weight Maxwell stress. The operator remains planar P1 at each
static angle; exact orientations and mesh coverage are checked before solving.
It does not reuse the previous position's nonlinear field as a frozen material
law. Coenergy differences and mesh/angle sensitivity remain necessary checks.

% SPDX-License-Identifier: LPPL-1.3c
% Copyright (C) 2026 Christophe Jorssen
% Author and Current Maintainer: Christophe Jorssen <christophe.jorssen@gmail.com>
% LPPL maintenance status: maintained
% This file is part of luafemm. See LICENSE and LICENSES.md for its terms.
\section{Release scope and convergence evidence}
\label{sec:status}
Version \luafemmversion{} adds axisymmetry, linear harmonic diffusion and DC-biased
operating-point responses and static moving-band positions to planar modelling and quantitative postprocessing. The following matrix distinguishes implemented capabilities from
approximations and unsupported models. ``Supported'' always refers to the stated
scope; it does not certify accuracy for an arbitrary user geometry.
The repository's \nolinkurl{docs/capabilities.md} provides the detailed matrix,
including individual test sources and the separate FEMM exchange restrictions.

\begin{longtable}{p{.24\linewidth}p{.13\linewidth}p{.53\linewidth}}
\toprule Capability & Status & Scope or restriction\\
\midrule\endhead
Generic formats & Supported & plain LuaTeX, LuaLaTeX and ConTeXt MkIV; pure Lua runtime.\\
TikZ geometry & Supported & Arbitrary paths, compound domains and parametric nodes; curves are flattened with a declared tolerance.\\
Planar nonlinear fields & Supported & Isotropic, single-valued B--H curves and permanent magnets.\\
Linear tensors & Supported & Rotated anisotropy and directional linear sheet mixtures.\\
Spatial magnetisation & Supported & Uniform, radial, tangential and angular laws.\\
Named frontiers & Supported & Dirichlet, Neumann, nonnegative Robin, periodic and antiperiodic paired traces.\\
Open exterior & Partial & Circular multipole Robin approximation; zero net longitudinal current and air at the boundary.\\
DC windings & Supported & Named series/parallel resistive branches with current or voltage drive.\\
Circuit measurements & Partial & Resistance, power, closed-winding flux linkage and supported secant inductance; incremental self and mutual responses; no automatic matrix.\\
Ideal circuits & Supported & Confined flux, junction conservation and mean contours; separate physical model.\\
Field plots & Supported & Oriented lines, component profiles, angular scans and spatial Fourier spectra.\\
Regional measurements & Supported & Area, volume, current, B/H averages, energy and coenergy.\\
Forces and torques & Partial & Planar Maxwell and conductor Lorentz methods, subject to selection and mesh assumptions.\\
Caches & Supported & Passive format 13; cold, warm and frozen reuse across the three formats.\\
FEM/ANS exchange & Partial & Documented planar/axisymmetric DC subset; ANS writing only, polygonal export and representability checks.\\
Axisymmetry & Partial & Static meridional formulation, full-revolution integrals and supported FEM/ANS exchange; see Section~\ref{sec:axisymmetric}.\\
Harmonic fields & Partial & Linear planar/axisymmetric diffusion, phased currents, snapshots and Joule losses; Section~\ref{sec:harmonic}.\\
Harmonic materials & Supported subset & Passive complex tensors, insulated-sheet homogenisation and separate intrinsic/sheet losses; Section~\ref{sec:losses}.\\
Harmonic circuits & Supported subset & Coupled current/voltage ports, series/parallel branches and solid-current constraints; Section~\ref{sec:ac-circuits}.\\
Eddy-current loops & Supported subset & Axial diffusion, normal glyphs and transverse axisymmetric sections; Section~\ref{sec:currents}.\\
Operating points & Partial & Static frozen/incremental and small harmonic responses on the exact bias mesh; Sections~\ref{sec:operating}--\ref{sec:biased-harmonic}.\\
Sliding air bands & Partial & Complete circular planar DC P1 band, retained bulk mesh, torque/linkage; Section~\ref{sec:sliding}.\\
Advanced magnetics & Deferred & Nonlinear harmonic laws, nonlinear anisotropy and periodic sliding gaps.\\
Additional rendering & Deferred & Scalar colour maps and general B/H vector arrows.\\
Other physics & Deferred & Electrostatics, electrical conduction and thermal solvers.\\
\bottomrule
\end{longtable}

The material catalogue is larger than the supported physics. A stored loss,
frequency or wire property does not activate its associated material model.
Positive |frequency| explicitly selects the supported linear harmonic solver. Similarly,
a spatial Fourier spectrum of a winding field is not a time-harmonic solve.
Native rotated tensors can be solved, but FEMM export rejects model-frame
cross terms. Native parallel and voltage-driven circuits export their resolved
operating currents, not an editable equivalent circuit controller.

\subsection{Three separate sensitivity questions}
Run |make static-convergence| from the source checkout. It writes
\nolinkurl{build/convergence/static.csv}; no native program is needed.
The study asks whether refining the field mesh reduces a known error, whether
moving an artificial boundary changes an observable, and whether a nonlinear
force stabilises. These questions require different comparisons.

For the manufactured potential $A_z=\sin(\pi x)$ on the unit square,
$\mathbf B=(0,-\pi\cos(\pi x))$. Antiperiodic left/right traces and homogeneous
Neumann top/bottom data exercise the signed constraints. Three-point triangle
quadrature estimates the squared field error, including variation inside each
triangle. The relative $L^2$ errors are:
\begin{center}
\begin{tabular}{rrr}
\toprule Spacing (m) & Triangles & Relative field error\\
\midrule
.16 & 110 & 12.823\%\\
.08 & 386 & 6.334\%\\
.04 & 1518 & 3.155\%\\
\bottomrule
\end{tabular}
\end{center}
Successive error ratios 2.02 and 2.01 agree with first-order convergence of
piecewise-constant P1 fields. The runner requires each ratio to exceed 1.6.

For the exterior study, a disk of radius .25 m centred at $(.4,0)$ m has
$\mu_r=1$ and $\mu_0H_c=1$ T. Its exact internal field in infinite air is
$B_x=.5$ T. The offset deliberately excites modes beyond the first circular
Robin multipole. The observable is the area-averaged internal $B_x$.
\begin{center}
\begin{tabular}{rrrr}
\toprule Outer radius (m) & Global spacing (m) & Mean $B_x$ (T) & Error\\
\midrule
.9 & .14 & .491457 & 1.709\%\\
1.5 & .14 & .484485 & 3.103\%\\
2.5 & .14 & .483158 & 3.368\%\\
1.5 & .07 & .495453 & .909\%\\
1.5 & .035 & .499307 & .139\%\\
\bottomrule
\end{tabular}
\end{center}
The radius sweep holds the 64-sided magnet polygon and its local spacing
.035 m fixed. Refinement reduces the global spacing, the local spacing
(one quarter of the global value) and the exterior polygon edge target;
the magnet polygon remains fixed. The small coarse domain appears more
accurate because truncation and discretisation errors partly cancel.
Enlarging it alone therefore does not establish accuracy. Refinement at fixed
radius reduces the combined error, but does not isolate the remaining
boundary error from polygonal geometry and finite-element errors.

Finally, a nonlinear U armature at 3000 ampere-turns, 20 mm depth and fixed
75 mm exterior half-width gives:
\begin{center}
\begin{tabular}{rrrr}
\toprule Spacing (mm) & Triangles & $F_y$ (N) & Successive change\\
\midrule
4 & 3374 & $-101.208792$ & ---\\
2 & 12768 & $-103.446281$ & 2.163\%\\
1.5 & 22144 & $-104.437018$ & .949\%\\
\bottomrule
\end{tabular}
\end{center}
The changes use the finer force as denominator. The final .949\% change
is sensitivity evidence, not a certified 1\% error bound. The first two
meshes also pass the independent xfemm postprocessor comparison for energy,
coenergy, force and torque. Agreement on the same mesh checks implementation;
it does not measure the mesh error.

\subsection{Release checks and remaining validation}
The release checks all three formats, cache reuse, the exact-predicate oracle,
the 44 complete manual source/output examples, and an extracted TDS
installation. Optional xfemm checks compare exported snapshots and independent
solves. Windows FEMM GUI acceptance remains unverified. The source repository
records the commands and measured tolerances in
\nolinkurl{doc/luatex/luafemm/validation.md}.

Old cache snapshots must be recomputed. Automatic mode rebuilds incompatible
results; frozen mode reports the incompatibility instead of guessing a
conversion. Keep the original TikZ source alongside any exported FEMM files:
file-format restrictions can prevent a lossless round trip of native semantics.

% SPDX-License-Identifier: LPPL-1.3c
% Copyright (C) 2026 Christophe Jorssen
% Author and Current Maintainer: Christophe Jorssen <christophe.jorssen@gmail.com>
% LPPL maintenance status: maintained
% This file is part of luafemm. See LICENSE and LICENSES.md for its terms.
\section{Lua API and documentation builds}
\label{sec:development}
\subsection{Module boundaries}
Every Lua file returns a module table and uses local implementation
variables; there is no deprecated |module()| call and no runtime change
to |package.path|. LuaTeX's or ConTeXt's TEXMF resolver finds installed
files. The TeX loader creates one deliberate Lua global, |luafemm|, to
hold the loaded module and current model. Numerical scripts can use
|local femm=require('luafemm')| without a running TeX interpreter.

\begin{description}
\item[luafemm.lua] Material laws, model ownership, assembly, Newton,
  preconditioning, field evaluation and contour topology.
\item[luafemm-axisymmetric.lua] Regular azimuthal potential, radial quadrature,
  cylindrical fields, flux contours and physical revolution integrals.
\item[luafemm-boundary.lua] Side/tagged conditions, corner constraints, component gauge
  compatibility and exact P1 boundary integrals.
\item[luafemm-frontiers.lua] Native domains, tagged paths, signed periodic
  constraints, trace synchronisation and circular exterior checks.
\item[luafemm-static.lua] Linear tensors, sheet mixtures, spatial magnetisation
  quadrature and named DC winding operating points and observables.
\item[luafemm-operating.lua] Owned static bias captures, same-mesh static frozen
  and static/harmonic incremental responses, material tangents and provenance checks.
\item[luafemm-static-exchange.lua] Representable static FEMM records, safe
  angular expressions, straight periodic pairs and resolved circuit exchange.
\item[luafemm-post.lua] Named selections, regional integrals, energy/coenergy,
  Maxwell and Lorentz forces and torques, result formatting and field weights.
\item[luafemm-ideal.lua] Confined-domain compilation, floating boundary
  degrees of freedom, exact section flux and component mean contours.
\item[luafemm-integrals.lua] Segment/triangle intersection and exact integration
  of the discrete field along profiles.
\item[luafemm-shapes.lua] E and toroidal geometries, tapered poles and magnet slices.
\item[luafemm-components.lua] Pure component geometry, physical anchors and
  area-based winding currents; no TeX or solver state.
\item[luafemm-fem.lua] Passive magnetic file codec, supported-physics checks,
  source metadata and the planar model adapter; computational problem export.
\item[luafemm-topology.lua] Split planar constraints, oriented face walks,
  nested face assignment, exclusion carving and active-domain topology audits.
\item[luafemm-bh.lua] Natural cubic slopes, bounded FEMM-compatible smoothing
  and constitutive derivatives without mutating the input table.
\item[luafemm-ans.lua] Static answer-file serialization from a converged
  problem/mesh snapshot, with no remeshing or solver invocation.
\item[luafemm-cache.lua] Passive versioned disk records, exact mesh/solution
  descriptors, validation, atomic replacement and on-demand restoration.
\item[luafemm-mesh.lua] PSLG subdivision, constrained Delaunay meshing,
  local density and final mesh audit.
\item[luafemm-refine.lua] Bounded segment-first quality refinement, a triangle
  priority queue and protected circumcentre insertion.
\item[luafemm-geometry.lua] Winding/parity membership shared by refinement
  and element material classification.
\item[luafemm-predicates.lua] Filtered orientation/incircle/disk signs with
  expansion fallback; diagnostic fallback counters.
\item[luafemm-path.lua] Affine frames and soft-path polygonalisation.
\item[luafemm-materials.lua] Catalogue lookup and static interpretation;
  the separate generated data module retains upstream values.
\item[luafemm-machines.lua] Independent rotor/stator geometry, compound iron
  domains, slot validation, conductor layers and balanced winding compilation.
\item[luafemm-profiles.lua] Named measurement paths, polar and exact circular
  sampling, cached rows and discrete spatial Fourier analysis.
\item[luafemm-tex.lua] TeX boundary: escaped numeric strings, picture
  initialisation, material resolution and emitted drawing commands.
\end{description}
Public functions have LDoc-style comments with units, arguments and
return values. Private numerical helpers document invariants and algorithmic
choices rather than restating each arithmetic operation. StyLua configuration
fixes Lua 5.3 syntax, four-space indentation and a 100-column target.

\subsection{Core calls}
The public constructor copies its input; material and region declarations
also take independent copies. The structures below are owned by their model.
Callbacks and logging functions are retained as function references.
\begin{codeexample}[code only]
local femm = require('luafemm')
local model = femm.new {
  unit = .001, mesher = 'delaunay', h = 2,
  xmin = -20, xmax = 20, ymin = -20, ymax = 20,
  tolerance = 1e-7, max_newton = 60
}
femm.material(model, 'magnet', {library='cast-alnico-5-lng37'})
femm.region(model, 'magnet', {{-3,-5},{3,-5},{3,5},{-3,5}}, 0, 90)
femm.solve(model)
local bx, by, az, hx, hy = femm.sample(model, 0, 8)
\end{codeexample}

|new(options)| accepts the numerical counterparts of problem keys:
|h|, |xmin|, |xmax|, |ymin|, |ymax|, |mesher|, |mesh_tolerance|,
|min_angle|, |max_area|, |max_steiner|, |tolerance|, |max_newton|,
|depth| and |field_model| (default |open|).
Quality options use degrees, model units squared and additional vertex counts,
with defaults 0, 0 and 5000 respectively. |unit| is metres per model coordinate unit,
default .001. Optional |boundary(x,y)| and |source(x,y)| callbacks receive
SI coordinates and return $A_z$ in T\,m and $J_z$ in A/m$^2$, respectively.
The source callback replaces region currents. |log(message)| receives
iteration diagnostics. The Lua-only mesher default is |grid|.

|material(model,name,properties)| accepts |mur|, |bh|, |coercivity|,
|remanence|, |current| and |library|. |region(model,name,points,current,angle,mesh_size)|
uses model coordinates, A/m$^2$, degrees and model-unit local spacing.
Omitted current inherits the material source; omitted angle and local
spacing are zero. |mesh(model)| prepares the discrete model, and
|solve(model)| meshes if needed and returns the same solved model.

Region declarations return the owned region record. Set its |ideal_domain=true|
before meshing to include it in |field_model='ideal'|. Use
|excitation(model,x,y,NI)| for an ideal-window seed, and
|flux(model,points)| for signed right-normal section flux in Wb; the latter
requires |depth| and validates that the section stays in the active domain.
|require('luafemm-ideal').mean(model,node,cycle,kind)| returns a model-coordinate
polyline for a registered component cycle. The numerical module does not create
TeX nodes or profile names. See Section~\ref{sec:ideal} for units and limitations.

For compound regions, use
\nolinkurl{region_contours(model,name,contours,current,angle,mesh_size,fill_rule)}.
Each contour is an array of vertex pairs without a repeated closing vertex.
The fill rule is |'nonzero'| (default) or |'even odd'|. Other arguments use
the same units and defaults as |region|. Inputs are copied; every contour
must be individually simple. Model regions store their geometry in |contours|;
use |region.contours[1]| to access the first contour.
This API defines material holes, never exclusions from the PDE domain.

|sample(model,x,y)| returns |Bx,By,Az,Hx,Hy| in SI; inputs use model
coordinates. |contours(model,count,levels)| returns oriented path arrays
with |x| and |y| in model units and a |level| in T\,m. Optional explicit
|levels| replace the automatic level count. |constitutive(material,B)|
returns $\mu_0H_{\rm table}/B$ and $\mu_0\,dH_{\rm table}/dB$.
These are scaled law coefficients, not the physical vector H of a magnet.

Solved |model.stats| has the names listed in Section~\ref{sec:field}.
|model.A| is the nodal potential in T\,m, |model.nodes| contains SI positions,
and |model.triangles| contains P1 connectivity, material references and
element fields. Treat these implementation structures as read-only.
|import_mesh|, |prepare_mesh| and |update_fields| are internal preparation hooks, not a
stable foreign-mesh import interface in this release.

\subsection{Boundary data}
|boundary(model,side,condition)| replaces one side or |all| sides before
meshing. The condition table has |type| (|dirichlet|, |neumann| or |robin|),
|value|, |dx|, |dy| and |coefficient|. Omitted type means Dirichlet; omitted
numbers mean zero. The affine data are |value+dx*x+dy*y|, with $x,y$ in
metres, independently of |options.unit|. Their physical units and signs
are given in Section~\ref{sec:boundaries}. Numeric side conditions support
persistent caching. They cannot be combined with the older |boundary(x,y)|
callback, which remains a prescribed potential on the entire exterior.

The independent module |luafemm-boundary.lua| owns validation, reconstruction
of exterior edge data, gauge compatibility and consistent P1 edge assembly.
The geometric set |model.exterior_nodes| is distinct from |model.fixed_values|,
the prescribed potential map. |model.gauge_node| is present only when a
reference was needed. These prepared structures are internal and read-only.

The generic TeX layer queues |femm/boundary| options until the model's
begin-picture hook runs. A picture-local flag prevents repeated
|femm/problem| keys from replacing that model. After initialisation,
the key uses the command directly. The queue and flag follow TeX grouping;
the resulting physical declarations belong to the Lua model.

\subsection{Component geometry}
|require('luafemm-components').make(kind,options)| constructs |u core| or
|u electromagnet| geometry without a running TeX interpreter. Option names
use underscores, for example |leg_thickness|, |ampere_turns| and |gap_mesh_size|.
The result owns its options, polygon regions, named anchors and envelope
half-dimensions. It does not register a model, resolve materials or solve.
The component key reference gives the corresponding dimensions and defaults.

|current(points,unit,turns)| returns uniform current density from signed
ampere-turns and the area of a transformed polygon. The TeX bridge calls it
after capture and rounding. The PGF shape saves a geometry identifier; anchor
lookup does not depend on the current model and has no registration effects.
Only its drawing hook emits physical paths, once at the final transform.

\subsection{Persistent cache options}
Pass |cache={file='results/magnet',mode='auto'}| to |new| to enable disk
reuse. Lua filenames are relative to the process working directory; the
TeX bridge additionally honours the engine's output directory.
Omitting |cache| disables disk access. An explicit table defaults to
|mode='auto'| and requires a filename stem unless its mode is |off|.
The four modes and validation rules are described in Section~\ref{sec:cache}.

The model's |cache_info.status| is |off|, |pending|, |miss|, |mesh|,
|solution| or |refresh|. |cache_info.file| contains the resolved filename;
|cache_info.reason| gives a diagnostic when applicable. The constructor
copies cache options just like other input tables. Cache configuration
and physical model data must be treated as immutable after meshing.
Calling |mesh| or |solve| again on an already prepared model is idempotent.
Create a new model to change physical inputs or to refresh a calculation.

Persistent caching rejects |source| and |boundary| callbacks: a closure's
captured state cannot be inferred safely from its function body. These
Lua-only callbacks remain supported with |mode='off'|. Logging callbacks
are allowed and excluded from cache descriptors. Cache records contain
no executable code and never reconstruct functions.

\subsection{Catalogue and profile calls}
|require('luafemm-materials').get(id)| returns an independent record with
|id|, |name|, |path|, |fields| and |bh|. Exact original names are also accepted
when unambiguous. |list()| returns all independent records in upstream
order. |properties(id)| returns the effective static material properties.

Load the profile module with |require('luafemm-profiles')|. Its
\nolinkurl{register(model,name,points,options)} function registers a polyline with |samples|,
|offset| and |outside| options. |sample(name)| returns cached rows containing
|s,t,x,y,B,Bx,By,Bt,Bn,H,Hx,Hy,Ht,Hn,Az|. |get(name)| exposes the internal
profile record and should be treated as read-only. \nolinkurl{coordinates(name,component,abscissa)}
formats the selected columns as PGFPlots coordinate pairs. Re-register
a profile after any manual solution modification; otherwise the cached
rows deliberately remain unchanged.

\subsection{Rebuilding the documentation}
The repository uses an installable |tex/| layout, English source comments,
English examples and documentation, and separate upstream data notices.
Run |python3 doc/luatex/luafemm/build-manual.py| from the browsable source
archive root. The builder executes the LuaLaTeX examples inside this manual
from the same sources shown in its listings, checks the reference fragments
in explicit contexts, and builds the index and cross-references. Only the
three named plain/ConTeXt format demonstrations use separate native-format
compilation. Their fresh PDFs are checked and cropped for readability.
All generated bodies, intermediate files and caches stay under |build/|.

\begin{command}{\luafemmversion}
Expandable semantic version, currently |0.11.0-dev|.
\end{command}
\begin{command}{\luafemmversiondate}
Expandable release date, currently |2026-11-12|.
\end{command}
\begin{command}{\luafemmversionnumber}
Expandable integer for simple compatibility comparisons, currently |110000|.
It increases when a release changes the public interface; it is not the
number of commits.
\end{command}

The original package code and documentation use the LaTeX Project Public
License, version 1.3c or later (LPPL-1.3c). Christophe Jorssen is the author and
Current Maintainer; the maintenance status is \texttt{maintained}. The
imported FEMM data retain their upstream terms. PGF's documentation
environments are loaded from TeX Live, not copied into this repository.
The manual's reference-and-example presentation was modelled on tikz-ext;
its prose and examples were written for luafemm.

\subsection{Installation and module ownership}
Load |require('luafemm')| before inspecting implementation modules in a Lua
script. The entry point pins all package modules to its own installed Lua
directory. The TeX loader selects that directory next to the actual loaded
|luafemm.tex| and inputs its TeX siblings explicitly. Missing components fail;
no second installation supplies a replacement. The selection is process-lifetime
state, whereas picture styles retain ordinary TeX grouping. Captured models,
profiles and named results have the lifetime documented by their respective APIs.

\appendix
% SPDX-License-Identifier: LPPL-1.3c
% Copyright (C) 2026 Christophe Jorssen
% Author and Current Maintainer: Christophe Jorssen <christophe.jorssen@gmail.com>
% LPPL maintenance status: maintained
% This file is part of luafemm. See LICENSE and LICENSES.md for its terms.
\section{Complete examples in the three formats}
\label{sec:format-examples}
The same magnetic scene can be compiled with plain LuaTeX, LuaLaTeX and
ConTeXt MkIV. The following listings show the shared geometry and every
complete wrapper, followed by its actual output. Place \texttt{scene.tex}
beside the selected wrapper and compile from that directory with an
installed luafemm TEXMF tree. The development build does this automatically.
The geometry, profile and solver are identical; page layout follows the format.
Only the plain and ConTeXt wrappers need separate native-format compilation;
their displayed sources are rebuilt and checked with the manual. All LuaLaTeX
examples execute inside the manual itself.

\subsection{A minimal plain-LuaTeX document}
This self-contained example uses the same picture keys and physical paths as
LaTeX and ConTeXt. It needs no format-specific wrapper or older command interface.
\luafemmformatexample{plain}

\subsection{Shared physical scene}
The alnico rectangle, iron return piece and passive copper region precede
the first field request. The curved measurement path is registered before
the field is solved. Each wrapper subsequently plots that same profile.
\luafemmsource{formats/scene}

\subsection{Plain LuaTeX}
Compile with \texttt{luatex plain.tex}. Plain environment commands are
used directly; the source requires no LaTeX wrapper.
\luafemmformatexample{formats/plain}

\subsection{LuaLaTeX}
Compile with \texttt{lualatex latex.tex}. The package wrapper and the
library entry point lead to the same generic implementation.
\luafemmexample{formats/latex}

\subsection{ConTeXt MkIV}
Compile with \texttt{context --luatex context.tex}. The explicit engine
selection matters on installations whose default ConTeXt engine is LMTX.
\luafemmformatexample{formats/context}

\section{Complete FEMM material catalogue}
\label{sec:catalogue}
The identifiers below are exact choices for |material| and |library|.
The original names are retained, including the two distinct Supermalloy
records. The full data and original folder paths are available through
the Lua catalogue and in \texttt{docs/materials-catalog.md}.
\input{build/doc/materials-catalog.tex}
\clearpage
\printindex
\end{document}
